% Traduction : verbes défectifs, phrase restrictive, discours indirect — Espagnol Tle A — Cours signé OnBuch+
% Programme officiel MINESEC. Compiler avec : tectonic E14-verbos-defectivos-restrictiva-estilo-indirecto.tex
\newcommand{\DOCMATIERE}{Espagnol}
\newcommand{\DOCNIVEAU}{Tle A}
\newcommand{\DOCMODULE}{Módulo 3 : Vie citoyenne et ouverture au monde}
\newcommand{\DOCLECON}{E14}
\newcommand{\DOCTITRE}{Traduction : verbes défectifs, phrase restrictive, discours indirect}
% OnBuch+ — préambule « cours signés » ENRICHI (compilé avec tectonic / XeLaTeX)
% Système visuel « Pop » d'OnBuch+ : fond crème (jamais blanc), bordures encre
% épaisses, ombres « dures » décalées, titres Archivo Black en capitales.
% Version MATHÉMATIQUES : graphes (pgfplots/tikz), boîtes pédagogiques
% (définition, propriété, méthode, attention, le savais-tu, exercice résolu,
% exercice, TP, bilan), page de garde et signature OnBuch+ ancrée en bas.
%
% Macros attendues dans main.tex AVANT \input{preamble} :
%   \DOCMATIERE  \DOCNIVEAU  \DOCMODULE  \DOCLECON (numéro)  \DOCTITRE
\documentclass[11pt]{article}

\usepackage{fontspec}
\usepackage[spanish,french]{babel} % français = langue principale ; l'espagnol via \esp{..} / espagnol
\usepackage[a4paper,margin=2cm,top=3.4cm,bottom=2.8cm,headheight=1.4cm,footskip=1.3cm]{geometry}
\usepackage{amsmath,amssymb,amsfonts}
\usepackage{mathtools} % \xmapsto, \coloneqq... souvent utilisés par les modèles
\usepackage{cancel} % simplifications barrées (bilans de forces)
% \abs{z} (module, valeur absolue) et \norm{u} (norme), utilisés par les modèles
\providecommand{\abs}{}\let\abs\relax\DeclarePairedDelimiter{\abs}{\lvert}{\rvert}
\providecommand{\norm}{}\let\norm\relax\DeclarePairedDelimiter{\norm}{\lVert}{\rVert}
\usepackage[table]{xcolor} % [table] : fournit aussi \rowcolors (alternance de lignes)
\usepackage{pagecolor}
\usepackage{tikz}
\usetikzlibrary{arrows.meta,positioning,calc,shapes.geometric,shapes.misc,shapes.arrows,decorations.pathmorphing,decorations.pathreplacing,decorations.markings,patterns,shadows,mindmap,trees,fit,backgrounds,matrix,intersections,angles,quotes}
\usepackage{pgfplots}
\usepackage[version=4]{mhchem} % équations nucléaires \ce{^{235}_{92}U}
\usepackage[european,siunitx]{circuitikz} % schémas électriques
\pgfplotsset{compat=1.18}
\usepgfplotslibrary{fillbetween,groupplots} % aires entre courbes (calcul intégral)
\usepackage{enumitem}
\usepackage{fancyhdr}
% [most] charge toutes les bibliothèques tcolorbox (listings, minted, raster,
% documentation...) et gonfle inutilement la mémoire TeX sur un document long
% (~300 boîtes) : on ne charge que ce qui est réellement utilisé.
\usepackage[skins,breakable]{tcolorbox}
\usepackage{graphicx}
\usepackage{colortbl}
\usepackage{booktabs}
\usepackage{tabularx}
\usepackage{diagbox} % case d'en-tête diagonale des tableaux à double entrée
% Colonnes extensibles centrée / gauche / droite (C, L, R), souvent utilisées par les modèles
\newcolumntype{C}{>{\centering\arraybackslash}X}
\newcolumntype{L}{>{\raggedright\arraybackslash}X}
\newcolumntype{R}{>{\raggedleft\arraybackslash}X}
\usepackage{array}
\usepackage{multicol}
\usepackage{siunitx}
% parse-numbers=false : accepte aussi \SI{2,5 \times 10^{-3}}{...}
\sisetup{locale=FR,output-decimal-marker={,},group-minimum-digits=4,parse-numbers=false}
\DeclareSIUnit\ohm{\ensuremath{\symup{\Omega}}} % U+2126 absent de la police
\DeclareSIUnit\division{div} % réglages d'oscilloscope (ms/div, V/div)
\DeclareSIUnit\tr{tr} % tours (vitesse de rotation, ex. \SI{1500}{\tr\per\minute})
\DeclareSIUnit\molar{M} % concentration molaire (ex. \SI{3}{\molar}, \SI{1}{\micro\molar})
\DeclareSIUnit\an{an} % année (le modèle confond parfois avec \year, un compteur TeX) — ex. \SI{5}{\kilo\meter\per\an}
\DeclareSIUnit\FCFA{FCFA} % franc CFA (ex. \SI{500}{\FCFA\per\kilo\gram})
\DeclareSIUnit\degreeBrix{\textdegree Brix} % teneur en sucre d'un sirop/jus (réfractométrie)
\DeclareSIUnit\month{mois} % le modèle utilise parfois \month, un compteur TeX, comme unité de durée
\usepackage{qrcode}
% \degree : commande fréquemment utilisée par les modèles pour un angle ou une
% température en dehors de \SI{}{} (non fournie par les paquets chargés).
\providecommand{\degree}{^{\circ}}
% \qed (fin de démonstration), utilisé par les modèles sans amsthm
\providecommand{\qed}{\hfill$\square$}
% \barre : double barre des valeurs interdites dans un tableau de variations (tkz-tab)
\providecommand{\barre}{\ensuremath{\|}}
% environnement proof (amsthm non chargé) utilisé par les modèles
\ifdefined\proof\else\newenvironment{proof}[1][Démonstration]{\par\noindent\textit{#1.}\ }{\qed\par}\fi
\usepackage{lastpage}
\usepackage{hyperref}
\hypersetup{hidelinks}

% ============================================================
% Palette « Pop » OnBuch+ — mêmes hex que la classe P (plus_ui.dart)
% ============================================================
\definecolor{poporange}{HTML}{F59321}
\definecolor{popdark}{HTML}{E07A0C}
\definecolor{popcream}{HTML}{FFF6E8}
\definecolor{popink}{HTML}{1C1714}
\definecolor{poppurple}{HTML}{7A5AE0}
\definecolor{popgreen}{HTML}{2E9E6A}
\definecolor{poppink}{HTML}{E0457B}
\definecolor{popblue}{HTML}{2F7DE1}
\definecolor{popgold}{HTML}{C98A12}
\definecolor{popmuted}{HTML}{8A7F76}
\definecolor{popline}{HTML}{EADFD2}
\definecolor{popgray}{HTML}{8A7F76}
\colorlet{popgrey}{popgray}
% Alias tolérants (le modèle nomme parfois les couleurs différemment)
\colorlet{popwhite}{white}
\colorlet{popblack}{popink}
\colorlet{popred}{poppink}
\colorlet{popyellow}{popgold}
% Teintes claires (fonds de tableaux/encadrés) — le modèle nomme parfois
% les couleurs avec un suffixe L (ex. popblueL) sans les avoir définies.
\colorlet{poporangeL}{poporange!20}
\colorlet{popdarkL}{popdark!20}
\colorlet{poppurpleL}{poppurple!20}
\colorlet{popgreenL}{popgreen!20}
\colorlet{poppinkL}{poppink!20}
\colorlet{popblueL}{popblue!20}
\colorlet{popgoldL}{popgold!20}
\colorlet{popredL}{popred!20}
\colorlet{popgrayL}{popgray!20}
% Teintes claires pour fonds de tableaux / zones
\colorlet{popblueL}{popblue!10!white}
\colorlet{poporangeL}{poporange!14!white}
\colorlet{popgreenL}{popgreen!12!white}
\colorlet{poppurpleL}{poppurple!12!white}
\colorlet{poppinkL}{poppink!12!white}
\colorlet{popgoldL}{popgold!14!white}

\pagecolor{popcream}

% ============================================================
% Typographie
% ============================================================
\defaultfontfeatures{Ligatures=TeX}
\setmainfont{PlusJakartaSans-Medium.ttf}[Path=../../../fonts/,BoldFont=PlusJakartaSans-Bold.ttf]
% Maths en sans-serif (Fira Math) avec chiffres et lettres droites de Plus
% Jakarta, pour que les formules se fondent dans le texte.
\usepackage[math-style=ISO,bold-style=ISO]{unicode-math}
\setmathfont{FiraMath-Regular.otf}[Path=../../../fonts/]
\setmathfont{PlusJakartaSans-Medium.ttf}[Path=../../../fonts/,range={up/num,up/{Latin,latin},\mathup}]
\usepackage{icomma} % virgule décimale sans espace en mode math
% Caractères mathématiques Unicode absents des polices de texte (Plus Jakarta,
% Archivo Black) : rendus par leur équivalent mathématique, en texte comme en
% formule — sinon ils s'affichent en carré vide (ex. « Arithmétique dans ℤ »).
\usepackage{newunicodechar}
\newunicodechar{ℝ}{\ensuremath{\mathbb{R}}}
\newunicodechar{ℤ}{\ensuremath{\mathbb{Z}}}
\newunicodechar{ℂ}{\ensuremath{\mathbb{C}}}
\newunicodechar{ℕ}{\ensuremath{\mathbb{N}}}
\newunicodechar{ℚ}{\ensuremath{\mathbb{Q}}}
\newunicodechar{𝔻}{\ensuremath{\mathbb{D}}}
\newunicodechar{α}{\ensuremath{\alpha}}
\newunicodechar{β}{\ensuremath{\beta}}
\newunicodechar{γ}{\ensuremath{\gamma}}
\newunicodechar{δ}{\ensuremath{\delta}}
\newunicodechar{ε}{\ensuremath{\varepsilon}}
\newunicodechar{θ}{\ensuremath{\theta}}
\newunicodechar{λ}{\ensuremath{\lambda}}
\newunicodechar{μ}{\ensuremath{\mu}}
\newunicodechar{σ}{\ensuremath{\sigma}}
\newunicodechar{τ}{\ensuremath{\tau}}
\newunicodechar{φ}{\ensuremath{\varphi}}
\newunicodechar{ψ}{\ensuremath{\psi}}
\newunicodechar{ω}{\ensuremath{\omega}}
\newunicodechar{Σ}{\ensuremath{\Sigma}}
% majuscules produites par \MakeUppercase dans les titres (α-aminés -> Α)
\newunicodechar{Α}{\ensuremath{\alpha}}
\newunicodechar{Β}{\ensuremath{\beta}}
\newunicodechar{Γ}{\ensuremath{\Gamma}}
\newunicodechar{Δ}{\ensuremath{\Delta}}
\newunicodechar{Ε}{\ensuremath{\varepsilon}}
\newunicodechar{Θ}{\ensuremath{\Theta}}
\newunicodechar{Λ}{\ensuremath{\Lambda}}
\newunicodechar{Μ}{\ensuremath{\mu}}
\newunicodechar{Τ}{\ensuremath{\tau}}
\newunicodechar{Φ}{\ensuremath{\Phi}}
\newunicodechar{Ψ}{\ensuremath{\Psi}}
\newunicodechar{Ω}{\ensuremath{\Omega}}
\newunicodechar{↦}{\ensuremath{\mapsto}}
\newunicodechar{∈}{\ensuremath{\in}}
\newunicodechar{∉}{\ensuremath{\notin}}
\newunicodechar{∧}{\ensuremath{\wedge}}
\newunicodechar{⇔}{\ensuremath{\Leftrightarrow}}
\newunicodechar{⟺}{\ensuremath{\Longleftrightarrow}}
\newunicodechar{⇒}{\ensuremath{\Rightarrow}}
\newunicodechar{⟹}{\ensuremath{\Longrightarrow}}
\newunicodechar{∘}{\ensuremath{\circ}}
\newunicodechar{∪}{\ensuremath{\cup}}
\newunicodechar{∩}{\ensuremath{\cap}}
\newunicodechar{≡}{\ensuremath{\equiv}}
\newunicodechar{⊂}{\ensuremath{\subset}}
\newunicodechar{⊥}{\ensuremath{\perp}}
\newunicodechar{∃}{\ensuremath{\exists}}
\newunicodechar{∀}{\ensuremath{\forall}}
\newunicodechar{∛}{\ensuremath{\sqrt[3]{\,}}}
\newunicodechar{∜}{\ensuremath{\sqrt[4]{\,}}}
\newunicodechar{⃗}{\ensuremath{{}^{\to}}}
\newunicodechar{ⁿ}{\ifmmode^{n}\else\textsuperscript{\ensuremath{n}}\fi}
\newunicodechar{ˣ}{\ifmmode^{x}\else\textsuperscript{\ensuremath{x}}\fi}
\newunicodechar{ᵇ}{\ifmmode^{b}\else\textsuperscript{\ensuremath{b}}\fi}
\newunicodechar{ʳ}{\ifmmode^{r}\else\textsuperscript{\ensuremath{r}}\fi}
\newunicodechar{ᵘ}{\ifmmode^{u}\else\textsuperscript{\ensuremath{u}}\fi}
\newunicodechar{ʸ}{\ifmmode^{y}\else\textsuperscript{\ensuremath{y}}\fi}
\newunicodechar{ᵃ}{\ifmmode^{a}\else\textsuperscript{\ensuremath{a}}\fi}
\newunicodechar{ᵉ}{\ifmmode^{e}\else\textsuperscript{\ensuremath{e}}\fi}
\newunicodechar{ᵏ}{\ifmmode^{k}\else\textsuperscript{\ensuremath{k}}\fi}
\newunicodechar{ᵖ}{\ifmmode^{p}\else\textsuperscript{\ensuremath{p}}\fi}
\newunicodechar{ᴺ}{\ifmmode^{N}\else\textsuperscript{\ensuremath{N}}\fi}
\newunicodechar{ᶜ}{\ifmmode^{c}\else\textsuperscript{\ensuremath{c}}\fi}
\newunicodechar{ʰ}{\ifmmode^{h}\else\textsuperscript{\ensuremath{h}}\fi}
\newunicodechar{ᵠ}{\ifmmode^{q}\else\textsuperscript{\ensuremath{q}}\fi}
\newunicodechar{ₙ}{\ifmmode_{n}\else\textsubscript{\ensuremath{n}}\fi}
\newunicodechar{ₐ}{\ifmmode_{a}\else\textsubscript{\ensuremath{a}}\fi}
\newunicodechar{ᵢ}{\ifmmode_{i}\else\textsubscript{\ensuremath{i}}\fi}
\newunicodechar{ᵣ}{\ifmmode_{r}\else\textsubscript{\ensuremath{r}}\fi}
\newunicodechar{ₖ}{\ifmmode_{k}\else\textsubscript{\ensuremath{k}}\fi}
\newunicodechar{ᵦ}{\ifmmode_{\beta}\else\textsubscript{\ensuremath{\beta}}\fi}
\newunicodechar{ₓ}{\ifmmode_{x}\else\textsubscript{\ensuremath{x}}\fi}
\newfontfamily\popdisplay{ArchivoBlack-Regular.ttf}[Path=../../../fonts/]
\newfontfamily\poplogo{SpaceGrotesk-Bold.ttf}[Path=../../../fonts/]
\newfontfamily\popsemi{PlusJakartaSans-SemiBold.ttf}[Path=../../../fonts/]
\newfontfamily\popxbold{PlusJakartaSans-ExtraBold.ttf}[Path=../../../fonts/]
\newfontfamily\popxboldls{PlusJakartaSans-ExtraBold.ttf}[Path=../../../fonts/,LetterSpace=3.0]

\setlist[enumerate]{leftmargin=1.7em,itemsep=5pt,topsep=4pt}
\setlist[itemize]{leftmargin=1.7em,itemsep=4pt,topsep=4pt,label={\color{poporange}\textbullet}}
\setlength{\parindent}{0pt}
\setlength{\parskip}{5pt}
\renewcommand{\arraystretch}{1.25}

% pgfplots : style Pop par défaut
\pgfplotsset{
  every axis/.append style={axis line style={popink,line width=1pt},tick style={popink},
    grid style={popline,line width=0.5pt},label style={font=\small},tick label style={font=\footnotesize}},
  popcurve/.style={poporange,line width=1.8pt,no markers,smooth},
}
% Même style disponible dans un \draw TikZ simple (hors pgfplots)
\tikzset{popcurve/.style={poporange,line width=1.8pt}}
\tikzset{popgrid/.style={popline,very thin}}
\colorlet{popcurve}{poporange} % le nom du style est parfois utilisé comme couleur

% ============================================================
% Pill / squiggle
% ============================================================
\newcommand{\poppill}[3]{%
  \tikz[baseline=-0.55ex]\node[draw=popink,line width=1.2pt,rounded corners=2.6pt,%
    fill=#2,inner xsep=6.5pt,inner ysep=3pt]{%
    \popxboldls\fontsize{7}{7}\selectfont\color{#3}\MakeUppercase{#1}};%
}
\newcommand{\popsquiggle}[1]{%
  \tikz[baseline=-0.4ex]{
    \draw[#1,line width=2.6pt,line cap=round]
      plot [smooth] coordinates {(0,0.09) (0.34,0.01) (0.68,0.21) (1.02,0.01) (1.36,0.10)};
  }%
}
% Mot-clé surligné (marqueur orange sous le texte)
\newcommand{\cle}[1]{{\popxbold\color{popdark}#1}}

% ============================================================
% En-tête / pied
% ============================================================
\pagestyle{fancy}
\fancyhf{}
\renewcommand{\headrulewidth}{0pt}
\renewcommand{\footrulewidth}{0pt}
\lhead{\raisebox{-0.15ex}{\poplogo\fontsize{13}{13}\selectfont\color{popink}OnBuch\color{poporange}+}}
\rhead{\poppill{\DOCMATIERE\ \textbullet\ \DOCNIVEAU\ \textbullet\ Leçon \DOCLECON}{white}{popink}}
\lfoot{\popxbold\fontsize{7.6}{7.6}\selectfont\color{popmuted}\MakeUppercase{Cours signé OnBuch+}\ \textbullet\ {\popsemi\DOCTITRE}}
\rfoot{\tikz[baseline=-0.6ex]\node[rounded corners=8.5pt,fill=popink,minimum height=17pt,minimum width=17pt,inner xsep=5pt,inner sep=0]{\popxbold\fontsize{8}{8}\selectfont\color{popcream}\thepage\,/\,\pageref*{LastPage}};}

% ============================================================
% Titres
% ============================================================
\newcommand{\chaptitre}[2]{%
  \begin{tcolorbox}[enhanced,colback=poporange,colframe=popink,boxrule=2.6pt,arc=11pt,
      drop shadow=popink,left=15pt,right=15pt,top=12pt,bottom=11pt,before skip=3pt,after skip=18pt]
    \poppill{#2}{popink}{popcream}\par\vspace{9pt}
    {\raggedright\hyphenpenalty=10000\exhyphenpenalty=10000\popdisplay\fontsize{22}{24}\selectfont\color{popcream}\MakeUppercase{#1}\par}
  \end{tcolorbox}
}
% Section numérotée : pastille orange avec le numéro + titre Archivo
\newcounter{popsec}
\newcommand{\coursec}[1]{%
  \stepcounter{popsec}\setcounter{popsub}{0}%
  \par\vspace{16pt}\noindent
  \begin{minipage}{\linewidth}
  \tikz[baseline=-0.7ex]\node[draw=popink,line width=1.6pt,fill=poporange,rounded corners=4pt,minimum size=22pt,inner sep=0]{\popdisplay\fontsize{12}{12}\selectfont\color{popcream}\thepopsec};%
  \hspace{8pt}{\popdisplay\fontsize{15}{17}\selectfont\color{popink}\MakeUppercase{#1}}\par
  \vspace{4pt}\noindent{\color{poporange}\rule{36pt}{3.6pt}}
  \end{minipage}\par\nopagebreak\vspace{8pt}
}
\newcounter{popsub}
\newcommand{\courssub}[1]{%
  \stepcounter{popsub}%
  \par\vspace{9pt}\noindent{\popxbold\fontsize{11.5}{13}\selectfont\color{popink}{\color{poporange}\thepopsec.\thepopsub}\hspace{6pt}#1}\par\nopagebreak\vspace{4pt}
}

% ============================================================
% Boîtes pédagogiques — même recette : fond blanc, bordure épaisse colorée,
% ombre dure encre, badge plein. Titre optionnel [#1].
% ============================================================
\tcbset{popbase/.style={breakable,enhanced,colback=white,boxrule=2.3pt,arc=9pt,
  shadow={3pt}{-3pt}{0pt}{popink},left=14pt,right=14pt,top=11pt,bottom=11pt,before skip=11pt,after skip=15pt,
  fontupper=\popsemi\fontsize{10.6}{14.5}\selectfont\color{popink},
  fontlower=\popsemi\fontsize{10.6}{14.5}\selectfont\color{popink}}}
% Coche dessinée (le glyphe ✓ n'existe pas dans les polices du document)
\newcommand{\popcheck}[1]{\tikz[baseline=-0.5ex]\draw[#1,line width=1.6pt,line cap=round,line join=round](-0.13,0.0)--(-0.04,-0.09)--(0.13,0.1);}
\providecommand{\coche}{\popcheck{popgreen}}
\providecommand{\resultat}[1]{\par\smallskip\noindent\fcolorbox{popgreen}{popgreenL}{\parbox{\dimexpr\linewidth-2\fboxsep-2\fboxrule\relax}{\textbf{Résultat.} #1}}\par\smallskip}
\newcommand{\popboxhead}[3]{\poppill{#1}{#2}{white}\ifx\relax#3\relax\else\hspace{6pt}{\popxbold\fontsize{10.6}{12}\selectfont\color{popink}#3}\fi\par\vspace{7pt}}

\newtcolorbox{aretenir}[1][]{popbase,colframe=popgreen,before upper={\popboxhead{À retenir}{popgreen}{#1}}}
% Texte en espagnol (document de lecture, dialogue, poème, extrait) : cadre
% d'étude. Un titre optionnel (source, auteur) s'affiche dans l'en-tête.
\newtcolorbox{textebox}[1][]{popbase,colframe=popblue,before upper={\popboxhead{Texto}{popblue}{#1}\begin{otherlanguage}{spanish}},after upper={\end{otherlanguage}}}
% Marques ✓ / ✗ (forme correcte / fautive) souvent utilisées par les modèles
\providecommand{\cross}{\textcolor{poppink}{\ensuremath{\times}}}
\providecommand{\xmark}{\cross}\providecommand{\crossmark}{\cross}\providecommand{\cmark}{\ensuremath{\checkmark}}
% Ligne de réponse / justification à compléter (inventée par les modèles)
\providecommand{\justification}[1]{\par\noindent\textit{Justification~:} \dotfill\par}
% \textipa (package tipa non chargé) : la police ne couvre pas l'API -> simple passage
\providecommand{\textipa}[1]{#1}
% Soulignement (ulem non chargé) : \uline, \ul
\providecommand{\uline}[1]{\underline{#1}}\providecommand{\ul}[1]{\underline{#1}}
% Mot ou phrase en espagnol à l'intérieur d'un paragraphe français
\newcommand{\esp}[1]{\ifmmode\text{\textit{#1}}\else\begingroup\selectlanguage{spanish}\itshape #1\endgroup\fi}
\newtcolorbox{exemplebox}[1][]{popbase,colframe=poporange,before upper={\popboxhead{Exemple}{poporange}{#1}}}
\newtcolorbox{definition}[1][]{popbase,colframe=popblue,before upper={\popboxhead{Définition}{popblue}{#1}}}
\newtcolorbox{propriete}[1][]{popbase,colframe=poppurple,before upper={\popboxhead{Propriété}{poppurple}{#1}}}
\newtcolorbox{methode}[1][]{popbase,colframe=popgold,before upper={\popboxhead{Méthode}{popgold}{#1}}}
\newtcolorbox{attention}[1][]{popbase,colframe=poppink,before upper={\popboxhead{Attention}{poppink}{#1}}}
\newtcolorbox{experience}[1][]{popbase,colframe=popblue,colback=popblueL,before upper={\popboxhead{Expérience}{popblue}{#1}}}
% « Le savais-tu ? » : carte violette pleine (contexte, culture, Cameroun)
\newtcolorbox{savaistu}[1][]{popbase,colback=poppurple,colframe=popink,
  fontupper=\popsemi\fontsize{10.4}{14.2}\selectfont\color{white},
  before upper={\poppill{Le savais-tu\,?}{white}{poppurple}\ifx\relax#1\relax\else\hspace{6pt}{\popxbold\color{white}#1}\fi\par\vspace{7pt}}}
% Exercice résolu : énoncé en haut, \tcblower puis la solution sur fond crème
\newcounter{popexo}\newcounter{popexr}
\newtcolorbox{exoresolu}[1][]{popbase,colframe=poporange,colbacklower=popcream,
  segmentation style={popink,line width=1.2pt,dashed},
  before upper={\stepcounter{popexr}\popboxhead{Exercice résolu \thepopexr}{poporange}{#1}},
  before lower={\poppill{Solution}{popink}{popcream}\par\vspace{6pt}}}
% Exercice (non résolu en place) — la correction est dans la partie Corrigés
\newtcolorbox{exercice}[1][]{popbase,colframe=popink,
  before upper={\stepcounter{popexo}\popboxhead{Exercice \thepopexo}{popink}{#1}}}
% Corrigé
\newtcolorbox{corrige}[1][]{popbase,colframe=popgreen,colback=popgreenL,
  before upper={\popboxhead{Corrigé}{popgreen}{#1}}}
% Objectifs / prérequis / bilan
\newtcolorbox{objectifs}{popbase,colframe=popink,colback=poporangeL,
  before upper={\popboxhead{Objectifs de la leçon}{popink}{}}}
\newtcolorbox{prerequis}{popbase,colframe=popink,colback=popblueL,
  before upper={\popboxhead{Prérequis}{popblue}{}}}
\newtcolorbox{bilan}{popbase,colframe=popink,colback=white,boxrule=2.8pt,
  before upper={\popboxhead{Fiche bilan}{poporange}{L'essentiel à connaître pour l'examen}}}
% Figure encadrée avec légende
\newtcolorbox{popfigure}[1][]{enhanced,breakable=false,colback=white,colframe=popline,boxrule=1.4pt,arc=8pt,
  left=10pt,right=10pt,top=10pt,bottom=8pt,before skip=10pt,after skip=12pt,halign=center,
  lower separated=false,halign lower=center,
  fontlower=\popsemi\fontsize{9}{11.5}\selectfont\color{popmuted},
  #1}
\newcommand{\legende}[1]{\par\vspace{6pt}{\popsemi\fontsize{9}{11.5}\selectfont\color{popmuted}#1}}

% ============================================================
% Signature OnBuch+
% ============================================================
\newcommand{\poplogoinline}[1]{{\poplogo\fontsize{#1}{#1}\selectfont\color{popink}OnBuch\color{poporange}+}}
\newcommand{\signatureonbuch}{%
  \begin{tcolorbox}[enhanced,colback=popink,colframe=popink,boxrule=2.2pt,arc=10pt,
      drop shadow=poporange,left=14pt,right=14pt,top=10pt,bottom=10pt,before skip=0pt,after skip=0pt]
    \tikz[baseline=-0.7ex]\node[circle,fill=popgreen,draw=popcream,line width=1.3pt,minimum size=22pt,inner sep=0]{\popcheck{white}};%
    \hspace{9pt}\begin{minipage}[c]{0.62\linewidth}
      {\popxbold\fontsize{12}{13}\selectfont\color{popcream}Signé\ \textbullet\ {\poplogo OnBuch\color{poporange}+}}\par\vspace{2pt}
      {\raggedright\popsemi\fontsize{8}{10.5}\selectfont\color{popline}\MakeUppercase{Équipe pédagogique OnBuch+}\\\MakeUppercase{Conforme au programme officiel MINESEC}\par}
    \end{minipage}\hfill
    {\poplogo\fontsize{22}{22}\selectfont\color{popcream}OnBuch\color{poporange}+}
  \end{tcolorbox}%
}

% ============================================================
% Page de garde
% #1 titre · #2 matière · #3 niveau · #4 module · #5 n° leçon · #6 durée ·
% #7 description / objectifs (texte ou itemize)
% ============================================================
\newcommand{\pagegarde}[7]{%
  \thispagestyle{empty}
  \newgeometry{a4paper,margin=1.7cm,top=1.6cm,bottom=1.4cm}%
  \begin{tikzpicture}[remember picture,overlay]
    \fill[poporange!18] ([xshift=-3.2cm,yshift=-3.6cm]current page.north east) circle (4.6cm);
    \fill[poppurple!14] ([xshift=2.4cm,yshift=7.6cm]current page.south west) circle (3.8cm);
  \end{tikzpicture}%
  {\poplogo\fontsize{30}{30}\selectfont\color{popink}OnBuch\color{poporange}+}\hfill
  \raisebox{0.6ex}{\poppill{Cours signé \textbullet\ Premium}{popink}{popcream}}\par
  \vspace{1pt}\popsquiggle{poppurple}\par
  \vspace{8pt}
  \begin{tcolorbox}[enhanced,colback=poporange,colframe=popink,boxrule=2.8pt,arc=13pt,
      drop shadow=popink,left=18pt,right=18pt,top=12pt,bottom=13pt,after skip=10pt]
    \poppill{#2 \textbullet\ #3}{popink}{popcream}\hspace{5pt}\poppill{#4}{white}{popink}\par\vspace{8pt}
    {\popdisplay\fontsize{14}{14}\selectfont\color{popink}LEÇON #5}\par\vspace{4pt}
    {\raggedright\hyphenpenalty=10000\exhyphenpenalty=10000\popdisplay\fontsize{25}{27}\selectfont\color{popcream}\MakeUppercase{#1}\par}
    \vspace{8pt}
    {\popxbold\fontsize{9.5}{9.5}\selectfont\color{popink}\MakeUppercase{Durée conseillée : #6}}
  \end{tcolorbox}
  \begin{tcolorbox}[enhanced,colback=white,colframe=popink,boxrule=2.2pt,arc=10pt,
      drop shadow=popink,left=16pt,right=16pt,top=10pt,bottom=10pt,after skip=8pt]
    \poppill{Dans cette leçon}{poppurple}{white}\par\vspace{5pt}
    \popsemi\fontsize{9.3}{12.6}\selectfont\color{popink}#7
  \end{tcolorbox}
  \noindent\begin{minipage}[t]{0.487\linewidth}
    \begin{tcolorbox}[enhanced,colback=popgreen,colframe=popink,boxrule=2.2pt,arc=10pt,
        drop shadow=popink,left=11pt,right=11pt,top=10pt,bottom=10pt,height=3.1cm,valign=center]
      \tcbox[colback=white,colframe=popink,boxrule=1.3pt,arc=5pt,boxsep=0pt,nobeforeafter,
        left=3pt,right=3pt,top=3pt,bottom=3pt,valign=center]{\qrcode[height=1.9cm]{https://play.google.com/store/apps/details?id=com.luvvix.onbuch&hl=fr}}%
      \hspace{8pt}\begin{minipage}[c]{0.5\linewidth}
        {\popdisplay\fontsize{11}{12.5}\selectfont\color{white}\MakeUppercase{Télécharge OnBuch}}\par\vspace{4pt}
        {\raggedright\popsemi\fontsize{8.4}{11}\selectfont\color{white}Gratuit \textbullet\ Google Play\\Tuteur IA, annales, concours, résultats\par}
      \end{minipage}
    \end{tcolorbox}
  \end{minipage}\hfill
  \begin{minipage}[t]{0.487\linewidth}
    \begin{tcolorbox}[enhanced,colback=poppurple,colframe=popink,boxrule=2.2pt,arc=10pt,
        drop shadow=popink,left=11pt,right=11pt,top=10pt,bottom=10pt,height=3.1cm,valign=center]
      \tikz[baseline=-0.7ex]\node[circle,fill=white,draw=popink,line width=1.3pt,minimum size=1.6cm,inner sep=0]{\popdisplay\fontsize{24}{24}\selectfont\color{poppurple}+};%
      \hspace{8pt}\begin{minipage}[c]{0.6\linewidth}
        {\popdisplay\fontsize{11}{12.5}\selectfont\color{white}\MakeUppercase{Passe à OnBuch+}}\par\vspace{4pt}
        {\raggedright\popsemi\fontsize{8.4}{11}\selectfont\color{white}Tous les cours signés, Tuteur IA illimité, fiches PDF — onglet OnBuch+ de l'app\par}
      \end{minipage}
    \end{tcolorbox}
  \end{minipage}
  \vspace{8pt}
  \popcontenu
  \vfill
  \signatureonbuch
  \restoregeometry
  \newpage
}

% Rangée « Ce cours contient » (page de garde)
\newcommand{\poptile}[3]{% #1 couleur #2 gros texte #3 légende
  \begin{tcolorbox}[enhanced,colback=#1,colframe=popink,boxrule=2pt,arc=9pt,shadow={2.5pt}{-2.5pt}{0pt}{popink},
      left=6pt,right=6pt,top=9pt,bottom=9pt,halign=center,valign=center,height=1.75cm,width=\linewidth,nobeforeafter]
    {\popdisplay\fontsize{12.5}{13}\selectfont\color{white}\MakeUppercase{#2}}\par\vspace{3pt}
    {\centering\popsemi\fontsize{7.8}{9.5}\selectfont\color{white}#3\par}
  \end{tcolorbox}}
\newcommand{\popcontenu}{%
  \par\noindent{\popxboldls\fontsize{7.5}{7.5}\selectfont\color{popmuted}CE COURS SIGNÉ CONTIENT}\par\vspace{7pt}
  \noindent\begin{minipage}[t]{0.235\linewidth}\poptile{popblue}{Cours}{complet, illustré, pas à pas}\end{minipage}\hfill
  \begin{minipage}[t]{0.235\linewidth}\poptile{poporange}{Méthodes}{et exercices résolus}\end{minipage}\hfill
  \begin{minipage}[t]{0.235\linewidth}\poptile{poppink}{Exercices}{type examen + corrigés}\end{minipage}\hfill
  \begin{minipage}[t]{0.235\linewidth}\poptile{popgreen}{Fiche bilan}{l'essentiel à retenir}\end{minipage}\par}

% Sommaire
\newtcolorbox{sommaire}{enhanced,colback=white,colframe=popink,boxrule=2.2pt,arc=10pt,
  drop shadow=popink,left=18pt,right=18pt,top=13pt,bottom=13pt,before skip=6pt,after skip=20pt,
  before upper={\poppill{Sommaire}{poppurple}{white}\par\vspace{9pt}\popsemi\fontsize{10.6}{14.5}\selectfont\color{popink}}}

% Signature de fin de cours — poussée tout en bas de la dernière page
\newcommand{\findecours}{%
  \par\vfill\nopagebreak
  \begin{center}\popsquiggle{poporange}\end{center}\vspace{4pt}
  \signatureonbuch
}
\AtBeginDocument{\def\llbracket{\mathopen{[\mkern-3.5mu[}}\def\rrbracket{\mathclose{]\mkern-3.5mu]}}} % intervalles d'entiers (glyphe ⟦ absent de la police)
% Glyphes absents de Fira Math : redéfinis à partir de glyphes disponibles
\AtBeginDocument{%
  \def\setminus{\mathbin{\backslash}}\let\smallsetminus\setminus
  \def\vdots{\mathord{\vbox{\baselineskip3.2pt\lineskiplimit0pt\kern4pt\hbox{.}\hbox{.}\hbox{.}}}}%
  \def\ddots{\mathinner{\mkern1mu\raise6.5pt\hbox{.}\mkern2mu\raise3.6pt\hbox{.}\mkern2mu\raise0.7pt\hbox{.}\mkern1mu}}%
  \def\triangle{\mathord{\scalebox{0.9}{$\symup{\Delta}$}}}%
  \def\nabla{\mathord{\raisebox{\depth}{\scalebox{1}[-1]{$\symup{\Delta}$}}}}%
  \def\checkmark{\popcheck{popdark}}%
  \def\star{\mathbin{\ast}}\def\bigstar{\mathord{\ast}}%
  \def\diamond{\mathbin{\scalebox{0.6}{$\Diamond$}}}%
  \def\bigcup{\mathop{\mathchoice{\raisebox{-0.3em}{\scalebox{1.7}{$\cup$}}}{\scalebox{1.25}{$\cup$}}{\cup}{\cup}}\displaylimits}%
  \def\bigcap{\mathop{\mathchoice{\raisebox{-0.3em}{\scalebox{1.7}{$\cap$}}}{\scalebox{1.25}{$\cap$}}{\cap}{\cap}}\displaylimits}%
  \def\lesseqqgtr{\mathrel{\lessgtr}}\def\bigoplus{\mathop{\oplus}\displaylimits}%
}
\providecommand{\ssi}{si et seulement si} % abréviation fréquente
\AtBeginDocument{\providecommand{\wideparen}{\overparen}} % arc de cercle

%% FIGFIX-BEGIN v5 — correctifs globaux des figures (docs/onbuch-generation/tools/figfix)
\usepackage{adjustbox}\usepackage{etoolbox}\tcbuselibrary{raster}
% toute figure TikZ/pgfplots plus large que la ligne est réduite à la largeur disponible
\BeforeBeginEnvironment{tikzpicture}{\begin{adjustbox}{max width=\linewidth}}
\AfterEndEnvironment{tikzpicture}{\end{adjustbox}}
% légendes pgfplots lisibles par-dessus les courbes
\pgfplotsset{every axis/.append style={legend style={fill=white,fill opacity=0.92,text opacity=1,draw=black!15,inner sep=3pt}}}
% étiquettes d'arêtes (syntaxe edge["..."]) avec fond blanc
\tikzset{every edge quotes/.append style={fill=white,inner sep=1.2pt}}
%% FIGFIX-END
\begin{document}

\pagegarde{Traduction : verbes défectifs, phrase restrictive, discours indirect}{Espagnol}{Tle A}{Módulo 3 : Vie citoyenne et ouverture au monde}{E14}{2 h}{Cette leçon de traduction entraîne les élèves de Terminale A à rendre en espagnol trois structures incontournables du Bac : les verbes défectifs (gustar, doler, parecer...), la phrase restrictive (sólo, no... más que, no... sino) et le discours indirect avec sa concordance des temps. Observation, règles systématiques, comparaison français/espagnol et exercices de thème et de version garantissent la maîtrise des savoir-faire.
\vspace{4pt}
\begin{multicols}{2}\small
\begin{itemize}
\item À la fin de cette leçon, je sais identifier et conjuguer les verbes défectifs gustar, doler, parecer, soler, atañer, concernir, llover
\item je sais construire correctement la phrase à verbe défectif (pronom d'objet indirect + verbe à la 3e personne + sujet)
\item je sais traduire la restriction avec sólo, únicamente, no... más que, no... sino et apenas
\item je sais choisir entre no... más que et no... sino selon le sens
\item je sais passer du discours direct au discours indirect en espagnol (personnes, temps, indicateurs spatio-temporels)
\item je sais appliquer la concordance des temps après un verbe introducteur au passé
\end{itemize}\end{multicols}}

\begin{sommaire}
\begin{multicols}{2}
\begin{enumerate}
\item Observation : les verbes défectifs en action
\item Règles et emplois des verbes défectifs
\item La phrase restrictive
\item Du discours direct au discours indirect : principes
\item Concordance des temps et verbes introducteurs
\item Entraînement intégré : thème et version
\item Activité d'intégration
\item Méthodes et exercices résolus
\item Exercices
\item Corrigés des exercices
\item Fiche bilan
\end{enumerate}
\end{multicols}
\end{sommaire}

\begin{objectifs}
\begin{itemize}
\item À la fin de cette leçon, je sais identifier et conjuguer les verbes défectifs gustar, doler, parecer, soler, atañer, concernir, llover
\item je sais construire correctement la phrase à verbe défectif (pronom d'objet indirect + verbe à la 3e personne + sujet)
\item je sais traduire la restriction avec sólo, únicamente, no... más que, no... sino et apenas
\item je sais choisir entre no... más que et no... sino selon le sens
\item je sais passer du discours direct au discours indirect en espagnol (personnes, temps, indicateurs spatio-temporels)
\item je sais appliquer la concordance des temps après un verbe introducteur au passé
\item je sais employer les verbes introducteurs variés : decir, afirmar, preguntar, ordenar, rogar, aconsejar
\item je sais éviter les pièges francophones (me gusto, solo que, concordance des temps calquée du français)
\item je sais traduire un court texte du français vers l'espagnol et inversement mobilisant ces trois notions
\end{itemize}
\end{objectifs}

\begin{prerequis}
\begin{itemize}
\item Les pronoms personnels d'objet indirect (me, te, le, nos, os, les) — Première
\item L'indicatif présent, l'imparfait, le passé simple et le futur — Première
\item Le subjonctif présent et imparfait (bases) — début de Terminale
\item Le conditionnel présent — Première
\item La négation simple (no) et les pronoms relatifs de base
\end{itemize}
\end{prerequis}

\newpage

\coursec{Observation : les verbes défectifs en action}

\courssub{Situation d'accroche : trois phrases qui résistent}

C'est la rentrée au lycée de Yaoundé. Un élève de Terminale A accueille Diego, correspondant mexicain, et l'emmène au marché Mokolo. Il veut lui dire trois choses très simples en apparence : « j'aime le ndolé », « seuls les élèves sérieux réussissent », et rapporter les paroles du proviseur : « Le directeur a dit que les examens commenceraient le lendemain ». En espagnol, ces trois phrases exigent des structures très particulières : \esp{Me gusta el ndolé}, \esp{Sólo los alumnos serios aprueban}, \esp{El director dijo que los exámenes empezarían al día siguiente}. Aucune ne se traduit mot à mot : la première renverse l'ordre logique français, la deuxième impose une restriction, la troisième déplace les temps et les repères du discours rapporté. Comment maîtriser ces mécanismes pour le thème et la version du Bac ? C'est l'objet de cette leçon. Commençons par observer les \cle{verbes défectifs} en action.

\courssub{Lecture : dialogue au marché}

\begin{textebox}[En el mercado Mokolo]
--- Diego, ¿te gusta la cocina camerunesa?\par
--- Sí, me encanta el ndolé, pero me duelen los ojos cuando pica mucho. A mi hermano no le parece bueno el pimentón.\par
--- ¿Y te gustan los plátanos fritos?\par
--- ¡Me gustan muchísimo! En México solemos comer plátanos también, pero no así.\par
--- Aquí suele llover cuando vamos al mercado por la tarde. ¿Te parece bien volver al colegio?\par
--- Me parece perfecto. Por cierto, ¿qué asuntos atañen al consejo de clase de mañana?\par
--- A nosotros no nos atañe ese asunto; sólo concierne a los profesores.
\end{textebox}

\textbf{Glossaire :} \esp{el ndolé} : le ndolé (plat camerounais) ; \esp{picar} : être pimenté ; \esp{el pimentón} : le piment ; \esp{los plátanos fritos} : les bananes plantain frites ; \esp{así} : ainsi, comme ça ; \esp{por cierto} : à propos ; \esp{el asunto} : l'affaire, la question ; \esp{el consejo de clase} : le conseil de classe ; \esp{concernir} : concerner.

\textbf{Questions de compréhension :}
\begin{enumerate}
\item Quel plat camerounais Diego aime-t-il beaucoup ?
\item Qu'est-ce qui lui fait mal quand le plat est très pimenté ?
\item Que se passe-t-il souvent l'après-midi au marché ?
\item À qui concerne l'affaire du conseil de classe ?
\end{enumerate}

\courssub{Repérage guidé : où est le sujet grammatical ?}

Regardons de près la phrase \esp{Me gusta el fútbol}. Un francophone pense : « j'aime le football », donc « je » est le sujet. Faux ! En espagnol, le verbe \esp{gusta} est à la troisième personne du singulier : il s'accorde avec \esp{el fútbol}. La preuve : si la chose aimée est au pluriel, le verbe passe au pluriel : \esp{Me gustan los libros} « j'aime les livres ». Le mot \esp{me} n'est pas un sujet : c'est un \cle{pronom complément d'objet indirect} (COI), comme dans la phrase française « le football \emph{me} plaît ».

\begin{aretenir}[Le constat fondamental]
La construction espagnole \textbf{inverse} la construction française. En français, la personne qui éprouve le goût est le sujet (« \emph{j'}aime le ndolé »). En espagnol, le \textbf{sujet réel est la chose aimée}, et la personne est exprimée par un pronom COI : \esp{Me gusta el ndolé} = littéralement « le ndolé me plaît ».
\end{aretenir}

\begin{popfigure}
\begin{center}
\begin{tikzpicture}[node distance=1.1cm, >=Stealth]
\node[draw=popgreen, fill=popgreenL, rounded corners, thick, align=center, minimum width=3.4cm, minimum height=1.1cm] (cosa) {\textbf{La chose aimée}\\ \esp{el ndolé}\\ {\small sujet réel}};
\node[draw=popblue, fill=popblueL, rounded corners, thick, align=center, minimum width=3.2cm, minimum height=1.1cm, right=of cosa] (verbo) {\textbf{Verbe 3\up{e} pers.}\\ \esp{gusta / gustan}\\ {\small accord avec la chose}};
\node[draw=poporange, fill=poporangeL, rounded corners, thick, align=center, minimum width=3.2cm, minimum height=1.1cm, right=of verbo] (coi) {\textbf{Pronom COI}\\ \esp{me, te, le...}\\ {\small la personne}};
\draw[->, very thick, popdark] (cosa) -- (verbo);
\draw[->, very thick, popdark] (verbo) -- (coi);
\node[below=0.35cm of verbo, align=center, text=popmuted] {\esp{Me gusta el ndolé} : la chose commande le verbe};
\end{tikzpicture}
\end{center}
\legende{Schéma de la construction inversée : le sujet réel est la chose aimée.}
\end{popfigure}

\courssub{Définition et règle d'accord}

\begin{definition}[Le verbe défectif]
Un \cle{verbe défectif} est un verbe qui ne se conjugue pas à toutes les personnes. Il s'emploie principalement à la \textbf{troisième personne du singulier ou du pluriel} (\esp{gustar}, \esp{doler}, \esp{parecer}) ou \textbf{uniquement à la troisième personne du singulier} : ce sont alors les verbes \emph{impersonnels} (\esp{llover} : il pleut, \esp{atañer} : concerner, \esp{nevar} : neiger, \esp{hacer frío} : faire froid). Le verbe \esp{soler} (avoir l'habitude de) se conjugue à toutes les personnes mais n'existe qu'à quelques temps : présent et imparfait de l'indicatif surtout.
\end{definition}

\begin{propriete}[Accord du verbe avec la chose aimée]
Le verbe s'accorde avec le \textbf{nom ou l'infinitif} qui le suit :
\begin{itemize}
\item chose au singulier : \esp{Me gusta el libro.} « J'aime le livre. »
\item chose au pluriel : \esp{Me gustan los libros.} « J'aime les livres. »
\item suivi d'un infinitif : toujours le singulier : \esp{Me gusta bailar.} « J'aime danser. »
\end{itemize}
Le pronom COI est \textbf{obligatoire} et placé devant le verbe : \esp{me, te, le, nos, os, les}. On peut renforcer avec \esp{a mí, a ti, a él, a nosotros...} : \esp{A mi hermano no le parece bueno el pimentón.}
\end{propriete}

\begin{exemplebox}[Exemples traduits]
\begin{itemize}
\item \esp{Me gusta bailar.} « J'aime danser. »
\item \esp{¿Te duelen los pies?} « Tu as mal aux pieds ? »
\item \esp{Me parece difícil.} « Cela me semble difficile. »
\item \esp{Suele llegar tarde.} « Il a l'habitude d'arriver en retard. »
\item \esp{Llueve mucho en Douala.} « Il pleut beaucoup à Douala. »
\item \esp{Eso no te atañe.} « Cela ne te concerne pas. »
\end{itemize}
\end{exemplebox}

\courssub{Premier tableau de correspondances français / espagnol}

\begin{center}
\begin{tabularx}{\linewidth}{|X|X|X|}
\hline
\rowcolor{popblueL} Français & Español & Remarque \\
\hline
J'aime le football. & Me gusta el fútbol. & sujet réel : \esp{el fútbol} \\
\hline
J'aime les mangues. & Me gustan los mangos. & verbe au pluriel \\
\hline
Tu as mal à la tête. & Te duele la cabeza. & \esp{doler} = « faire mal à » \\
\hline
Nous avons mal aux dents. & Nos duelen los dientes. & COI \esp{nos} + pluriel \\
\hline
Cela me semble juste. & Me parece justo. & \esp{parecer} = « sembler » \\
\hline
Il a l'habitude de se lever tôt. & Suele levantarse temprano. & \esp{soler} + infinitif \\
\hline
Il pleut souvent en juillet. & Suele llover en julio. & verbe impersonnel \\
\hline
Cela concerne les élèves. & Eso atañe a los alumnos. & 3\up{e} pers. sing. uniquement \\
\hline
\end{tabularx}
\end{center}

\courssub{Carte mentale des verbes défectifs}

\begin{popfigure}
\begin{center}
\begin{tikzpicture}[mindmap, grow cyclic, every node/.style={concept, align=center}, root concept/.append style={fill=popblue, text=white, font=\bfseries}, level 1/.append style={level distance=3.1cm, sibling angle=72}, concept color=popblueL]
\node[root concept, align=center]{Verbes\\ défectifs}
  child[concept color=popgreenL]{ node[align=center]{gustar\\ « plaire »\\ \esp{me gusta}} }
  child[concept color=poporangeL]{ node[align=center]{doler\\ « faire mal »\\ \esp{me duelen}} }
  child[concept color=poppinkL]{ node[align=center]{parecer\\ « sembler »\\ \esp{me parece}} }
  child[concept color=popgoldL]{ node[align=center]{soler\\ « avoir l'habitude »\\ \esp{suele + inf.}} }
  child[concept color=poppurpleL]{ node[align=center]{impersonnels\\ \esp{llover, nevar,}\\ \esp{atañer}} };
\end{tikzpicture}
\end{center}
\legende{Carte mentale : les grandes familles de verbes défectifs.}
\end{popfigure}

\begin{exoresolu}[Traduire avec un verbe défectif]
Énoncé : traduis en espagnol : « Mes amis aiment le cacao, mais ils ont mal au ventre quand ils mangent trop de chocolat. »
\tcblower
Solution détaillée : « le cacao » est la chose aimée, donc sujet réel au singulier : \esp{gusta}. La personne est « mes amis » : COI pluriel \esp{les}. « Avoir mal au ventre » se dit avec \esp{doler} : \esp{les duele el vientre} (sujet : \esp{el vientre}, singulier). Traduction : \esp{A mis amigos les gusta el cacao, pero les duele el vientre cuando comen demasiado chocolate.} On remarque le renforcement \esp{a mis amigos}, facultatif mais très fréquent pour préciser la personne.
\end{exoresolu}

\begin{attention}[L'erreur n°1 des francophones]
Ne dis jamais \esp{yo gusto} pour « j'aime » : cette forme n'existe pas dans cet emploi. On dit obligatoirement \esp{ME GUSTA} + nom ou infinitif. Autres pièges : oublier le pronom COI (\esp{gusta el fútbol} seul ne veut rien dire) ; accorder le verbe avec la personne au lieu de la chose (\esp{me gusto los libros} est faux : \esp{me gustan los libros}) ; traduire « il me semble » par \esp{me semblo} : le verbe français « sembler » se dit \esp{parecer}.
\end{attention}

\begin{exercice}[À toi de jouer]
Traduis en espagnol : 1. J'aime les cours d'espagnol. 2. Est-ce que tu as mal aux yeux ? 3. Nous avons l'habitude d'aller au marché le samedi. 4. Cette affaire ne concerne pas les élèves. 5. Il pleut souvent à Douala en août.
\end{exercice}

\begin{corrige}[Exercice « À toi de jouer »]
1. \esp{Me gustan las clases de español.} (chose au pluriel : \esp{gustan}). 2. \esp{¿Te duelen los ojos?} 3. \esp{Solemos ir al mercado los sábados.} (\esp{soler} se conjugue ici à la 1\iere{} personne du pluriel). 4. \esp{Ese asunto no atañe a los alumnos.} 5. \esp{Suele llover en Douala en agosto.}
\end{corrige}

\coursec{Règles et emplois des verbes défectifs}

Dans la section précédente, nous avons \emph{observé} des phrases comme \esp{Me gusta el debate} ou \esp{A los alumnos les duelen las piernas} et nous avons remarqué quelque chose d'étrange : le sujet français (\emph{je}, \emph{ils}) semble avoir disparu, et le verbe espagnol se conjugue avec la \emph{chose} et non avec la personne. Il est temps maintenant de transformer ces observations en \textbf{règles rigoureuses} : c'est la condition pour traduire sans faute au Bac.

\courssub{La structure générale du verbe défectif}

\begin{definition}[Verbe défectif]
On appelle \cle{verbe défectif} (ou verbe à construction inversée) un verbe espagnol dont le \textbf{sujet grammatical est la chose} qui provoque le sentiment, la douleur ou l'opinion, tandis que la \textbf{personne} qui ressent est exprimée par un \textbf{pronom complément d'objet indirect (COI)}. Le français dit « \emph{quelqu'un aime quelque chose} » ; l'espagnol dit littéralement « \emph{quelque chose plaît à quelqu'un} ».
\end{definition}

La construction complète comporte quatre éléments, dans cet ordre logique :

\begin{center}
\begin{tabularx}{\linewidth}{|X|X|X|X|}
\hline
\rowcolor{popblueL} \textbf{(a + nom/pronom tonique)} & \textbf{Pronom COI} & \textbf{Verbe (3\up{e} pers.)} & \textbf{Sujet réel} \\
\hline
A mí & me & gusta & el arroz \\
\hline
A ti & te & duelen & las muelas \\
\hline
A Juan & le & interesa & el debate \\
\hline
A nosotros & nos & parece & difícil \\
\hline
A los alumnos & les & encanta & el fútbol \\
\hline
\end{tabularx}
\end{center}

\begin{propriete}[Règle d'accord : le verbe s'accorde avec la CHOSE]
Le verbe défectif se conjugue à la \textbf{3\up{e} personne} et s'accorde avec le \textbf{sujet réel} (la chose), placé le plus souvent \emph{après} le verbe :
\begin{itemize}
\item sujet singulier $\rightarrow$ \esp{me gusta el arroz} « j'aime le riz » ;
\item sujet pluriel $\rightarrow$ \esp{me gustan las frituras} « j'aime les beignets » ;
\item devant un \textbf{infinitif}, le verbe est \textbf{toujours au singulier} : \esp{me gusta leer} « j'aime lire » ; \esp{¿Te gusta bailar?} « Tu aimes danser ? ».
\end{itemize}
\textbf{Exception apparente :} avec \esp{gustar} au sens affectif entre personnes, le sujet reste la personne « aimée » : \esp{Tú me gustas} « Tu me plais » (litt. « tu plais à moi »).
\end{propriete}

\begin{attention}[L'erreur n° 1 du francophone]
Ne jamais traduire « j'aime beaucoup » par \esp{me gusto mucho} : cela signifierait « je me plais beaucoup à moi-même » ! On dit \esp{me gusta mucho}. De même, « nous aimons le cacao » se traduit \esp{nos gusta el cacao} et jamais \esp{nos gustamos el cacao}. Retenez le réflexe : le pronom COI (\esp{me, te, le, nos, os, les}) n'est \textbf{jamais} un sujet.
\end{attention}

\courssub{Gustar et sa famille}

\esp{Gustar} n'est pas seul : toute une famille de verbes de sentiment suit exactement la même construction.

\begin{center}
\begin{tabularx}{\linewidth}{|X|X|X|}
\hline
\rowcolor{popblueL} \textbf{Verbe} & \textbf{Français} & \textbf{Exemple} \\
\hline
gustar & aimer, plaire & Me gusta el arroz. = J'aime le riz. \\
\hline
encantar & adorer & Le encanta viajar. = Il adore voyager. \\
\hline
interesar & intéresser & Nos interesa la historia. = L'histoire nous intéresse. \\
\hline
apetecer & avoir envie de & ¿Te apetece un zumo? = Tu as envie d'un jus ? \\
\hline
importar & importer, soucier & No me importa. = Ça m'est égal. \\
\hline
fascinar & passionner & Me fascinan las lenguas. = Les langues me passionnent. \\
\hline
molestar & déranger, gêner & Le molesta el ruido. = Le bruit le dérange. \\
\hline
\end{tabularx}
\end{center}

\begin{exemplebox}[Exemples traduits]
\begin{itemize}
\item \esp{A los alumnos les interesa el debate.} « Les élèves s'intéressent au débat. »
\item \esp{A mi hermana le encantan los bendskins.} « Ma sœur adore les motos-taxis. »
\item \esp{¿Os apetece cenar fuera esta noche?} « Vous avez envie de dîner dehors ce soir ? »
\item \esp{A nosotros no nos importa esperar.} « Attendre ne nous dérange pas. »
\end{itemize}
\end{exemplebox}

\begin{methode}[Traduire « j'aime X » en quatre étapes]
\begin{enumerate}
\item \textbf{Identifier X} : est-ce un nom singulier, un nom pluriel ou un infinitif ?
\item \textbf{Choisir la forme du verbe} : \esp{gusta} si X est singulier ou un infinitif ; \esp{gustan} si X est pluriel.
\item \textbf{Placer le pronom COI} correspondant à la personne française : je $\rightarrow$ \esp{me}, tu $\rightarrow$ \esp{te}, il/elle $\rightarrow$ \esp{le}, nous $\rightarrow$ \esp{nos}, vous $\rightarrow$ \esp{os}, ils $\rightarrow$ \esp{les}.
\item \textbf{Ajouter \esp{a + pronom tonique}} pour insister ou lever l'ambiguïté : \esp{a mí me gusta}, \esp{a ti te gusta}, \esp{a Juan le gusta}.
\end{enumerate}
\end{methode}

\begin{attention}[Le redoublement obligatoire ou utile]
Le pronom COI \esp{le} / \esp{les} est ambigu (à lui ? à elle ? à eux ?). Il faut alors préciser : \esp{A Juan le gusta el fútbol} « Jean aime le football ». Avec un nom de personne, le \esp{a} est \textbf{obligatoire}. En revanche, \esp{a mí, a ti, a nosotros} sert seulement à insister : \esp{A mí no me gusta, pero a ti sí} « Moi je n'aime pas, mais toi oui ». Attention à l'accent : \esp{mí} (pronom) avec accent, \esp{mi} (possessif) sans accent.
\end{attention}

\begin{popfigure}
\begin{center}
\begin{tikzpicture}[every node/.style={font=\small}]
\node[draw=popblue, fill=popblueL, rounded corners, align=center, minimum width=5.2cm] (fr) at (0,0) {\textbf{Français}\\ je / tu / il \ldots\\ \textbf{aime(nt)} quelque chose};
\node[draw=poporange, fill=poporangeL, rounded corners, align=center, minimum width=5.2cm] (es) at (8.6,0) {\textbf{Espagnol}\\ me / te / le \ldots\\ \textbf{gusta(n)} + sujet réel};
\draw[-{Stealth}, very thick, poppurple] (fr) -- node[above, align=center, font=\footnotesize]{inversion\\ sujet $\leftrightarrow$ COI} (es);
\node[align=center, font=\footnotesize, popmuted] at (4.3,-1.4) {« J'aime le riz » $\rightarrow$ \esp{Me gusta el arroz} (litt. « le riz plaît à moi »)};
\end{tikzpicture}
\end{center}
\legende{Transformation français $\rightarrow$ espagnol : le sujet français devient COI, la chose devient sujet.}
\end{popfigure}

\courssub{Doler : la douleur et l'accord avec la partie du corps}

\esp{Doler} (radical \esp{duel-}, comme \esp{dormir} $\rightarrow$ \esp{duermo}) signifie « faire mal ». Le français dit « j'ai mal à la tête » ; l'espagnol dit littéralement « la tête me fait mal ».

\begin{exemplebox}[Doler en action]
\begin{itemize}
\item \esp{Me duele la cabeza.} « J'ai mal à la tête. » (singulier)
\item \esp{Me duelen las muelas.} « J'ai mal aux dents. » (pluriel)
\item \esp{Nos duele la espalda.} « Nous avons mal au dos. »
\item \esp{¿Te duele la garganta?} « Tu as mal à la gorge ? »
\item \esp{Ayer me dolía el estómago.} « Hier, j'avais mal au ventre. »
\end{itemize}
\end{exemplebox}

\begin{attention}[Article défini, pas de possessif]
En espagnol, la partie du corps prend l'\textbf{article défini} (\esp{la cabeza}, \esp{las muelas}) et jamais le possessif : on dit \esp{me duele la cabeza}, pas \esp{me duele mi cabeza}. Le pronom COI indique déjà à qui appartient la tête ! Même logique que dans \esp{me lavo las manos} « je me lave les mains ».
\end{attention}

\courssub{Parecer : l'opinion}

\esp{Parecer} (« sembler, paraître ») suit la même construction et sert à exprimer l'opinion :

\begin{itemize}
\item \esp{Me parece que el examen es fácil.} « Il me semble que l'examen est facile. » (+ \textbf{indicatif} quand on affirme)
\item \esp{Me parece interesante.} « Je trouve cela intéressant. » (\esp{parecer} + adjectif)
\item \esp{No me parece que sea justo.} « Il ne me semble pas que ce soit juste. » (négation + \textbf{subjonctif}, complément utile au Bac)
\end{itemize}

\begin{propriete}[Parecer : indicatif ou subjonctif ?]
Après \esp{me parece que} (opinion affirmée), on emploie l'\textbf{indicatif} ; après \esp{no me parece que} (opinion niée), le subordonné passe au \textbf{subjonctif}, comme après \esp{no creo que} : \esp{No me parece que tengas razón.} « Je ne pense pas que tu aies raison. »
\end{propriete}

\courssub{Soler + infinitif : l'habitude}

\esp{Soler} (radical \esp{suel-}) suivi d'un infinitif traduit « avoir l'habitude de, faire habituellement ». Attention : c'est un verbe à construction \emph{normale} (le sujet est la personne), mais il est \textbf{défectif en français} car il n'a pas d'équivalent verbal direct.

\begin{exemplebox}[Soler]
\begin{itemize}
\item \esp{Suelo estudiar por la noche.} « J'ai l'habitude d'étudier le soir. »
\item \esp{Solíamos ir al mercado los sábados.} « Nous avions l'habitude d'aller au marché le samedi. »
\item \esp{¿Sueles levantarte temprano?} « Tu as l'habitude de te lever tôt ? »
\end{itemize}
\end{exemplebox}

\begin{attention}[Soler : un verbe vraiment défectif]
\esp{Soler} ne s'emploie pratiquement qu'au \textbf{présent} (\esp{suelo, sueles, suele\ldots}) et à l'\textbf{imparfait} (\esp{solía, solías, solía\ldots}). Les autres temps (passé simple, futur) sont théoriques et quasiment jamais utilisés : au Bac, limitez-vous au présent et à l'imparfait. Pour « j'avais l'habitude de », dites \esp{solía} + infinitif, et non un passé simple.
\end{attention}

\courssub{Atañer et concernir : « concerner »}

Ces deux verbes, fréquents dans les textes de vie citoyenne, n'existent pratiquement qu'à la 3\up{e} personne :

\begin{itemize}
\item \esp{Eso no te atañe.} « Cela ne te concerne pas. »
\item \esp{Eso no me concierne.} « Cela ne me concerne pas. »
\item \esp{En lo que concierne a la educación\ldots} « En ce qui concerne l'éducation\ldots » (formule de dissertation très appréciée au Bac)
\end{itemize}

\courssub{Les verbes impersonnels : le temps qu'il fait}

Les verbes météorologiques n'ont qu'une seule forme : la 3\up{e} personne du singulier, sans sujet.

\begin{center}
\begin{tabularx}{\linewidth}{|X|X|}
\hline
\rowcolor{popblueL} \textbf{Espagnol} & \textbf{Français} \\
\hline
Llueve / Llovía / Llovió & Il pleut / Il pleuvait / Il a plu \\
\hline
Nieva & Il neige \\
\hline
Truena & Il tonne \\
\hline
Hace calor / Hace frío / Hace viento & Il fait chaud / Il fait froid / Il y a du vent \\
\hline
Hay que + infinitif & Il faut + infinitif \\
\hline
\end{tabularx}
\end{center}

\begin{exemplebox}[Exemples traduits]
\begin{itemize}
\item \esp{Llovía cuando salimos.} « Il pleuvait quand nous sommes sortis. »
\item \esp{En diciembre nieva en los Pirineos.} « En décembre, il neige dans les Pyrénées. »
\item \esp{Hay que respetar las leyes.} « Il faut respecter les lois. »
\end{itemize}
\end{exemplebox}

\courssub{Tableau récapitulatif des formes}

\begin{center}
\begin{tabular}{|l|l|l|l|}
\hline
\rowcolor{popblueL} \textbf{Verbe} & \textbf{Présent} & \textbf{Imparfait} & \textbf{Passé simple} \\
\hline
gustar & gusta / gustan & gustaba / gustaban & gustó / gustaron \\
\hline
doler & duele / duelen & dolía / dolían & dolió / dolieron \\
\hline
soler & suele / suelen & solía / solían & (inusité) \\
\hline
llover & llueve & llovía & llovió \\
\hline
\end{tabular}
\end{center}

\begin{popfigure}
\begin{center}
\begin{tikzpicture}[font=\small]
\draw[-{Stealth}, thick, popdark] (0,0) -- (10.6,0);
\foreach \x/\t/\f in {0.8/{passé simple}/gustó, 3.6/{imparfait}/gustaba, 6.4/{présent}/gusta, 9.2/{futur}/gustará}{
  \fill[poporange] (\x,0) circle (2.2pt);
  \node[above, align=center, font=\footnotesize] at (\x,0.1) {\t\\\esp{\f}};
}
\node[below, font=\footnotesize, popmuted] at (5.3,-0.25) {3\up{e} personne du singulier de \esp{gustar} : \esp{me gustó}, \esp{me gustaba}, \esp{me gusta}, \esp{me gustará}};
\end{tikzpicture}
\end{center}
\legende{Frise des temps de gustar (3\up{e} pers. sing.) : du passé simple au futur.}
\end{popfigure}

\courssub{Comparaison systématique français / espagnol}

\begin{center}
\begin{tabularx}{\linewidth}{|X|X|X|}
\hline
\rowcolor{popblueL} \textbf{Français} & \textbf{Espagnol} & \textbf{Remarque} \\
\hline
J'aime le riz. & Me gusta el arroz. & sujet fr. $\rightarrow$ COI esp. \\
\hline
Tu aimes les langues. & Te gustan las lenguas. & pluriel $\rightarrow$ gustan \\
\hline
Il adore voyager. & Le encanta viajar. & infinitif $\rightarrow$ singulier \\
\hline
Nous avons mal au dos. & Nos duele la espalda. & article défini, pas de possessif \\
\hline
Ils s'intéressent au débat. & Les interesa el debate. & redoublement possible : a ellos \\
\hline
Il pleuvait. & Llovía. & impersonnel : sans sujet \\
\hline
\end{tabularx}
\end{center}

\begin{exoresolu}[Thème guidé : mettez la règle en application]
Traduisez en espagnol : « Mes camarades aiment beaucoup le football, mais moi, j'ai mal aux jambes après le match. D'habitude, il pleut le samedi à Yaoundé, alors nous avons l'habitude de jouer le dimanche. »
\tcblower
\textbf{Solution détaillée.}
\begin{enumerate}
\item « Mes camarades aiment le football » : sujet français pluriel = COI \esp{les} ; la chose aimée (\esp{el fútbol}) est singulière $\rightarrow$ \esp{gusta} ; redoublement obligatoire : \esp{A mis compañeros les gusta mucho el fútbol}.
\item « moi, j'ai mal aux jambes » : \esp{las piernas} est pluriel $\rightarrow$ \esp{duelen} : \esp{pero a mí me duelen las piernas después del partido}.
\item « il pleut le samedi » : impersonnel : \esp{Normalmente llueve los sábados en Yaundé}.
\item « nous avons l'habitude de jouer » : \esp{soler} + infinitif : \esp{así que solemos jugar los domingos}.
\end{enumerate}
\textbf{Traduction complète :} \esp{A mis compañeros les gusta mucho el fútbol, pero a mí me duelen las piernas después del partido. Normalmente llueve los sábados en Yaundé, así que solemos jugar los domingos.}
\end{exoresolu}

\begin{exercice}[À vous de jouer]
Traduisez en espagnol :
\begin{enumerate}
\item Mes parents adorent la musique makossa.
\item Tu as mal à la tête ? Repose-toi un peu.
\item Il ne me semble pas que cette loi soit juste.
\item Quand j'étais enfant, j'avais l'habitude d'aller au marché avec ma grand-mère.
\item En ce qui concerne les transports, les motos-taxis ne nous intéressent pas.
\end{enumerate}
\end{exercice}

\begin{corrige}[Exercice : À vous de jouer]
\begin{enumerate}
\item \esp{A mis padres les encanta la música makossa.} (redoublement \esp{a mis padres} obligatoire ; \esp{la música} singulier $\rightarrow$ \esp{encanta})
\item \esp{¿Te duele la cabeza? Descansa un poco.} (article défini \esp{la cabeza} ; impératif \esp{descansa})
\item \esp{No me parece que esta ley sea justa.} (opinion niée $\rightarrow$ subjonctif \esp{sea})
\item \esp{Cuando era niño, solía ir al mercado con mi abuela.} (habitude passée $\rightarrow$ imparfait \esp{solía})
\item \esp{En lo que concierne a los transportes, los bendskins no nos interesan.} (\esp{los bendskins} pluriel $\rightarrow$ \esp{interesan})
\end{enumerate}
\end{corrige}

\begin{savaistu}[Gustar et doler : les stars du Bac]
Au Bac, \esp{gustar} et \esp{doler} apparaissent presque chaque année dans le thème : les correcteurs vérifient systématiquement l'accord \esp{gusta/gustan}, \esp{duele/duelen} et la présence du pronom COI. De même, \esp{soler} + infinitif est la traduction attendue de « avoir l'habitude de » : le calque \esp{tener la costumbre de} est accepté mais moins élégant. Sachez enfin qu'en Guinée équatoriale, pays hispanophone voisin du Cameroun, on entend partout \esp{¿Te gusta el malamba?} : la construction inversée y est aussi vivante qu'en Espagne !
\end{savaistu}

\coursec{La phrase restrictive}

Après avoir étudié les verbes défectifs, ces verbes « à construction spéciale » qui obligent à repenser la place du sujet, nous abordons un deuxième obstacle majeur de la traduction : la \cle{restriction}. En français, elle s'exprime presque toujours par la petite structure « ne... que » (\emph{Je n'ai que cinq mille francs}). En espagnol, il n'existe pas une mais plusieurs structures concurrentes, et choisir la mauvaise est une faute éliminatoire au Baccalauréat. Cette section vous donne une méthode de décision infaillible.

\courssub{Observation : une même idée, deux structures}

\begin{textebox}[Au marché de Mokolo — dialogue]
— ¿Cuánto dinero llevas para el mercado, Amina?\\
— \textbf{Sólo tengo cinco mil francos}. Con eso compro únicamente tomates y cebollas.\\
— ¿Y no quieres comprar también pescado?\\
— \textbf{No tengo más que cinco mil francos}, te digo. Además, el vendedor \textbf{no es camerunés, sino nigeriano}, y \textbf{no vende pescado, sino carne}.\\
— Pues yo \textbf{apenas tengo mil francos}. ¡Compraremos poco hoy!\\
— No importa: \textbf{no trabajamos los lunes, sino los sábados}, así que volveremos.
\end{textebox}

\begin{definition}[Glossaire du dialogue]
\esp{llevar} : avoir sur soi ; \esp{cebollas} : oignons ; \esp{el vendedor} : le vendeur ; \esp{así que} : donc ; \esp{volveremos} : nous reviendrons (futur de \esp{volver}) ; \esp{únicamente} : uniquement ; \esp{también} : aussi.
\end{definition}

\begin{definition}[La restriction]
La \cle{restriction} est un procédé qui limite la portée d'un énoncé à un seul élément (une personne, une quantité, un jour, une action). Le français utilise « ne... que » ou « seulement ». L'espagnol dispose de \textbf{deux familles de formes} :
\begin{itemize}
\item des formes \textbf{positives} (phrase affirmative) : \esp{sólo}, \esp{únicamente}, \esp{solamente} ;
\item des formes \textbf{négatives} (phrase négative) : \esp{no... más que} et \esp{no... sino}.
\end{itemize}
\end{definition}

Observons les deux premières répliques d'Amina : \esp{Sólo tengo cinco mil francos} et \esp{No tengo más que cinco mil francos}. Le sens est identique (« Je n'ai que cinq mille francs »), mais la construction change : phrase affirmative avec \esp{sólo}, ou phrase négative avec \esp{no... más que}. C'est la première grande leçon de cette section : \textbf{le français « ne... que » a toujours deux traductions possibles en espagnol pour exprimer une quantité limitée}.

\courssub{La restriction positive : sólo, únicamente, solamente}

\begin{propriete}[Forme et place de sólo]
\esp{Sólo} (et ses synonymes plus soutenus \esp{únicamente}, \esp{solamente}) est un adverbe qui se place \textbf{immédiatement devant l'élément restreint}, ou devant le verbe quand c'est l'action entière qui est limitée. La phrase reste affirmative : pas de \esp{no}.
\begin{itemize}
\item \esp{Sólo quedan dos días.} \esp{« Il ne reste que deux jours. »} (restriction sur le sujet \emph{dos días})
\item \esp{Trabajo sólo los lunes.} \esp{« Je ne travaille que le lundi. »} (restriction sur le complément \emph{los lunes})
\item \esp{Únicamente los alumnos de Terminale pueden entrar.} \esp{« Seuls les élèves de Terminale peuvent entrer. »} (restriction sur le sujet, ton formel)
\end{itemize}
\end{propriete}

\begin{attention}[La place change le sens]
Comme en français, la place de la restriction modifie le sens de la phrase. Comparez :
\begin{itemize}
\item \esp{Sólo los lunes trabajo.} \esp{« Uniquement le lundi, je travaille »} (pas les autres jours).
\item \esp{Los lunes sólo trabajo.} \esp{« Le lundi, je ne fais que travailler »} (rien d'autre ce jour-là).
\end{itemize}
En traduction, repérez d'abord \textbf{quel élément} le « que » français restreint, puis placez \esp{sólo} juste devant cet élément.
\end{attention}

\courssub{La restriction négative : no... más que}

\begin{propriete}[no... más que = « ne... que » (quantité limitée)]
La structure \esp{no + verbe + más que + élément} exprime une quantité ou un ensemble limité. Elle est parfaitement interchangeable avec \esp{sólo} / \esp{únicamente} :
\begin{itemize}
\item \esp{No tengo más que una hermana.} \esp{« Je n'ai qu'une sœur. »} (= \esp{Sólo tengo una hermana.})
\item \esp{No leo más que novelas.} \esp{« Je ne lis que des romans. »}
\item \esp{No quedan más que dos días.} \esp{« Il ne reste que deux jours. »}
\end{itemize}
Attention à l'orthographe : \esp{más} (avec accent) signifie « plus » (quantité) ; \esp{mas} (sans accent) signifie « mais » (conjonction littéraire, synonyme de \esp{sino}).
\end{propriete}

\courssub{La restriction d'opposition : no... sino}

\begin{propriete}[no... sino = « non pas... mais »]
Quand la phrase ne limite pas une quantité mais \textbf{corrige ou remplace une affirmation} (opposition), l'espagnol impose \esp{no... sino} (et jamais \esp{más que}).
\begin{itemize}
\item Devant un \textbf{nom, pronom, adjectif, infinitif} : \esp{sino}.
\item Devant un \textbf{verbe conjugué} (subordonnée) : \esp{sino que}.
\end{itemize}
Exemples :
\begin{itemize}
\item \esp{No es médico, sino enfermero.} \esp{« Il n'est pas médecin, mais infirmier. »} (noms)
\item \esp{No vino a trabajar, sino a descansar.} \esp{« Il n'est pas venu travailler, mais se reposer. »} (infinitifs)
\item \esp{No estudia, sino que trabaja.} \esp{« Il n'étudie pas, mais il travaille. »} (verbes conjugués)
\item \esp{No quiero que vengas, sino que llames.} \esp{« Je ne veux pas que tu viennes, mais que tu appelles. »} (subjonctif après verbes de volonté)
\end{itemize}
\end{propriete}

\begin{aretenir}[Le critère de décision infaillible]
Posez-vous la question : le français « ne... que » peut-il être remplacé par « \emph{mais} » sans changer le sens ?
\begin{itemize}
\item Si \textbf{OUI} (« Il n'est pas médecin, \emph{mais} infirmier ») $\rightarrow$ \textbf{Opposition} $\rightarrow$ \esp{no... sino (que)}.
\item Si \textbf{NON} (« Je n'ai \emph{que} cinq mille francs » — on ne dit pas « mais cinq mille francs ») $\rightarrow$ \textbf{Quantité limitée} $\rightarrow$ \esp{sólo} ou \esp{no... más que}.
\end{itemize}
\end{aretenir}

\begin{popfigure}
\begin{center}
\begin{tikzpicture}[
  node distance=0.9cm and 1.6cm,
  startstop/.style={rectangle, rounded corners, draw=popblue, fill=popblueL, align=center, font=\small, inner sep=6pt},
  decision/.style={diamond, draw=poporange, fill=poporangeL, align=center, font=\small, inner sep=2pt, aspect=2.2},
  result/.style={rectangle, draw=popgreen, fill=popgreenL, align=center, font=\small, inner sep=6pt},
  arrow/.style={-{Stealth[length=3mm]}, thick, popdark}
]
\node[startstop, align=center] (start){Le français dit\\ « ne... que » ?};
\node[decision, below=of start, align=center] (test){Peut-on dire\\ « mais » ?};
\node[result, below left=1.1cm and 0.4cm of test, align=center] (quant){Quantité limitée\\ \esp{sólo} ou \esp{no... más que}\\ \esp{No tengo más que cinco mil francos.}};
\node[result, below right=1.1cm and 0.4cm of test, align=center] (oppos){Opposition / correction\\ \esp{no... sino} (+ nom, infinitif)\\ \esp{sino que} (+ verbe conjugué)\\ \esp{No es médico, sino enfermero.}};
\draw[arrow] (start) -- (test);
\draw[arrow] (test) -| node[above left, font=\small]{non} (quant);
\draw[arrow] (test) -| node[above right, font=\small]{oui} (oppos);
\end{tikzpicture}
\end{center}
\legende{Organigramme de décision : choisir entre \esp{más que} et \esp{sino}}
\end{popfigure}

\begin{popfigure}
\begin{center}
\begin{tikzpicture}[
  box/.style={rectangle, rounded corners, draw=popblue, fill=popblueL, align=center, font=\small, inner sep=6pt},
  box2/.style={rectangle, rounded corners, draw=poppurple, fill=poppurpleL, align=center, font=\small, inner sep=6pt},
  arrow/.style={-{Stealth[length=3mm]}, thick, popdark}
]
\node[box] (no) {\esp{NO} + verbe};
\node[box2, right=1.6cm of no, align=center] (link){\esp{más que}\\ ou \esp{sino (que)}};
\node[box, right=1.6cm of link, align=center] (elem){élément restreint\\ (nom, quantité, infinitif,\\ proposition)};
\draw[arrow] (no) -- (link);
\draw[arrow] (link) -- (elem);
\node[below=0.5cm of link, align=center, font=\small\itshape, text=popmuted] {No vino \textbf{más que} con cien francos. / No vino a trabajar, \textbf{sino} a descansar.};
\end{tikzpicture}
\end{center}
\legende{Schéma de la structure négative : \esp{no} + verbe ... \esp{más que / sino} + élément restreint}
\end{popfigure}

\courssub{Apenas : la restriction d'intensité}

\begin{definition}[Apenas]
\esp{Apenas} signifie « à peine » et limite l'intensité, la quantité ou le degré. Il se place devant le verbe.
\begin{itemize}
\item \esp{Apenas habla español.} \esp{« Il parle à peine espagnol. »}
\item \esp{Apenas tengo tiempo.} \esp{« J'ai à peine le temps. »}
\item \esp{Apenas comimos nada.} \esp{« Nous avons à peine mangé quelque chose. »} (valeur négative)
\end{itemize}
Au passé, \esp{apenas + pretérito indefinido} correspond à la construction française « à peine... que » (antériorité immédiate) :
\esp{Apenas llegó, empezó a llover.} \esp{« À peine fut-il arrivé qu'il se mit à pleuvoir. »} (point de programme Bac : corrélation temporelle).
\end{definition}

\courssub{Tableau comparatif français / espagnol}

\begin{center}
\begin{tabularx}{\linewidth}{|X|X|X|}
\hline
\rowcolor{popblueL} Français & Español & Remarque \\
\hline
Je n'ai que cinq mille francs. & Sólo tengo cinco mil francos. / No tengo más que cinco mil francos. & Deux traductions équivalentes (quantité) \\
\hline
Il ne reste que deux jours. & Sólo quedan dos días. / No quedan más que dos días. & Quantité limitée \\
\hline
Il n'est pas médecin, mais infirmier. & No es médico, sino enfermero. & Opposition : \esp{sino} + nom \\
\hline
Il n'étudie pas, mais il travaille. & No estudia, sino que trabaja. & Opposition : \esp{sino que} + indicatif \\
\hline
Je ne travaille que le lundi. & Sólo trabajo los lunes. / No trabajo más que los lunes. & \esp{sólo} devant l'élément restreint \\
\hline
Il parle à peine espagnol. & Apenas habla español. & Restriction d'intensité \\
\hline
À peine arrivé, il partit. & Apenas llegó, se fue. & Valeur temporelle (Bac) \\
\hline
\end{tabularx}
\end{center}

\begin{methode}[Traduire une phrase restrictive : 3 étapes]
\begin{enumerate}
\item \textbf{Identifiez l'élément restreint} par le « que » français (sujet, COD, complément de temps, verbe...).
\item \textbf{Appliquez le test du « mais »} : peut-on remplacer « ne... que » par « non pas... mais » ?
\item \textbf{Choisissez la structure} :
\begin{itemize}
\item Test \textbf{NON} (quantité) $\rightarrow$ \esp{Sólo + [élément] + verbe} \textbf{OU} \esp{No + verbe + más que + [élément]}.
\item Test \textbf{OUI} (opposition) $\rightarrow$ \esp{No + verbe + sino (+ que) + [élément corrigé]}.
\item Cas \textbf{Apenas} $\rightarrow$ si le sens est « à peine » (peu, difficilement) ou « à peine... que » (temps).
\end{itemize}
\end{enumerate}
\end{methode}

\begin{exoresolu}[Thème : la restriction au quotidien]
Traduisez en espagnol :
\begin{enumerate}
\item « Je n'ai qu'une sœur. »
\item « Elle n'est pas venue pour discuter, mais pour travailler. »
\item « Nous ne sortons que le samedi. »
\item « Il n'écrit pas des lettres, mais il envoie des messages. »
\end{enumerate}
\tcblower
1. \textbf{Quantité limitée} (test « mais » impossible) : \esp{Sólo tengo una hermana.} ou \esp{No tengo más que una hermana.} \\
2. \textbf{Opposition} sur des infinitifs (test « mais » possible) : \esp{No ha venido para discutir, sino para trabajar.} \\
3. \textbf{Quantité limitée} sur le complément de temps : \esp{Sólo salimos los sábados.} ou \esp{No salimos más que los sábados.} \\
4. \textbf{Opposition} sur deux verbes conjugués : \esp{No escribe cartas, sino que envía mensajes.} (attention à la diphtongaison \esp{enviar} $\rightarrow$ \esp{envía}, 3\textsuperscript{e} personne sg. présent).
\end{exoresolu}

\begin{attention}[Fautes éliminatoires au Bac]
\begin{itemize}
\item Ne traduisez \textbf{jamais} « ne... que » par \esp{no... que} seul : \esp{*No tengo que cinco mil francos} est une faute grave. Il faut \esp{no... más que} ou \esp{sólo}.
\item Ne confondez pas \esp{sólo que} en début de subordonnée, qui signifie « sauf que / seulement » : \esp{Iría contigo, sólo que estoy enfermo.} \esp{« J'irais avec toi, sauf que je suis malade. »}
\item Devant un verbe conjugué, \esp{sino} seul est impossible : on dit \esp{No estudia, sino que trabaja}, jamais \esp{*sino trabaja}.
\item Oublier l'accent sur \esp{más} (plus) dans \esp{no... más que} : \esp{*No tengo mas que...} change le sens en « Je n'ai pas mais que... » (absurde).
\item Utiliser l'article indéfini devant un nom de profession/identité après \esp{sino} : \esp{*No es un médico, sino un enfermero} $\rightarrow$ Correct : \esp{No es médico, sino enfermero}.
\end{itemize}
\end{attention}

\begin{savaistu}[¿Sólo ou solo ?]
Selon la Real Academia Española (RAE, 2010), l'accent sur \esp{sólo} (adverbe, « seulement ») est aujourd'hui \textbf{facultatif} : on peut écrire \esp{solo} dans tous les cas pour éviter l'ambiguïté graphique. Cependant, à l'écrit académique et surtout à l'examen du Baccalauréat, l'accent reste \textbf{fortement recommandé} pour lever l'ambiguïté immédiate avec l'adjectif \esp{solo} (« seul, unique ») : comparez \esp{Está solo} \esp{« Il est seul »} et \esp{Sólo está} \esp{« Il est seulement là »}.
\end{savaistu}

\begin{exercice}[À vous de traduire]
Traduisez en espagnol :
\begin{enumerate}
\item « Il ne reste que trois places. »
\item « Ce n'est pas un professeur, mais un étudiant. »
\item « Je ne lis que des romans policiers. »
\item « Elle ne chante pas, mais elle danse. »
\item « Nous avons à peine le temps de manger. »
\item « À peine eut-il fini, il sortit. » (indice : \esp{apenas} + pretérito)
\end{enumerate}
\end{exercice}

\begin{corrige}[Exercice : la restriction]
\begin{enumerate}
\item \esp{Sólo quedan tres plazas.} ou \esp{No quedan más que tres plazas.} (quantité)
\item \esp{No es profesor, sino estudiante.} (opposition, noms sans article après \esp{ser})
\item \esp{Sólo leo novelas policíacas.} ou \esp{No leo más que novelas policíacas.} (quantité sur COD)
\item \esp{No canta, sino que baila.} (opposition, verbe conjugué $\rightarrow$ \esp{sino que})
\item \esp{Apenas tenemos tiempo de comer.} (restriction d'intensité)
\item \esp{Apenas terminó, salió.} (valeur temporelle, \esp{pretérito indefinido})
\end{enumerate}
\end{corrige}

\textbf{Transition.} Maîtriser la restriction, c'est éviter les pièges du « ne... que » français. Le troisième pilier de cette leçon de traduction, le \textbf{discours indirect}, mobilisera la même rigueur : il s'agira de transposer les personnes, les temps et les indicateurs spatio-temporels sans trahir le sens original. C'est l'objet de la section suivante.

\coursec{Du discours direct au discours indirect : principes}

Dans la section précédente, nous avons appris à traduire la restriction (\esp{sólo}, \esp{no... más que}, \esp{apenas}), c'est-à-dire à \emph{limiter} ce que dit une phrase. Nous allons maintenant apprendre à \emph{rapporter} ce que dit quelqu'un. En effet, dans la presse, dans un récit ou dans un commentaire de texte — exercices centraux de l'épreuve de Terminale A —, on ne cite presque jamais les paroles mot à mot : on les \textbf{rapporte}. C'est l'objet du \cle{discours indirect}, en espagnol \esp{estilo indirecto}.

\courssub{Observation : une paire de phrases révélatrice}

Observons attentivement les deux phrases suivantes :

\begin{center}
\esp{Dijo: «Vendré mañana».} \quad $\longrightarrow$ \quad \esp{Dijo que vendría al día siguiente.}
\end{center}

« Il a dit : "Je viendrai demain". » $\longrightarrow$ « Il a dit qu'il viendrait le lendemain. »

\begin{exemplebox}[Qu'est-ce qui a changé ?]
Comparons mot à mot :
\begin{itemize}
\item \esp{Vendré} (futur, 1\iere{} personne) devient \esp{vendría} (conditionnel, 3\ieme{} personne) : la \textbf{personne} et le \textbf{temps} changent.
\item \esp{mañana} (demain) devient \esp{al día siguiente} (le lendemain) : le \textbf{repère temporel} change.
\item Les guillemets et les deux-points disparaissent ; la parole rapportée est introduite par \esp{que} : la \textbf{structure} change.
\end{itemize}
Le sens reste identique, mais tout a été \emph{transposé} parce que le locuteur qui rapporte n'est plus celui qui a parlé, et que le moment du rapport n'est plus celui de la parole.
\end{exemplebox}

\begin{textebox}[Un anuncio del profesor]
\textbf{Style direct :} El profesor anunció: «El examen será mañana aquí, en esta sala. Traed vuestros cuadernos y no lleguéis tarde.»

\textbf{Style indirect :} El profesor anunció que el examen sería al día siguiente allí, en aquella sala. Mandó que llevaran sus cuadernos y que no llegaran tarde.

\emph{Traduction :} Le professeur a annoncé : « L'examen aura lieu demain ici, dans cette salle. Apportez vos cahiers et n'arrivez pas en retard. » / Le professeur a annoncé que l'examen aurait lieu le lendemain là-bas, dans cette salle-là. Il a ordonné qu'ils apportent leurs cahiers et qu'ils n'arrivent pas en retard.

\textbf{Glossaire :} \esp{anunciar} = annoncer ; \esp{la sala} = la salle ; \esp{traer / llevar} = apporter / emporter ; \esp{llegar tarde} = arriver en retard ; \esp{mandar} = ordonner.
\end{textebox}

\textbf{Questions d'observation :} 1) Relevez dans le texte trois mots qui ont changé entre le style direct et le style indirect. 2) Quel mot introduit la déclaration rapportée ? 3) Comment l'impératif \esp{traed} est-il rapporté ?

\begin{definition}[Le discours indirect (estilo indirecto)]
Le \cle{discours indirect} rapporte les paroles d'un locuteur \textbf{sans les citer mot à mot} : il les intègre dans une \textbf{proposition subordonnée} dépendant d'un verbe introducteur (\esp{decir}, \esp{preguntar}, \esp{mandar}...). Cette intégration impose une \textbf{adaptation} de trois catégories d'éléments :
\begin{enumerate}
\item les \textbf{personnes} grammaticales et les possessifs ;
\item les \textbf{temps} verbaux (concordance des temps) ;
\item les \textbf{repères spatio-temporels} (déictiques : \esp{hoy}, \esp{aquí}, \esp{este}...).
\end{enumerate}
La subordonnée est introduite par \esp{que} (déclaration), par \esp{si} (question fermée, sans mot interrogatif) ou par un \textbf{mot interrogatif} (\esp{quién}, \esp{dónde}, \esp{cuándo}...) pour les questions ouvertes.
\end{definition}

\courssub{La transposition des personnes et des possessifs}

\begin{propriete}[Règle de transposition des personnes]
Au discours indirect, les personnes sont réinterprétées \textbf{du point de vue de celui qui rapporte} :
\begin{itemize}
\item \esp{yo} (je) devient en général \esp{él / ella} (il/elle) : \esp{Dijo: «Yo estoy cansado»} $\rightarrow$ \esp{Dijo que él estaba cansado.} « Il a dit qu'il était fatigué. »
\item \esp{tú} (tu) devient \esp{él / ella}, ou \esp{yo} si le rapporteur était l'interlocuteur visé : \esp{Me dijo: «Tú tienes razón»} $\rightarrow$ \esp{Me dijo que (yo) tenía razón.} « Il m'a dit que j'avais raison. »
\item Les possessifs suivent la même logique : \esp{mi} $\rightarrow$ \esp{su} ; \esp{tu} $\rightarrow$ \esp{su} ; \esp{nuestro} $\rightarrow$ \esp{su} (ou reste si le rapporteur est inclus).
\end{itemize}
\textbf{Exemple :} \esp{Ana dijo: «Mi hermano vendrá conmigo».} $\rightarrow$ \esp{Ana dijo que su hermano vendría con ella.} « Ana a dit que son frère viendrait avec elle. »
\end{propriete}

\begin{attention}[Le piège du possessif]
Le français « son frère » et l'espagnol \esp{su hermano} sont ambigus (son frère à elle ? à lui ?). L'espagnol peut lever l'ambiguïté avec \esp{el hermano de ella}. En traduction, choisissez toujours la personne \textbf{en fonction du contexte}, jamais mécaniquement.
\end{attention}

\courssub{La transposition des indicateurs temporels et spatiaux (déictiques)}

Les \cle{déictiques} sont des mots dont le sens dépend de la situation d'énonciation (qui parle ? où ? quand ?). Quand on rapporte les paroles plus tard et ailleurs, ces repères « reculent » :

\begin{center}
\begin{tabularx}{\linewidth}{|X|X|X|}
\hline
\rowcolor{popblueL} \textbf{Discours direct} & \textbf{Discours indirect} & \textbf{Français (direct $\rightarrow$ indirect)} \\
\hline
\esp{hoy} & \esp{aquel día} & aujourd'hui $\rightarrow$ ce jour-là \\
\hline
\esp{mañana} & \esp{al día siguiente} & demain $\rightarrow$ le lendemain \\
\hline
\esp{ayer} & \esp{el día anterior} & hier $\rightarrow$ la veille \\
\hline
\esp{ahora} & \esp{entonces} & maintenant $\rightarrow$ alors \\
\hline
\esp{esta noche} & \esp{aquella noche} & ce soir $\rightarrow$ ce soir-là \\
\hline
\esp{la semana próxima} & \esp{la semana siguiente} & la semaine prochaine $\rightarrow$ la semaine suivante \\
\hline
\esp{aquí} & \esp{allí / allá} & ici $\rightarrow$ là / là-bas \\
\hline
\esp{este / esta} & \esp{ese / aquel} & ce(tte) $\rightarrow$ ce(tte)...-là \\
\hline
\esp{venir} & \esp{ir} & venir $\rightarrow$ aller \\
\hline
\esp{traer} & \esp{llevar} & apporter $\rightarrow$ emporter \\
\hline
\end{tabularx}
\end{center}

\begin{exemplebox}[Exemples traduits]
\begin{itemize}
\item \esp{«Vengo aquí todos los días», dijo.} $\rightarrow$ \esp{Dijo que iba allí todos los días.} « Il a dit qu'il allait là tous les jours. »
\item \esp{«Traeré este libro mañana», prometió.} $\rightarrow$ \esp{Prometió que llevaría ese libro al día siguiente.} « Elle a promis qu'elle emporterait ce livre le lendemain. »
\item \esp{«Ayer compré cacao en el mercado de Yaoundé», contó.} $\rightarrow$ \esp{Contó que el día anterior había comprado cacao en el mercado de Yaoundé.} « Il a raconté que la veille il avait acheté du cacao au marché de Yaoundé. »
\end{itemize}
\end{exemplebox}

\begin{popfigure}
\begin{center}
\begin{tikzpicture}[
  direct/.style={rectangle, rounded corners, draw=popblue, fill=popblueL, align=center, minimum height=0.8cm, font=\small},
  indirect/.style={rectangle, rounded corners, draw=poporange, fill=poporangeL, align=center, minimum height=0.8cm, font=\small},
  fleche/.style={-{Stealth[length=2.5mm]}, thick, popdark}
]
\node[direct, font=\small\bfseries] (t1) at (0,0) {Discours direct};
\node[indirect, font=\small\bfseries] (t2) at (7.5,0) {Discours indirect};
\node[direct] (a1) at (0,-1.1) {hoy};
\node[indirect] (b1) at (7.5,-1.1) {aquel día};
\node[direct] (a2) at (0,-2.1) {mañana};
\node[indirect] (b2) at (7.5,-2.1) {al día siguiente};
\node[direct] (a3) at (0,-3.1) {ayer};
\node[indirect] (b3) at (7.5,-3.1) {el día anterior};
\node[direct] (a4) at (0,-4.1) {aquí};
\node[indirect] (b4) at (7.5,-4.1) {allí / allá};
\node[direct] (a5) at (0,-5.1) {este};
\node[indirect] (b5) at (7.5,-5.1) {ese / aquel};
\node[direct] (a6) at (0,-6.1) {yo / mi};
\node[indirect] (b6) at (7.5,-6.1) {él / su};
\foreach \x in {1,2,3,4,5,6} {\draw[fleche] (a\x) -- (b\x);}
\end{tikzpicture}
\end{center}
\legende{Tableau à double entrée : les déictiques et les personnes au passage du style direct au style indirect.}
\end{popfigure}

\courssub{Les trois types de phrases rapportées}

\begin{propriete}[Déclaratives, interrogatives, impératives]
\textbf{1. Phrases déclaratives} : verbe introducteur (\esp{decir}, \esp{afirmar}, \esp{anunciar}, \esp{contestar}) + \esp{que} + indicatif (avec concordance des temps).

\esp{Dijo: «Estoy contento».} $\rightarrow$ \esp{Dijo que estaba contento.} « Il a dit qu'il était content. »

\textbf{2. Phrases interrogatives} :
\begin{itemize}
\item Question \textbf{fermée} (réponse oui/non) : \esp{preguntar si} + indicatif. \emph{Attention :} \esp{si} ne prend \textbf{jamais d'accent} ici.

\esp{Me preguntó: «¿Vienes?»} $\rightarrow$ \esp{Me preguntó si (yo) venía.} « Il m'a demandé si je venais. »
\item Question \textbf{ouverte} : \esp{preguntar} + mot interrogatif \textbf{accentué} (\esp{quién}, \esp{qué}, \esp{dónde}, \esp{cuándo}, \esp{cómo}, \esp{por qué}, \esp{cuánto}).

\esp{Preguntó: «¿Dónde vives?»} $\rightarrow$ \esp{Preguntó dónde vivía.} « Il a demandé où j'habitais. »
\end{itemize}

\textbf{3. Phrases impératives} (ordres, conseils, prières) : deux constructions possibles.
\begin{itemize}
\item \esp{mandar / ordenar / decir / rogar / pedir} + \esp{que} + \textbf{subjonctif imparfait} : \esp{El profesor ordenó que abrieran los cuadernos.} « Le professeur a ordonné qu'ils ouvrent les cahiers. »
\item \esp{mandar / ordenar / pedir / decir a alguien} + \textbf{infinitif} : \esp{El profesor nos mandó abrir los cuadernos.} « Le professeur nous a ordonné d'ouvrir les cahiers. »
\end{itemize}
\end{propriete}

\begin{exemplebox}[Les trois types en contexte camerounais]
\begin{itemize}
\item \textbf{Déclarative :} \esp{«El partido de los Leones Indomables será mañana», anunció el periodista.} $\rightarrow$ \esp{El periodista anunció que el partido de los Leones Indomables sería al día siguiente.} « Le journaliste a annoncé que le match des Lions Indomptables aurait lieu le lendemain. »
\item \textbf{Interrogative fermée :} \esp{«¿Has visto el mercado de Douala?», me preguntó.} $\rightarrow$ \esp{Me preguntó si había visto el mercado de Douala.} « Il m'a demandé si j'avais vu le marché de Douala. »
\item \textbf{Interrogative ouverte :} \esp{«¿Cuándo llega el bendskin?», preguntó ella.} $\rightarrow$ \esp{Ella preguntó cuándo llegaba el bendskin.} « Elle a demandé quand arrivait le bendskin. »
\item \textbf{Impérative :} \esp{«No lleguéis tarde al colegio», dijo la madre.} $\rightarrow$ \esp{La madre dijo que no llegaran tarde al colegio.} « La mère a dit qu'ils n'arrivent pas en retard au collège. »
\end{itemize}
\end{exemplebox}

\begin{attention}[Question indirecte : accent, ordre des mots et « ce que »]
Trois pièges attendent le francophone dans la question indirecte :
\begin{itemize}
\item \textbf{L'accent du mot interrogatif est obligatoire}, même sans point d'interrogation : \esp{Preguntó qué hacía.} « Il a demandé ce qu'il faisait. » Sans accent, *\esp{Preguntó que hacía} se lit comme une déclaration (« il a demandé qu'il faisait... ») : sens faussé !
\item \textbf{Pas d'inversion calquée du français.} En espagnol, le sujet se place naturellement après le verbe ou est sous-entendu : \esp{Preguntó quién era.} « Il a demandé qui c'était. » On peut dire \esp{Preguntó quién era él}, mais jamais une structure figée à la française.
\item \textbf{« Ce que / ce qui » se traduit par \esp{lo que / lo cual}}, jamais par \esp{que} seul : \esp{Preguntó lo que hacía.} « Il a demandé ce qu'il faisait. » Comparez : \esp{Preguntó qué hacía} (avec \esp{qué}, on attend une réponse précise : « qu'est-ce qu'il faisait ? »).
\end{itemize}
\end{attention}

\begin{methode}[Transposer du style direct au style indirect en 4 étapes]
\begin{enumerate}
\item \textbf{Identifier le verbe introducteur et son temps.} S'il est au passé (\esp{dijo}, \esp{preguntó}...), la concordance des temps s'appliquera (présent $\rightarrow$ imparfait, futur $\rightarrow$ conditionnel, parfait $\rightarrow$ plus-que-parfait — voir section suivante). S'il est au présent (\esp{dice}, \esp{pregunta}), les temps ne changent pas : \esp{Dice: «Vendré mañana»} $\rightarrow$ \esp{Dice que vendrá mañana.}
\item \textbf{Adapter les personnes et les possessifs.} \esp{yo} $\rightarrow$ \esp{él/ella} ; \esp{mi} $\rightarrow$ \esp{su} ; conjuguer le verbe rapporté à la nouvelle personne.
\item \textbf{Adapter les temps} selon la concordance : \esp{vendré} $\rightarrow$ \esp{vendría} ; \esp{estoy} $\rightarrow$ \esp{estaba}.
\item \textbf{Adapter les déictiques} : \esp{hoy} $\rightarrow$ \esp{aquel día} ; \esp{mañana} $\rightarrow$ \esp{al día siguiente} ; \esp{aquí} $\rightarrow$ \esp{allí} ; \esp{este} $\rightarrow$ \esp{ese/aquel} ; \esp{venir} $\rightarrow$ \esp{ir}.
\end{enumerate}
Vérification finale : relire la phrase rapportée en se demandant « qui parle, où, quand ? » du point de vue du rapporteur.
\end{methode}

\begin{popfigure}
\begin{center}
\begin{tikzpicture}[
  etape/.style={rectangle, rounded corners, draw=poppurple, fill=poppurpleL, align=center, font=\small, text width=3.4cm, minimum height=1.1cm},
  fleche/.style={-{Stealth[length=2.5mm]}, very thick, poporange}
]
\node[etape, align=center] (p1) at (0,0){\textbf{Parole directe}\\ \esp{«Vendré mañana aquí»}};
\node[etape, right=1.1cm of p1, align=center] (p2){\textbf{Verbe introducteur}\\ \esp{dijo que / si} + temps au passé};
\node[etape, right=1.1cm of p2, align=center] (p3){\textbf{Transpositions}\\ personne + temps\\ + espace};
\node[etape, draw=popgreen, fill=popgreenL, right=1.1cm of p3, align=center] (p4){\textbf{Style indirect}\\ \esp{Dijo que vendría}\\ \esp{al día siguiente allí}};
\draw[fleche] (p1) -- (p2);
\draw[fleche] (p2) -- (p3);
\draw[fleche] (p3) -- (p4);
\end{tikzpicture}
\end{center}
\legende{Schéma de transformation : de la parole directe à la parole rapportée.}
\end{popfigure}

\begin{exoresolu}[Thème : transposer trois phrases]
\textbf{Énoncé.} Traduisez en espagnol au discours indirect (verbe introducteur au passé) : 1) « Je viendrai demain », dit-il. 2) « Où habites-tu ? », demanda-t-elle. 3) « Ouvrez vos cahiers », ordonna le professeur.
\tcblower
\textbf{Solution détaillée.}
\begin{enumerate}
\item Verbe introducteur \esp{dijo} (passé simple) ; \esp{yo} $\rightarrow$ \esp{él} ; futur \esp{vendré} $\rightarrow$ conditionnel \esp{vendría} ; \esp{mañana} $\rightarrow$ \esp{al día siguiente}.

\esp{Dijo que vendría al día siguiente.}
\item Question ouverte : \esp{preguntar} + mot interrogatif accentué \esp{dónde} ; \esp{tú} $\rightarrow$ 3\ieme{} personne ; présent \esp{vives} $\rightarrow$ imparfait \esp{vivía} ; pas d'inversion à la française.

\esp{Ella preguntó dónde vivía.}
\item Impératif rapporté : \esp{ordenar que} + subjonctif imparfait ; \esp{abrid} $\rightarrow$ \esp{abrieran} ; \esp{vuestros} $\rightarrow$ \esp{sus} (ou possessif implicite).

\esp{El profesor ordenó que abrieran los cuadernos.}
\end{enumerate}
\end{exoresolu}

\begin{exercice}[À vous de transposer]
Mettez au discours indirect (commencez par \esp{Dijo que...}, \esp{Preguntó...}, \esp{Mandó que...}) :
\begin{enumerate}
\item \esp{«Hoy tengo un examen aquí», dijo María.}
\item \esp{«¿Has terminado tu trabajo?», preguntó el jefe.}
\item \esp{«¿Cuántos alumnos hay en esta clase?», preguntó el inspector.}
\item \esp{«Traed vuestros libros mañana», dijo la profesora.}
\item \esp{«Ayer visité el mercado central», contó el turista.}
\end{enumerate}
\end{exercice}

\begin{corrige}[Exercice : transposition au discours indirect]
\begin{enumerate}
\item \esp{María dijo que aquel día tenía un examen allí.} (María a dit que ce jour-là elle avait un examen là-bas.)
\item \esp{El jefe preguntó si había terminado su trabajo.} (Le chef a demandé s'il avait terminé son travail.)
\item \esp{El inspector preguntó cuántos alumnos había en aquella clase.} (L'inspecteur a demandé combien d'élèves il y avait dans cette classe-là.)
\item \esp{La profesora dijo que llevaran sus libros al día siguiente.} (La professeure a dit qu'ils apportent leurs livres le lendemain.)
\item \esp{El turista contó que el día anterior había visitado el mercado central.} (Le touriste a raconté que la veille il avait visité le marché central.)
\end{enumerate}
\end{corrige}

\begin{aretenir}[L'essentiel de la section]
\begin{itemize}
\item Le discours indirect rapporte sans citer : subordonnée introduite par \esp{que} (déclaration), \esp{si} (question fermée) ou un mot interrogatif \textbf{accentué} (question ouverte).
\item Trois transpositions obligatoires quand le verbe introducteur est au passé : \textbf{personnes} (\esp{yo} $\rightarrow$ \esp{él}, \esp{mi} $\rightarrow$ \esp{su}), \textbf{temps} (futur $\rightarrow$ conditionnel, présent $\rightarrow$ imparfait), \textbf{déictiques} (\esp{hoy} $\rightarrow$ \esp{aquel día}, \esp{mañana} $\rightarrow$ \esp{al día siguiente}, \esp{aquí} $\rightarrow$ \esp{allí}, \esp{venir} $\rightarrow$ \esp{ir}).
\item Si le verbe introducteur est au présent, les temps ne changent pas : \esp{Dice que vendrá mañana.}
\item Les ordres se rapportent avec \esp{mandar/ordenar que} + subjonctif imparfait, ou avec l'infinitif.
\item Jamais d'inversion calquée du français dans la question indirecte, et « ce que » se traduit \esp{lo que} : \esp{Preguntó lo que hacía.}
\end{itemize}
\end{aretenir}

La section suivante approfondira le point délicat que nous venons d'effleurer : la \textbf{concordance des temps} et le choix précis des \textbf{verbes introducteurs} selon l'intention du locuteur.

\coursec{Concordance des temps et verbes introducteurs}

Dans la section précédente, nous avons appris à transformer les \cle{personnes} (je $\rightarrow$ il/elle), les \cle{possessifs} et les \cle{indicateurs spatio-temporels} (hoy $\rightarrow$ aquel día, aquí $\rightarrow$ allí, mañana $\rightarrow$ al día siguiente). Il reste la transformation la plus délicate, celle que le correcteur du Bac surveille en priorité : le \cle{temps des verbes}. En espagnol comme en français, le choix du temps dans la subordonnée dépend du temps du \cle{verbe introducteur}. C'est la \emph{concordance des temps}.

\courssub{Observation : un même discours, deux situations}

Lisons d'abord ce court dialogue de la vie scolaire à Yaoundé.

\begin{textebox}[En la sala de profesores]
La profesora de español llega al centro a las siete y media. Antes de entrar en clase, habla con el director.

— Mis alumnos de Terminale \textbf{trabajan} mucho —dice la profesora—. \textbf{Han terminado} el programa de traducción y \textbf{harán} un examen blanco la semana próxima.

Por la tarde, el director resume la conversación a un inspector de visita:

— La profesora me \textbf{dijo} que sus alumnos de Terminale \textbf{trabajaban} mucho. Me \textbf{explicó} que \textbf{habían terminado} el programa de traducción y me \textbf{anunció} que \textbf{harían} un examen blanco la semana siguiente. También me \textbf{pidió} que \textbf{preparara} más ejemplares de la prueba.
\end{textebox}

\textbf{Glossaire :} \esp{el centro} = l'établissement scolaire ; \esp{el examen blanco} = l'examen blanc ; \esp{el inspector de visita} = l'inspecteur en visite ; \esp{los ejemplares} = les exemplaires ; \esp{la prueba} = l'épreuve.

\textbf{Questions d'observation :}
\begin{enumerate}
\item Quels temps la professeure utilise-t-elle dans son discours direct ?
\item Quels temps le directeur emploie-t-il pour rapporter les mêmes paroles ?
\item Pourquoi \esp{preparara} est-il au subjonctif imparfait ?
\end{enumerate}

On constate le phénomène essentiel : quand le verbe introducteur passe au \cle{passé} (\esp{dijo}, \esp{explicó}, \esp{anunció}, \esp{pidió}), chaque temps de la subordonnée \emph{recule d'un cran vers le passé}. C'est exactement ce qui se produit en français : « mes élèves \emph{travaillent} » $\rightarrow$ « elle a dit que ses élèves \emph{travaillaient} ».

\courssub{Le principe fondamental}

\begin{propriete}[Règle de concordance des temps]
\begin{itemize}
\item \textbf{Verbe introducteur au présent ou au futur} (\esp{dice que}, \esp{dirá que}) : les temps de la subordonnée \textbf{restent inchangés}. \esp{Dice que está cansado.} « Il dit qu'il est fatigué. »
\item \textbf{Verbe introducteur au passé} (\esp{dijo}, \esp{preguntó}, \esp{ordenó}, \esp{había dicho}, \esp{decía}) : le temps de la subordonnée \textbf{recule d'un cran} :
\begin{itemize}
\item présent $\rightarrow$ \textbf{imparfait} : \esp{Dijo que estaba cansado.} « Il a dit qu'il était fatigué. »
\item passé composé / parfait simple $\rightarrow$ \textbf{plus-que-parfait} : \esp{Dijo que había terminado.} « Il a dit qu'il avait terminé. »
\item futur $\rightarrow$ \textbf{conditionnel} : \esp{Prometió que lo haría.} « Il a promis qu'il le ferait. »
\item subjonctif présent $\rightarrow$ \textbf{subjonctif imparfait} : \esp{Me pidió que fuera.} « Il m'a demandé de venir / que je vienne. »
\item impératif $\rightarrow$ \esp{que} + \textbf{subjonctif imparfait} (ou infinitif) : \esp{Mandó que lo hicieran.} / \esp{Mandó hacerlo.} « Il ordonna qu'on le fasse / de le faire. »
\end{itemize}
\end{itemize}
\end{propriete}

\begin{popfigure}
\begin{center}
\begin{tikzpicture}[font=\small]
\draw[thick, popline] (0,0) -- (12,0);
\foreach \x/\lab in {1/{plus-que-parfait}, 4/{imparfait}, 7/{présent}, 10/{futur}} {
  \fill[popblue] (\x,0) circle (2pt);
  \node[below, popdark] at (\x,-0.1) {\lab};
}
\node[below, popmuted] at (1,-0.7) {\esp{había terminado}};
\node[below, popmuted] at (4,-0.7) {\esp{venía}};
\node[below, popmuted] at (7,-0.7) {\esp{viene}};
\node[below, popmuted] at (10,-0.7) {\esp{vendrá}};
\draw[-{Stealth}, very thick, poporange] (7,0.35) to[bend left=35] node[above, popdark]{présent $\rightarrow$ imparfait} (4,0.35);
\draw[-{Stealth}, very thick, popgreen] (10,0.9) to[bend left=35] node[above, popdark]{futur $\rightarrow$ conditionnel} (5.6,0.9);
\draw[-{Stealth}, very thick, poppurple] (8.6,1.5) to[bend left=30] node[above, popdark]{passé $\rightarrow$ plus-que-parfait} (1,1.5);
\node[popdark, align=center] at (6,2.6) {\textbf{Après un verbe introducteur au passé : recul d'un cran}};
\end{tikzpicture}
\end{center}
\legende{La frise du « recul des temps » : chaque temps glisse vers la gauche quand le verbe introducteur est au passé.}
\end{popfigure}

\courssub{Verbe introducteur au présent ou au futur : rien ne change}

C'est le cas le plus simple : on transpose les personnes et les indicateurs, mais \textbf{jamais les temps}.

\begin{exemplebox}[Temps inchangés après un présent]
\begin{itemize}
\item \esp{« Estoy cansado. » $\rightarrow$ Dice que está cansado.} « Je suis fatigué » $\rightarrow$ Il dit qu'il est fatigué.
\item \esp{« He terminado el deber. » $\rightarrow$ Dice que ha terminado el deber.} « J'ai fini le devoir » $\rightarrow$ Il dit qu'il a fini le devoir.
\item \esp{« Vendré mañana. » $\rightarrow$ Dice que vendrá mañana.} « Je viendrai demain » $\rightarrow$ Il dit qu'il viendra demain.
\item \esp{« Quiero que me ayudes. » $\rightarrow$ Dice que quiere que le ayudes.} « Je veux que tu m'aides » $\rightarrow$ Il dit qu'il veut que tu l'aides.
\end{itemize}
\end{exemplebox}

\courssub{Verbe introducteur au passé : le recul des temps}

\textbf{a) Présent $\rightarrow$ imparfait.}

\esp{« Vengo mañana. » $\rightarrow$ Dijo que venía al día siguiente.} « Je viens demain » $\rightarrow$ Il a dit qu'il venait le lendemain.

\textbf{b) Passé composé / parfait simple $\rightarrow$ plus-que-parfait.}

\esp{« He terminado. » / « Terminé ayer. » $\rightarrow$ Dijo que había terminado el día anterior.} « J'ai fini » / « J'ai fini hier » $\rightarrow$ Il a dit qu'il avait fini la veille.

\textbf{c) Futur $\rightarrow$ conditionnel.}

\esp{« Lo haré. » $\rightarrow$ Prometió que lo haría.} « Je le ferai » $\rightarrow$ Il a promis qu'il le ferait.

\textbf{d) Impératif $\rightarrow$ \esp{que} + subjonctif imparfait (ou infinitif).}

\esp{« ¡Ven! » $\rightarrow$ Me pidió que fuera.} « Viens ! » $\rightarrow$ Il m'a demandé de venir.

\esp{« ¡Hacedlo! » $\rightarrow$ Mandó que lo hicieran.} / \esp{Mandó hacerlo.} « Faites-le ! » $\rightarrow$ Il ordonna qu'on le fasse / de le faire.

\textbf{e) Subjonctif présent $\rightarrow$ subjonctif imparfait.}

\esp{« Quiero que vengas. » $\rightarrow$ Dijo que quería que fuera.} « Je veux que tu viennes » $\rightarrow$ Il a dit qu'il voulait que je vienne.

\begin{popfigure}
\begin{center}
\begin{tikzpicture}[font=\small]
\node[draw, rounded corners, fill=popgreenL, align=center, minimum width=5.2cm, minimum height=3.6cm] (pres) at (0,0) {
\textbf{Verbe introducteur PRÉSENT}\\[2pt]
\esp{dice que está cansado}\\
\esp{dice que ha terminado}\\
\esp{dice que vendrá}\\
\esp{quiere que vengas}\\[2pt]
\emph{temps inchangés}
};
\node[draw, rounded corners, fill=poporangeL, align=center, minimum width=5.2cm, minimum height=3.6cm] (pas) at (8.5,0) {
\textbf{Verbe introducteur PASSÉ}\\[2pt]
\esp{dijo que estaba cansado}\\
\esp{dijo que había terminado}\\
\esp{dijo que vendría}\\
\esp{quería que fueras}\\[2pt]
\emph{recul d'un cran}
};
\draw[-{Stealth}, very thick, popdark] (pres) -- (pas) node[midway, above, align=center]{passage au passé\\ \esp{dijo}, \esp{preguntó}, \esp{ordenó}};
\end{tikzpicture}
\end{center}
\legende{Deux colonnes à mémoriser : à gauche, aucun changement ; à droite, chaque temps recule.}
\end{popfigure}

\begin{center}
\begin{tabularx}{\linewidth}{|X|X|X|}
\hline
\rowcolor{popblueL} Discours direct & Après \esp{dice} (présent) & Après \esp{dijo} (passé) \\
\hline
\esp{estoy} (présent) & \esp{dice que está} & \esp{dijo que estaba} \\
\hline
\esp{he terminado} (p. composé) & \esp{dice que ha terminado} & \esp{dijo que había terminado} \\
\hline
\esp{terminé} (parfait simple) & \esp{dice que terminó} & \esp{dijo que había terminado} \\
\hline
\esp{vendré} (futur) & \esp{dice que vendrá} & \esp{dijo que vendría} \\
\hline
\esp{¡ven!} (impératif) & \esp{me pide que vaya} & \esp{me pidió que fuera} \\
\hline
\esp{quiero que vengas} (subj. prés.) & \esp{dice que quiere que vengas} & \esp{dijo que quería que fueras} \\
\hline
\end{tabularx}
\end{center}

\courssub{Les verbes introducteurs classés par fonction}

Le choix du verbe introducteur dépend de la \cle{nature de la phrase rapportée} : déclarative, interrogative ou injonctive.

\begin{center}
\begin{tabularx}{\linewidth}{|l|X|X|}
\hline
\rowcolor{popblueL} Fonction & Verbes espagnols & Construction \\
\hline
Déclarer & \esp{decir, afirmar, anunciar, explicar, añadir, responder, comentar} & \esp{que} + indicatif \\
\hline
Interroger & \esp{preguntar, querer saber} & \esp{si} (oui/non) ou mot interrogatif (\esp{quién, cuándo, dónde}) \\
\hline
Ordonner / demander & \esp{ordenar, mandar, pedir, rogar, aconsejar, prohibir} & \esp{que} + subjonctif (ou infinitif) \\
\hline
\end{tabularx}
\end{center}

\begin{exemplebox}[Un verbe introducteur pour chaque fonction]
\begin{itemize}
\item Déclarer : \esp{Afirmó que el proyecto era bueno.} « Il affirma que le projet était bon. »
\item Interroger (question totale) : \esp{Me preguntó si estaba listo.} « Il m'a demandé si j'étais prêt. »
\item Interroger (question partielle) : \esp{Preguntó cuándo llegaba el autobús.} « Il demanda quand arrivait le car. »
\item Ordonner : \esp{El profesor ordenó que guardáramos silencio.} « Le professeur ordonna que nous gardions le silence. »
\item Conseiller : \esp{Me aconsejó que estudiara más.} « Il me conseilla d'étudier davantage. »
\item Interdire : \esp{Prohibió que usáramos el teléfono en clase.} « Il interdit que nous utilisions le téléphone en classe. »
\end{itemize}
\end{exemplebox}

\begin{attention}[Question indirecte : pas de point d'interrogation, mais l'accent reste]
Dans la question indirecte, on supprime les signes ¿ ?, mais le mot interrogatif \textbf{garde son accent} : \esp{Preguntó dónde vivía.} « Il demanda où j'habitais. » Et pour une question fermée (oui/non), on utilise \esp{si} : \esp{Preguntó si tenía hambre.} « Il demanda si j'avais faim. » Ne traduisez jamais « si » par \esp{sí} (avec accent) : \esp{sí} signifie « oui » !
\end{attention}

\courssub{Cas particulier : la vérité permanente (complément Bac)}

\begin{propriete}[Exception : la vérité permanente]
Quand la subordonnée exprime une \cle{vérité générale}, un fait scientifique ou une réalité toujours valable, le \textbf{présent peut se maintenir} même après un verbe au passé :
\esp{El maestro dijo que la tierra es redonda.} « Le maître a dit que la terre est ronde. »
\esp{Explicó que el agua hierve a cien grados.} « Il expliqua que l'eau bout à cent degrés. »
L'imparfait reste possible (\esp{dijo que la tierra era redonda}), mais le présent souligne que le fait est encore vrai aujourd'hui. Le français fonctionne pareil : « il a démontré que la terre \emph{tourne} ».
\end{propriete}

\courssub{Comparaison avec le français}

Bonne nouvelle : la concordance des temps est \textbf{presque identique} dans les deux langues. « Il dit qu'il vient » / \esp{dice que viene} ; « il a dit qu'il venait » / \esp{dijo que venía} ; « il a dit qu'il viendrait » / \esp{dijo que vendría}. Deux différences à retenir :

\begin{itemize}
\item Après un verbe de \textbf{volonté au passé}, l'espagnol exige le \textbf{subjonctif imparfait}, là où le français hésite entre subjonctif imparfait (soutenu) et présent : \esp{Me pidió que fuera.} « Il m'a demandé que je vienne / de venir. »
\item Le français remplace souvent \esp{que} + subjonctif par un \textbf{infinitif} : \esp{Mandó que lo hicieran.} = « Il ordonna de le faire. » Au thème, les deux constructions espagnoles sont acceptées (\esp{mandó hacerlo} est même plus élégant quand le sujet est le même).
\end{itemize}

\begin{methode}[Thème : la concordance en quatre étapes]
\begin{enumerate}
\item \textbf{Repérer le verbe introducteur français} et son temps : « il dit » (présent) ou « il a dit » (passé) ? C'est lui qui déclenche ou non le recul.
\item \textbf{Traduire le verbe introducteur} : \esp{dice} / \esp{dijo} ; « demander » = \esp{pedir} (jamais \esp{demandar}, qui signifie « exiger » ou « poursuivre en justice »).
\item \textbf{Appliquer la concordance} : présent inchangé, passé $\rightarrow$ recul d'un cran (présent $\rightarrow$ imparfait, futur $\rightarrow$ conditionnel, passé $\rightarrow$ plus-que-parfait).
\item \textbf{Vérifier la nature de la phrase} : déclarative (\esp{que} + indicatif), interrogative (\esp{si} ou mot interrogatif accentué), injonctive (\esp{que} + subjonctif ou infinitif).
\end{enumerate}
\end{methode}

\begin{exoresolu}[Thème guidé : la concordance en action]
Traduisez en espagnol : « Le directeur a annoncé que les cours commenceraient à huit heures et il a demandé aux élèves d'arriver à l'heure. Il a ajouté qu'il avait prévenu les professeurs la veille. »
\tcblower
\textbf{Étape 1 :} verbes introducteurs au passé : « a annoncé » $\rightarrow$ \esp{anunció} ; « a demandé » $\rightarrow$ \esp{pidió} ; « a ajouté » $\rightarrow$ \esp{añadió}. Le recul des temps s'impose partout.

\textbf{Étape 2 :} « commenceraient » : futur dans le discours direct (« les cours commenceront ») $\rightarrow$ conditionnel espagnol : \esp{empezarían}. « d'arriver » : injonction $\rightarrow$ \esp{que} + subjonctif imparfait : \esp{que llegaran} (ou infinitif \esp{les pidió llegar a tiempo}). « il avait prévenu » : déjà un plus-que-parfait français $\rightarrow$ \esp{había avisado}. « la veille » $\rightarrow$ \esp{el día anterior}.

\textbf{Traduction complète :} \esp{El director anunció que las clases empezarían a las ocho y pidió a los alumnos que llegaran a tiempo. Añadió que había avisado a los profesores el día anterior.}
\end{exoresolu}

\begin{attention}[Le piège n° 1 : conserver le futur après « il a dit que... »]
L'erreur la plus fréquente des francophones : écrire \esp{dijo que vendrá} pour « il a dit qu'il viendrait ». Si l'action est rapportée au passé, le futur est \textbf{impossible} : il faut le \textbf{conditionnel} \esp{dijo que vendría}. Même vigilance pour le présent : « il a dit qu'il était fatigué » = \esp{dijo que estaba cansado}, jamais \esp{dijo que está cansado} (sauf vérité permanente). Enfin, n'oubliez pas que \esp{pedir} se construit sans préposition : \esp{me pidió que fuera}, et non « me pidió de que fuera ».
\end{attention}

\begin{savaistu}[Le discours indirect, star de la presse hispanophone]
Ouvrez n'importe quel journal espagnol ou latino-américain : les articles de politique sont remplis de discours rapporté (\esp{el ministro anunció que..., fuentes oficiales afirmaron que...}). Les journalistes utilisent aussi le \cle{discours indirect libre}, qui mélange la voix du narrateur et celle du personnage sans verbe introducteur apparent. Au Bac, la version propose régulièrement deux ou trois phrases de discours rapporté : entraînez-vous à repérer le temps du verbe introducteur avant même de traduire la première subordonnée — c'est lui qui commande toute la phrase.
\end{savaistu}

\begin{exercice}[À vous de jouer]
Transposez au discours indirect en commençant par \esp{Dijo que...} ou un verbe adapté :
\begin{enumerate}
\item \esp{« Estoy cansado y he trabajado mucho. »}
\item \esp{« Iré a Douala mañana. »}
\item \esp{« ¡Cierra la puerta! », me dijo mi madre.}
\item \esp{« ¿Dónde vives? », me preguntó.}
\item Traduisez : « Elle a promis qu'elle m'écrirait et m'a demandé de l'attendre. »
\end{enumerate}
\end{exercice}

\begin{corrige}[Exercice : discours indirect et concordance]
\begin{enumerate}
\item \esp{Dijo que estaba cansado y que había trabajado mucho.} (présent $\rightarrow$ imparfait ; passé composé $\rightarrow$ plus-que-parfait)
\item \esp{Dijo que iría a Douala al día siguiente.} (futur $\rightarrow$ conditionnel ; \esp{mañana} $\rightarrow$ \esp{al día siguiente})
\item \esp{Mi madre me dijo que cerrara la puerta.} (impératif $\rightarrow$ \esp{que} + subjonctif imparfait)
\item \esp{Me preguntó dónde vivía.} (question partielle : accent conservé, présent $\rightarrow$ imparfait)
\item \esp{Prometió que me escribiría y me pidió que la esperara.} (futur $\rightarrow$ conditionnel ; injonction $\rightarrow$ subjonctif imparfait)
\end{enumerate}
\end{corrige}

\begin{aretenir}[L'essentiel de la section]
\begin{itemize}
\item Verbe introducteur au \textbf{présent/futur} : temps inchangés.
\item Verbe introducteur au \textbf{passé} : recul d'un cran — présent $\rightarrow$ imparfait ; passé $\rightarrow$ plus-que-parfait ; futur $\rightarrow$ conditionnel ; subjonctif présent $\rightarrow$ subjonctif imparfait ; impératif $\rightarrow$ \esp{que} + subjonctif imparfait (ou infinitif).
\item Verbes introducteurs : déclarer (\esp{decir, afirmar, anunciar} + \esp{que} + indicatif), interroger (\esp{preguntar} + \esp{si} ou mot accentué), ordonner (\esp{ordenar, mandar, pedir, prohibir} + \esp{que} + subjonctif).
\item Exception : la vérité permanente garde le présent (\esp{dijo que la tierra es redonda}).
\item Au thème : repérer d'abord le temps du verbe introducteur français — c'est lui qui décide de tout.
\end{itemize}
\end{aretenir}

\coursec{Entraînement intégré : thème et version}

Nous savons désormais conjuguer les verbes défectifs, traduire la restriction et appliquer la concordance des temps au discours indirect. Il reste à réunir ces trois compétences dans la situation réelle de l'examen : un \cle{thème} (français $\rightarrow$ espagnol) et une \cle{version} (espagnol $\rightarrow$ français). C'est ici que tout se joue : au Bac, le correcteur ne note pas ce que vous savez, mais ce que vous écrivez. Entraînons-nous donc méthodiquement.

\courssub{Stratégie générale : observer avant de traduire}

\begin{methode}[Stratégie de version : repérer avant de traduire]
\begin{enumerate}
\item \textbf{Lire} le texte espagnol entier une première fois, sans traduire, pour saisir le sens global.
\item \textbf{Souligner} d'une couleur tous les verbes défectifs (\esp{gustar, doler, parecer, soler, llover}) et vérifier leur construction : qui est le sujet grammatical ?
\item \textbf{Entourer} d'une autre couleur toutes les restrictions (\esp{sólo, únicamente, no... más que, no... sino, apenas}) et identifier ce qui est restreint.
\item \textbf{Encadrer} les verbes introducteurs du discours indirect (\esp{dijo que, anunció que, preguntó si}) et vérifier le temps de la subordonnée.
\item \textbf{Traduire} phrase par phrase, puis \textbf{vérifier} chaque construction repérée dans la traduction française : le sujet du verbe défectif est-il bien devenu un complément en français ? La restriction est-elle rendue par « seulement / ne... que » ? Les temps français suivent-ils la concordance ?
\item \textbf{Relire} la traduction à voix haute : une version doit être du français correct et naturel, pas un calque mot à mot.
\end{enumerate}
\end{methode}

\begin{attention}[Les trois erreurs les plus sanctionnées au Bac]
Avant de rendre votre copie, relisez-vous en ciblant ces trois points :
\begin{enumerate}
\item \esp{me gusto} au lieu de \esp{me gusta} : le verbe s'accorde avec la chose qui plaît, jamais avec la personne. \esp{Me gusta el cacao.} « Le cacao me plaît. »
\item \esp{no... que} sans \esp{más} : la restriction exige \esp{no... más que}. \esp{No tengo más que dos hermanos.} « Je n'ai que deux frères. »
\item Le futur conservé après un verbe au passé : après \esp{dijo que}, le futur devient \textbf{conditionnel}. \esp{Dijo que vendría.} « Il a dit qu'il viendrait. »
\end{enumerate}
Ces trois fautes coûtent à elles seules plusieurs points : une relecture ciblée de deux minutes peut les éliminer toutes.
\end{attention}

\courssub{Thème guidé n°1 : phrases isolées (défectifs et restriction)}

\begin{exemplebox}[Phrase modèle de thème]
\esp{El delegado anunció que sólo los alumnos de Terminale participarían en el debate al día siguiente.}\\
« Le délégué a annoncé que seuls les élèves de Terminale participeraient au débat le lendemain. »\\[4pt]
Analyse : verbe introducteur au passé simple (\esp{anunció}) $\rightarrow$ conditionnel dans la subordonnée (\esp{participarían}) ; restriction \esp{sólo} placée devant l'élément restreint ; déictique \esp{mañana} $\rightarrow$ \esp{al día siguiente}.
\end{exemplebox}

\begin{exoresolu}[Thème guidé : trois phrases]
Énoncé. Traduisez en espagnol : 1. « Seuls les délégués ont assisté à la réunion. » 2. « Cette décision ne concerne que les élèves de Première. » 3. « Il ne reste que deux jours avant les vacances. »
\tcblower
Solution détaillée.\\
1. \esp{Sólo los delegados asistieron a la reunión.} — Restriction sur le sujet : \esp{sólo} se place devant \esp{los delegados}. Passé simple car action ponctuelle achevée.\\
2. \esp{Esta decisión no atañe más que a los alumnos de Primera.} — Verbe défectif \esp{atañer} : le sujet grammatical est \esp{esta decisión} (singulier), la personne concernée est introduite par \esp{a}. Restriction \esp{no... más que} autour du verbe.\\
3. \esp{No quedan más que dos días antes de las vacaciones.} — \esp{quedar} se construit ici comme un verbe défectif : le sujet est \esp{dos días} (pluriel), d'où \esp{quedan}. Piège : ne pas écrire \esp{no queda que...}.
\end{exoresolu}

\begin{exercice}[À vous : thème de phrases]
Traduisez en espagnol :
\begin{enumerate}
\item « Mes jambes me font mal après le match de football. »
\item « Nous n'avons reçu que trois candidatures pour le poste de délégué. »
\item « Il pleut souvent à Douala en juillet. »
\item « Ce problème ne regarde que les membres du bureau. » (employer \esp{concernir})
\item « Elle a l'habitude d'arriver la première au lycée. » (employer \esp{soler})
\end{enumerate}
\end{exercice}

\begin{corrige}[Exercice : thème de phrases]
1. \esp{Me duelen las piernas después del partido de fútbol.} — Sujet pluriel \esp{las piernas} $\rightarrow$ \esp{duelen} ; pronom \esp{me} pour la personne.\\
2. \esp{No hemos recibido más que tres candidaturas para el puesto de delegado.} — \esp{no... más que} ; passé composé.\\
3. \esp{Llueve a menudo en Duala en julio.} — \esp{llover} à la troisième personne du singulier, sans sujet exprimé ; \esp{Duala} est la forme espagnole de « Douala ».\\
4. \esp{Este problema no concierne más que a los miembros de la junta directiva.} — \esp{concernir} se conjugue comme \esp{atañer} (troisième personne) ; « le bureau » d'une association se dit \esp{la junta directiva} (attention au faux ami : \esp{buró} désigne une table de nuit !).\\
5. \esp{Suele llegar la primera al liceo.} — \esp{soler} + infinitif ; accord de \esp{la primera} avec le sujet féminin sous-entendu.
\end{corrige}

\courssub{Thème guidé n°2 : transposition d'un dialogue au discours indirect}

\begin{textebox}[Diálogo en la sala de clase]
— \textbf{El delegado}: «Mañana habrá una reunión aquí, en esta sala. Sólo los alumnos de Terminale pueden participar.»\\
— \textbf{Amina}: «¿Tenemos que traer nuestros cuadernos?»\\
— \textbf{El delegado}: «No traigáis nada más que un bolígrafo. La reunión durará dos horas.»\\
— \textbf{Amina}: «Me gusta esta idea, pero me duele la cabeza hoy. ¿Puedo venir pasado mañana?»\\
— \textbf{El delegado}: «No, la reunión es mañana. Descansa esta tarde.»
\end{textebox}

\textbf{Glossaire :} \esp{sala} : salle ; \esp{cuaderno} : cahier ; \esp{bolígrafo} : stylo ; \esp{durar} : durer ; \esp{descansar} : se reposer ; \esp{pasado mañana} : après-demain.

\begin{exoresolu}[Du direct à l'indirect]
Énoncé. Transposez le dialogue ci-dessus au discours indirect, en commençant par : \esp{El delegado anunció que...}
\tcblower
Solution détaillée.\\
\esp{El delegado anunció que al día siguiente habría una reunión allí, en esa sala, y que sólo los alumnos de Terminale podían participar. Amina preguntó si tenían que llevar sus cuadernos. El delegado respondió que no llevaran nada más que un bolígrafo y que la reunión duraría dos horas. Amina dijo que le gustaba esa idea, pero que le dolía la cabeza ese día, y preguntó si podía ir dos días después. El delegado contestó que no, que la reunión era al día siguiente, y le aconsejó que descansara esa tarde.}\\[4pt]
Justifications : \esp{mañana} $\rightarrow$ \esp{al día siguiente} ; \esp{pasado mañana} $\rightarrow$ \esp{dos días después} ; \esp{aquí} $\rightarrow$ \esp{allí} ; \esp{esta} $\rightarrow$ \esp{esa} ; futur \esp{habrá} $\rightarrow$ conditionnel \esp{habría} ; présent \esp{pueden} $\rightarrow$ imparfait \esp{podían} ; question totale introduite par \esp{si} ; impératif \esp{no traigáis} $\rightarrow$ \esp{que} + subjonctif imparfait \esp{no llevaran} ; impératif de conseil rendu par \esp{le aconsejó que descansara}.
\end{exoresolu}

\begin{propriete}[Rappel des transformations direct $\rightarrow$ indirect]
\begin{center}
\begin{tabularx}{\linewidth}{|X|X|X|}
\hline
\rowcolor{popblueL} Style direct & Style indirect & Remarque \\
\hline
\esp{mañana} & \esp{al día siguiente} & déictique temporel \\
\hline
\esp{pasado mañana} & \esp{dos días después} & déictique temporel \\
\hline
\esp{ayer} & \esp{el día anterior} & déictique temporel \\
\hline
\esp{hoy} & \esp{ese día / aquel día} & déictique temporel \\
\hline
\esp{aquí} & \esp{allí} & déictique spatial \\
\hline
\esp{este / esta} & \esp{ese / esa} & déictique spatial \\
\hline
futur \esp{hará} & conditionnel \esp{haría} & concordance des temps \\
\hline
présent \esp{hace} & imparfait \esp{hacía} & concordance des temps \\
\hline
impératif \esp{haz} & \esp{que} + subjonctif imparfait \esp{hiciera} & ordre rapporté \\
\hline
\end{tabularx}
\end{center}
\end{propriete}

\courssub{Version : un texte à traduire}

\begin{methode}[Stratégie de version appliquée]
Appliquez au texte suivant la méthode en six étapes : lisez d'abord, soulignez les défectifs (\esp{gustar, doler, soler}), entourez les restrictions (\esp{sólo, no... más que, apenas}), encadrez les verbes introducteurs (\esp{dijo, explicó, preguntó}), puis traduisez et vérifiez.
\end{methode}

\begin{textebox}[La reunión del consejo de delegados]
Ayer por la tarde, el consejo de delegados se reunió en el aula de español. Sólo asistieron doce alumnos, porque a muchos no les gustaba la hora elegida. El delegado principal abrió la sesión y dijo que aquel año el colegio organizaría una semana cultural. Explicó que no había más que tres días para presentar los proyectos y que apenas quedaba dinero en la caja del club. Amina, la delegada de Terminale, declaró que a sus compañeros les gustaban mucho las actividades teatrales, pero que les dolía no poder participar por falta de tiempo. Añadió que solía ensayar con ellos los sábados por la mañana. El profesor de español, presente en la reunión, afirmó que apoyaría todos los proyectos y preguntó si los alumnos querían invitar a una escuela de Guinea Ecuatorial. Todos respondieron que sí, y la reunión terminó con aplausos. Al salir, Amina sonrió: no había más que una cosa segura, aquella semana cultural sería inolvidable.
\end{textebox}

\textbf{Glossaire :} \esp{consejo} : conseil ; \esp{aula} : salle de classe ; \esp{elegida} : choisie ; \esp{sesión} : séance ; \esp{caja} : caisse ; \esp{ensayar} : répéter ; \esp{apoyar} : soutenir ; \esp{aplausos} : applaudissements ; \esp{inolvidable} : inoubliable.

\begin{exoresolu}[Version corrigée]
Énoncé. Traduisez le texte « La reunión del consejo de delegados » en français.
\tcblower
Solution détaillée.\\
« Hier après-midi, le conseil des délégués s'est réuni dans la salle d'espagnol. Seuls douze élèves y ont assisté, car l'heure choisie ne plaisait pas à beaucoup d'entre eux. Le délégué principal a ouvert la séance et a dit que cette année-là le collège organiserait une semaine culturelle. Il a expliqué qu'il ne restait que trois jours pour présenter les projets et qu'il ne restait presque plus d'argent dans la caisse du club. Amina, la déléguée de Terminale, a déclaré que les activités théâtrales plaisaient beaucoup à ses camarades, mais que cela leur faisait mal de ne pas pouvoir participer faute de temps. Elle a ajouté qu'elle avait l'habitude de répéter avec eux le samedi matin. Le professeur d'espagnol, présent à la réunion, a affirmé qu'il soutiendrait tous les projets et a demandé si les élèves voulaient inviter une école de Guinée équatoriale. Tous ont répondu que oui, et la réunion s'est terminée par des applaudissements. En sortant, Amina a souri : une seule chose était sûre, cette semaine culturelle serait inoubliable. »\\[4pt]
Points de grammaire vérifiés : \esp{no les gustaba} $\rightarrow$ « ne plaisait pas » (imparfait conservé, verbe au singulier car le sujet est \esp{la hora}) ; \esp{les gustaban} $\rightarrow$ « plaisaient » (sujet pluriel \esp{las actividades}) ; \esp{les dolía} $\rightarrow$ « cela leur faisait mal » ; \esp{no había más que} $\rightarrow$ « il ne restait que » ; \esp{apenas quedaba} $\rightarrow$ « il ne restait presque plus » ; futurs rapportés \esp{organizaría, apoyaría, sería} $\rightarrow$ conditionnels français.
\end{exoresolu}

\courssub{Thème : un texte de vie citoyenne à traduire}

\begin{exercice}[Thème : le discours du délégué]
Traduisez en espagnol le texte suivant.\\[4pt]
« Hier, le délégué de classe a réuni tous les élèves dans la cour du lycée. Il a annoncé que seuls les élèves de Terminale participeraient au débat sur la vie citoyenne le lendemain. Il a expliqué que cette réunion ne concernait que les candidats au baccalauréat et qu'il ne restait que deux jours pour s'inscrire. Puis il a ajouté que le sujet du débat plaisait beaucoup aux professeurs. Une élève a demandé si elle pouvait inviter son frère. Le délégué a répondu que non, car il n'y avait que cinquante chaises dans la salle. Enfin, il a conseillé à tous de préparer leurs arguments ce soir-là. »
\end{exercice}

\begin{corrige}[Exercice : thème « le discours du délégué »]
\esp{Ayer, el delegado de clase reunió a todos los alumnos en el patio del liceo. Anunció que sólo los alumnos de Terminale participarían en el debate sobre la vida ciudadana al día siguiente. Explicó que esa reunión no concernía más que a los candidatos al bachillerato y que no quedaban más que dos días para inscribirse. Luego añadió que el tema del debate les gustaba mucho a los profesores. Una alumna preguntó si podía invitar a su hermano. El delegado respondió que no, porque no había más que cincuenta sillas en la sala. Por fin, aconsejó a todos que prepararan sus argumentos esa noche.}\\[4pt]
Justifications : futur « participeraient » $\rightarrow$ conditionnel \esp{participarían} après \esp{anunció} ; « ne concernait que » $\rightarrow$ \esp{no concernía más que a} ; « plaisait aux professeurs » $\rightarrow$ \esp{les gustaba mucho a los profesores} (verbe au singulier car le sujet est \esp{el tema}) ; question indirecte avec \esp{si} + imparfait \esp{podía} ; conseil rapporté $\rightarrow$ \esp{aconsejó que prepararan} (subjonctif imparfait) ; « le lendemain » $\rightarrow$ \esp{al día siguiente} ; « ce soir-là » $\rightarrow$ \esp{esa noche}.
\end{corrige}

\courssub{Grille d'auto-évaluation et correction collective}

Avant de rendre votre thème ou votre version, interrogez votre production à l'aide de cette grille. Chaque « oui » vaut un point ; visez 10 sur 10.

\begin{center}
\begin{tabularx}{\linewidth}{|X|X|}
\hline
\rowcolor{popblueL} Critère de vérification & Question à se poser \\
\hline
Construction du verbe défectif & Le verbe \esp{gustar/doler} est-il accordé avec la chose, et le pronom (\esp{me, te, le}) avec la personne ? \\
\hline
Personne concernée & Ai-je bien mis la préposition \esp{a} devant la personne (\esp{a los alumnos}) ? \\
\hline
Restriction \esp{más que} & Ai-je écrit \esp{no... más que} complet, sans oublier \esp{más} ? \\
\hline
Choix \esp{más que / sino} & Ai-je utilisé \esp{sino} uniquement pour une opposition après négation ? \\
\hline
\esp{sólo / únicamente} & La restriction est-elle placée juste devant l'élément restreint ? \\
\hline
Verbe introducteur & Le verbe introducteur est-il au passé ? Si oui, la subordonnée est-elle au conditionnel ou à l'imparfait ? \\
\hline
Concordance des temps & Aucun futur ne subsiste après un verbe au passé ? \\
\hline
Déictiques temporels & \esp{mañana} $\rightarrow$ \esp{al día siguiente} ; \esp{hoy} $\rightarrow$ \esp{ese día} ? \\
\hline
Déictiques spatiaux & \esp{aquí} $\rightarrow$ \esp{allí} ; \esp{este} $\rightarrow$ \esp{ese} ? \\
\hline
Relecture finale & La phrase traduite est-elle naturelle et correcte dans la langue d'arrivée ? \\
\hline
\end{tabularx}
\end{center}

\begin{experience}[Correction collective commentée]
Le professeur projette la copie anonyme d'un élève contenant trois erreurs types : \esp{me gusto} au lieu de \esp{me gusta}, \esp{no tengo que más que} (forme fautive mélangeant « devoir » et restriction) au lieu de \esp{no tengo más que}, et un futur \esp{vendrá} conservé après \esp{dijo que} au lieu du conditionnel \esp{vendría}. Par groupes de quatre, les élèves : 1. repèrent les trois erreurs à l'aide de la grille ; 2. proposent une correction justifiée en citant la règle ; 3. rédigent au tableau la version corrigée de la phrase. Chaque groupe doit justifier oralement chaque choix : accord du verbe défectif, forme complète de la restriction, concordance des temps. La classe vote ensuite pour la justification la plus claire.
\end{experience}

\begin{savaistu}[Ces structures au Bac camerounais]
Au baccalauréat camerounais, le sujet de thème compte presque toujours une phrase contenant une restriction (« ne... que », « seuls ») et une phrase au discours indirect introduite par un verbe au passé (« il a dit que », « elle a demandé si »). Ce n'est pas un hasard : ces deux structures permettent au correcteur de vérifier en une seule phrase la concordance des temps, les déictiques et la syntaxe de la subordonnée. Maîtriser ces deux mécanismes, c'est donc sécuriser plusieurs points à chaque épreuve de traduction.
\end{savaistu}

\courssub{Carte mentale récapitulative de la leçon}

\begin{popfigure}
\begin{tikzpicture}[mindmap, grow cyclic,
  every node/.style={concept, text=white, align=center},
  root concept/.append style={concept color=popdark, font=\bfseries},
  level 1/.append style={level distance=3.4cm, sibling angle=120},
  level 2/.append style={level distance=2.4cm, sibling angle=50},
  concept connection/.append style={color=popmuted}]
\node[root concept, align=center]{Traduction\\Tle A}
  child[concept color=popblue]{ node[align=center]{Verbes\\défectifs}
    child { node[concept color=popblue!70, font=\small, align=center]{accord avec\\la chose} }
    child { node[concept color=popblue!70, font=\small, align=center]{personne :\\pronom + a} }
  }
  child[concept color=poporange]{ node {Restriction}
    child { node[concept color=poporange!70, font=\small, align=center]{no... más que\\(quantité)} }
    child { node[concept color=poporange!70, font=\small, align=center]{no... sino\\(opposition)} }
  }
  child[concept color=popgreen]{ node[align=center]{Discours\\indirect}
    child { node[concept color=popgreen!70, font=\small, align=center]{futur $\rightarrow$\\conditionnel} }
    child { node[concept color=popgreen!70, font=\small, align=center]{mañana $\rightarrow$\\al día siguiente} }
  };
\end{tikzpicture}
\legende{Carte mentale récapitulative : les trois branches de la leçon et leur règle-clé.}
\end{popfigure}

\begin{aretenir}[L'essentiel de l'entraînement intégré]
\begin{itemize}
\item En version, \textbf{repérez avant de traduire} : soulignez les verbes défectifs, entourez les restrictions, encadrez les verbes introducteurs.
\item En thème, vérifiez systématiquement : accord du verbe défectif avec la \textbf{chose}, forme complète \esp{no... más que}, conditionnel après un verbe introducteur au passé, transformation des déictiques.
\item La grille d'auto-évaluation à dix critères est votre meilleure alliée : deux minutes de relecture ciblée valent plusieurs points.
\end{itemize}
\end{aretenir}

\newpage

\coursec{Activité d'intégration}

\coursec{Leçon E14 : Traduction — Verbes défectifs, phrase restrictive, discours indirect}

\courssub{Objectifs de l'activité}

À l'issue de cette activité d'intégration, l'élève sera capable de :

\begin{itemize}
\item rédiger en espagnol un court article de presse lycéenne (12 à 15 lignes) rapportant un événement réel de la vie du lycée ;
\item mobiliser correctement au moins \cle{trois verbes défectifs} (\esp{gustar}, \esp{doler}, \esp{parecer}, \esp{soler}, \esp{concernir}...) avec la construction à l'indirect (\esp{a los alumnos les gusta...}) ;
\item exprimer la \cle{restriction} par deux procédés différents (\esp{no... más que} et \esp{no... sino}) ;
\item transposer au \cle{discours indirect} au moins quatre paroles rapportées, avec des verbes introducteurs variés et une concordance des temps correcte ;
\item relire et évaluer la production d'un camarade à l'aide d'une grille d'auto-évaluation.
\end{itemize}

\courssub{Situation}

Tu es membre du \esp{Club Hispanófilo} de ton lycée à Yaoundé et tu collabores au journal du lycée, \emph{La Voz del Patio}. Hier après-midi, le club s'est réuni en salle 12 pour préparer la \esp{Semana de la Lengua Española}. Le président du club, Etienne, a fait plusieurs annonces importantes ; les élèves ont réagi, exprimé leurs goûts et leurs gênes ; et le débat a achoppé sur une question délicate : le budget, très limité.

Ton rédacteur en chef te confie la rédaction d'une \cle{crónica} (article-reportage) de 12 à 15 lignes pour le prochain numéro du journal.

Voici le compte rendu que ta camarade Aïcha a pris pendant la réunion, au \cle{discours direct} :

\begin{textebox}[Apuntes de Aïcha — reunión del Club Hispanófilo, ayer a las 16 h]
El presidente, Etienne, abre la sesión: «La Semana de la Lengua tendrá lugar del 12 al 16 de mayo. Mañana recibiremos a una delegación de estudiantes de Guinea Ecuatorial.»

Añade: «Este año no organizaremos una feria del libro, sino un concurso de poesía y una noche de teatro.»

Precisa: «No tenemos más que 25 000 francos para toda la semana.»

Reacciones de los alumnos: a muchos les gusta la idea del concurso de poesía; a algunos les duele que no haya feria del libro; a otros no les parece suficiente el presupuesto. Los miembros del club suelen reunirse los jueves, pero esta semana se reunirán también el lunes.
\end{textebox}

\textbf{Glossaire :} \esp{apuntes} : notes ; \esp{feria del libro} : foire du livre ; \esp{presupuesto} : budget ; \esp{suelen reunirse} : ont l'habitude de se réunir ; \esp{Guinea Ecuatorial} : Guinée équatoriale.

\courssub{Consignes}

\begin{enumerate}
\item \textbf{Compréhension.} Lis attentivement les notes d'Aïcha. Souligne : (a) toutes les paroles prononcées par Etienne au discours direct ; (b) tous les verbes défectifs présents dans les réactions des élèves ; (c) les deux expressions de restriction.
\item \textbf{Grammaire — discours indirect.} Réécris sur ton cahier les quatre annonces d'Etienne au discours indirect, en commençant par : \esp{El presidente dijo que...} / \esp{añadió que...} / \esp{precisó que...} / \esp{anunció que...} Attention à la concordance des temps (\esp{tendrá} $\rightarrow$ \esp{tendría} ; \esp{recibiremos} $\rightarrow$ \esp{recibirían}...) et aux indicateurs spatio-temporels (\esp{mañana} $\rightarrow$ \esp{al día siguiente} ; \esp{este año} $\rightarrow$ \esp{ese año}).
\item \textbf{Grammaire — verbes défectifs.} Rédige trois phrases sur les réactions des élèves en utilisant \esp{gustar}, \esp{doler} et \esp{parecer} avec le pronom d'objet indirect correct (\esp{les gusta}, \esp{les duele}, \esp{no les parece}).
\item \textbf{Grammaire — restriction.} Intègre dans ton article : (a) une restriction avec \esp{no... más que} pour le budget ; (b) une restriction avec \esp{no... sino} pour le changement de programme.
\item \textbf{Production.} Rédige ta \cle{crónica} de 12 à 15 lignes : un titre, une phrase d'introduction (qui, quoi, où, quand), le corps de l'article avec les annonces au discours indirect, les réactions et la conclusion. Respecte le contrat : au moins 3 verbes défectifs, 2 restrictions, 4 phrases au discours indirect avec des verbes introducteurs variés.
\item \textbf{Relecture croisée.} Échange ton article avec ton binôme. Évalue sa production avec la grille ci-dessous, puis corrige le tien avant l'affichage au mur de la salle d'espagnol.
\end{enumerate}

\begin{methode}[Grille d'auto-évaluation (relecture croisée)]
\begin{enumerate}
\item L'article a un titre et une phrase d'introduction complète (qui, quoi, où, quand). \hfill /2
\item Au moins 3 verbes défectifs, correctement construits avec \esp{a} + nom et pronom d'objet indirect. \hfill /3
\item Deux restrictions : une avec \esp{no... más que}, une avec \esp{no... sino}. \hfill /2
\item Quatre paroles rapportées au discours indirect, avec 4 verbes introducteurs différents. \hfill /4
\item Concordance des temps et indicateurs spatio-temporels correctement transposés. \hfill /4
\item Orthographe, accents, ponctuation espagnole (¿ ¡), longueur respectée. \hfill /5
\end{enumerate}
\end{methode}

\courssub{Ressources à mobiliser}

\begin{aretenir}[Boîte à outils de la crónica]
\textbf{1. Verbes défectifs} (construction : \esp{a} + personne + pronom COI + verbe accordé avec la chose) :
\esp{a los alumnos les gusta la idea} ; \esp{les duele la ausencia de la feria} ; \esp{no les parece suficiente} ; \esp{el club suele reunirse los jueves} ; \esp{ese asunto no les concierne}.

\textbf{2. La restriction :}
\esp{no tenemos más que 25 000 francos} (= nous n'avons que) ;
\esp{no organizaremos una feria, sino un concurso} (= non pas... mais) ;
\esp{sólo / únicamente} (= seulement) ;
\esp{apenas} (= à peine).

\textbf{3. Discours indirect :} verbes introducteurs variés :
\esp{decir que}, \esp{anunciar que}, \esp{añadir que}, \esp{precisar que}, \esp{explicar que}, \esp{recordar que}.
Concordance des temps (style indirect au passé) :
\begin{center}
\begin{tabularx}{\linewidth}{|X|X|X|}\hline
\rowcolor{popblueL} Direct (présent/futur) & Indirect (imparfait/conditionnel) & Exemple \\ \hline
Présent indicatif & Imparfait indicatif & \esp{tenemos} $\rightarrow$ \esp{tenían} \\ \hline
Futur simple & Conditionnel présent & \esp{tendrá} $\rightarrow$ \esp{tendría} \\ \hline
Passé composé & Plus-que-parfait & \esp{ha dicho} $\rightarrow$ \esp{había dicho} \\ \hline
Impératif / Subjonctif présent & Subjonctif imparfait & \esp{venid} / \esp{que vengan} $\rightarrow$ \esp{que vinieran} \\ \hline
\end{tabularx}
\end{center}
Indicateurs spatio-temporels :
\begin{center}
\begin{tabular}{|l|l|}\hline
\rowcolor{popblueL} Direct & Indirect \\ \hline
\esp{hoy} & \esp{aquel día} \\ \hline
\esp{mañana} & \esp{al día siguiente} \\ \hline
\esp{ayer} & \esp{el día anterior} \\ \hline
\esp{aquí} & \esp{allí} \\ \hline
\esp{este / esta} & \esp{ese / aquella} \\ \hline
\esp{nosotros} & \esp{ellos / el club} \\ \hline
\end{tabular}
\end{center}
\end{aretenir}

\begin{popfigure}
\begin{tikzpicture}[mindmap, grow cyclic, every node/.style=concept, concept color=popblue,
    level 1/.append style={level distance=4.5cm, sibling angle=120},
    level 2/.append style={level distance=3cm, sibling angle=60}]
\node[concept color=popdark] {Leçon E14}
    child[concept color=poporange] { node {Verbos defectivos}
        child { node {gustar / doler} }
        child { node {parecer / soler} }
        child { node {concernir / atañer} }
        child { node {llover / nevar} }
    }
    child[concept color=popgreen] { node {Restricción}
        child { node {no... más que} }
        child { node {no... sino} }
        child { node {sólo / únicamente} }
        child { node {apenas} }
    }
    child[concept color=poppurple] { node {Discurso indirecto}
        child { node {Personnes} }
        child { node {Temps (Concordance)} }
        child { node {Indicateurs} }
        child { node {Verbes introducteurs} }
    };
\end{tikzpicture}
\legende{Carte mentale : les trois piliers grammaticaux de la leçon E14}
\end{popfigure}

\begin{attention}[Pièges à éviter]
Ne dis jamais \esp{los alumnos gustan} pour « les élèves aiment » : avec \esp{gustar}, le sujet grammatical est la chose aimée : \esp{a los alumnos les gusta la idea}. N'oublie pas la transformation des personnes : \esp{tenemos} (nous) devient \esp{tenían} (ils) dans la bouche du journaliste. Enfin, \esp{sino} (et non \esp{si no}) s'emploie après une négation pour introduire la rectification. Attention à l'accent sur \esp{sólo} (adverbe = seulement) vs \esp{solo} (adjectif = seul) — l'usage académique conserve l'accent pour lever l'ambiguïté. Vérifie l'accord du verbe défectif : \esp{me duele la cabeza} (sg) mais \esp{me duelen los pies} (pl).
\end{attention}

\courssub{Proposition de production et corrigé commenté}

\begin{exoresolu}[Crónica modelo — La Voz del Patio]
\textbf{Énoncé.} Rédige la crónica de la reunión del Club Hispanófilo (12-15 líneas) en respetando el contrato : 3 verbos defectivos, 2 restricciones, 4 frases en estilo indirecto.

\tcblower

\textbf{Production modèle :}

\emph{El Club Hispanófilo prepara la Semana de la Lengua}

Ayer a las cuatro de la tarde, el Club Hispanófilo se reunió en la sala 12 para preparar la Semana de la Lengua Española. El presidente, Etienne, abrió la sesión y anunció que la Semana tendría lugar del 12 al 16 de mayo. Añadió que al día siguiente recibirían a una delegación de estudiantes de Guinea Ecuatorial, lo que alegró a todos los miembros. Después, precisó que ese año no organizarían una feria del libro, sino un concurso de poesía y una noche de teatro. Sobre el presupuesto, recordó que el club no tenía más que 25 000 francos para toda la semana, y explicó que apenas alcanzaba para los premios.

Las reacciones fueron variadas. A muchos alumnos les gusta la idea del concurso de poesía, porque les encanta escribir. Sin embargo, a algunos les duele que no haya feria del libro, pues era su actividad favorita. A otros no les parece suficiente el presupuesto y proponen buscar patrocinadores en Douala. Como el club suele reunirse sólo los jueves, el presidente pidió que todos volvieran el lunes para ultimar los preparativos. La Semana de la Lengua promete ser inolvidable.

\textbf{Commentaire des choix de langue :}

\begin{itemize}
\item \textbf{Discours indirect (4 occurrences, 4 introducteurs différents) :} \esp{anunció que... tendría} (futur $\rightarrow$ conditionnel) ; \esp{añadió que... recibirían} (futur $\rightarrow$ conditionnel, \esp{mañana} $\rightarrow$ \esp{al día siguiente}) ; \esp{precisó que... no organizarían... sino...} ; \esp{recordó que... no tenía más que} (présent $\rightarrow$ imparfait, \esp{tenemos} $\rightarrow$ \esp{tenían}) ; plus \esp{explicó que... alcanzaba} et \esp{pidió que... volvieran} (subjonctif imparfait après verbe de volonté au passé).
\item \textbf{Verbes défectifs (4) :} \esp{les gusta la idea} (accord avec \esp{la idea}, singulier) ; \esp{les duele que no haya feria} (subjonctif après \esp{doler} exprimant le sentiment) ; \esp{no les parece suficiente el presupuesto} ; \esp{suele reunirse} (verbe d'habitude).
\item \textbf{Restrictions (2 procédés) :} \esp{no tenía más que 25 000 francos} et \esp{no organizarían una feria del libro, sino un concurso} ; en bonus, \esp{apenas alcanzaba} et \esp{sólo los jueves}.
\item \textbf{Personnes et indicateurs transposés :} \esp{este año} $\rightarrow$ \esp{ese año} ; \esp{mañana} $\rightarrow$ \esp{al día siguiente} ; le journaliste dit \esp{el club} et \esp{ellos}, jamais \esp{nosotros}.
\item \textbf{Structure journalistique :} titre, phrase d'attaque (qui, quoi, où, quand), corps, réactions, conclusion ouverte.
\end{itemize}
\end{exoresolu}

\courssub{Bilan de l'activité}

Cette activité t'a permis de mobiliser dans une production authentique les trois savoir-faire de la leçon. Retiens l'essentiel :

\begin{itemize}
\item Les \cle{verbes défectifs} comme \esp{gustar}, \esp{doler}, \esp{parecer} se construisent avec \esp{a} + la personne concernée et un pronom d'objet indirect ; le verbe s'accorde avec la chose, pas avec la personne.
\item La \cle{restriction} se traduit selon le sens : \esp{no... más que} ou \esp{sólo} pour « seulement », \esp{no... sino} pour « non pas... mais », \esp{apenas} pour « à peine ».
\item Le passage au \cle{discours indirect} exige trois transformations simultanées : les personnes, les temps (concordance) et les indicateurs spatio-temporels ; varier les verbes introducteurs enrichit le style journalistique.
\end{itemize}

Les meilleurs articles, corrigés après la relecture croisée, seront affichés au mur de la salle d'espagnol et publiés dans le prochain numéro de \emph{La Voz del Patio}. Garde ta grille d'auto-évaluation : elle te servira de liste de vérification pour l'épreuve de thème du baccalauréat.

\newpage

\coursec{Méthodes et exercices résolus}

\courssub{Savoir-faire 1 : Conjuguer et employer les verbes défectifs}

\begin{definition}[Qu'est-ce qu'un verbe défectif ?]
Un verbe \cle{défectif} est un verbe qui ne se conjugue pas à toutes les personnes, ou dont la construction diffère du français. En espagnol, les plus importants sont : \esp{gustar} « plaire à », \esp{doler} « faire mal à », \esp{parecer} « sembler », \esp{soler} « avoir l'habitude de », \esp{concernir} et \esp{atañer} « concerner », \esp{encantar} « adorer », \esp{faltar} « manquer », \esp{quedar} « rester », et les verbes impersonnels comme \esp{llover} « pleuvoir », \esp{nevar} « neiger ».
\end{definition}

\begin{methode}[Traduire une phrase avec un verbe défectif]
Les verbes défectifs (\esp{gustar}, \esp{doler}, \esp{parecer}, \esp{soler}, \esp{concernir}, \esp{atañer}, \esp{llover}…) ne fonctionnent pas comme les verbes français : le sujet grammatical est souvent la \emph{chose}, et la personne devient complément d'objet indirect (COI). Voici la démarche en cinq étapes.
\begin{enumerate}
\item \textbf{Repérer la structure française.} Dans « ce livre me plaît », le sujet français est « ce livre » : l'espagnol suit la même logique que « me plaît » et non que « j'aime ».
\item \textbf{Identifier le sujet espagnol.} C'est la chose qui plaît, qui fait mal, qui semble : \esp{el libro}, \esp{la cabeza}, \esp{la propuesta}. Le verbe s'accorde avec elle (singulier ou pluriel).
\item \textbf{Placer le pronom COI devant le verbe.} \esp{me, te, le, nos, os, les}. On peut renforcer par \esp{a mí, a ti, a él, a María…} : \esp{A mí me gusta el café.}
\item \textbf{Choisir la bonne forme verbale.} \esp{gustar} et \esp{doler} s'emploient surtout à la 3\ieme{} personne : \esp{me gusta / me gustan}, \esp{me duele / me duelen}. \esp{soler} se conjugue normalement mais signifie « avoir l'habitude de » : \esp{Suelo levantarme temprano.} \esp{llover} n'existe qu'à la 3\ieme{} personne du singulier : \esp{Llueve mucho en Douala.}
\item \textbf{Vérifier le temps et l'accord.} Au passé : \esp{me gustó / me gustaron}. Avec un infinitif, toujours le singulier : \esp{Me gusta bailar.}
\end{enumerate}
Structures à retenir : \esp{Me duele la cabeza} « j'ai mal à la tête » ; \esp{Esto no te concierne} « cela ne te concerne pas » ; \esp{El asunto atañe a todos} « l'affaire concerne tout le monde » ; \esp{Suele llover en octubre} « il pleut souvent en octobre ».
\end{methode}

\begin{propriete}[Conjugaison utile des principaux verbes défectifs]
\begin{center}
\begin{tabular}{|l|l|l|l|}
\hline
\rowcolor{popblueL} Verbe & Présent & Imparfait & Passé simple \\
\hline
\esp{gustar} & \esp{gusta / gustan} & \esp{gustaba / gustaban} & \esp{gustó / gustaron} \\
\hline
\esp{doler} (o $\to$ ue) & \esp{duele / duelen} & \esp{dolía / dolían} & \esp{dolió / dolieron} \\
\hline
\esp{parecer} & \esp{parece / parecen} & \esp{parecía / parecían} & \esp{pareció / parecieron} \\
\hline
\esp{soler} (o $\to$ ue) & \esp{suelo, sueles, suele,} \esp{solemos, soléis, suelen} & \esp{solía, solías, solía…} & \esp{--} (rare) \\
\hline
\esp{llover} (impersonnel) & \esp{llueve} & \esp{llovía} & \esp{llovió} \\
\hline
\end{tabular}
\end{center}
Remarques : \esp{soler} ne s'emploie guère au passé simple ; on lui préfère l'imparfait \esp{solía} « j'avais l'habitude de ». \esp{parecer} à la 1\iere{} personne fait \esp{parezco} (comme \esp{conocer} $\to$ \esp{conozco}), mais cet emploi reste rare.
\end{propriete}

\begin{exemplebox}[Les verbes défectifs en contexte]
\esp{¿Te gustan los platos cameruneses? — Sí, me encanta el ndolé.} « Est-ce que tu aimes les plats camerounais ? — Oui, j'adore le ndolé. »

\esp{A María le duelen los pies después del partido.} « Marie a mal aux pieds après le match. »

\esp{Nos parece muy interesante la propuesta del alcalde.} « La proposition du maire nous semble très intéressante. »

\esp{Me falta dinero para el transporte.} « Il me manque de l'argent pour le transport. »

\esp{Quedan dos días para el examen.} « Il reste deux jours avant l'examen. »
\end{exemplebox}

\begin{exoresolu}[Thème : verbes défectifs]
Traduisez en espagnol : 1. J'aime beaucoup les romans de cet écrivain. 2. Mes jambes me font mal après le match. 3. D'habitude, nous nous levons à six heures. 4. Cette affaire ne vous concerne pas. 5. Il pleut souvent à Douala en août.
\tcblower
\textbf{Raisonnement phrase par phrase.}
\begin{enumerate}
\item « J'aime » avec une chose : structure \esp{gustar}. Sujet espagnol = \esp{las novelas} (pluriel) donc \esp{gustan} ; personne = COI \esp{me}. Solution : \esp{Me gustan mucho las novelas de este escritor.}
\item « Mes jambes me font mal » : \esp{doler}, sujet pluriel \esp{las piernas} donc \esp{duelen} (verbe en \esp{-er} à diphtongaison o $\to$ ue). Solution : \esp{Me duelen las piernas después del partido.}
\item « D'habitude » : \esp{soler} + infinitif, conjugué normalement à la 1\iere{} personne du pluriel, diphtongaison : \esp{solemos}. Solution : \esp{Solemos levantarnos a las seis.}
\item « Concerner » : \esp{concernir} (ou \esp{atañer}), sujet = \esp{este asunto}, COI \esp{te} renforcé par \esp{a ti}. Solution : \esp{Este asunto no te concierne (a ti).}
\item Verbe impersonnel : 3\ieme{} personne du singulier uniquement. Solution : \esp{Llueve a menudo en Douala en agosto.} (ou \esp{Suele llover en Douala en agosto.})
\end{enumerate}
\textbf{Conclusion.} Le réflexe essentiel : chercher d'abord le \emph{sujet espagnol} (la chose), accorder le verbe avec lui, puis placer le pronom COI. Ne jamais calquer « je » $\to$ \esp{yo} avec \esp{gustar}.
\end{exoresolu}

\courssub{Savoir-faire 2 : Traduire la restriction}

\begin{methode}[Traduire « seulement », « ne… que », « à peine »]
\begin{enumerate}
\item \textbf{Repérer le mot français de restriction} : seulement, uniquement, ne… que, ne… rien d'autre que, à peine.
\item \textbf{Choisir l'équivalent espagnol selon le registre} :
\begin{itemize}
\item \esp{sólo} (ou \esp{solamente}) = seulement : \esp{Sólo tengo cinco mil francos.}
\item \esp{únicamente} = uniquement (registre soutenu).
\item \esp{no… más que} = ne… que : \esp{No tengo más que cinco mil francos.}
\item \esp{no… sino} = ne… que / mais (après négation, avec opposition) : \esp{No quiero dinero, sino justicia.}
\item \esp{apenas} = à peine : \esp{Apenas tengo tiempo para comer.}
\end{itemize}
\item \textbf{Placer le mot restrictif juste devant l'élément restreint} : \esp{Sólo vino María} (seule Marie est venue) $\neq$ \esp{María sólo vino} (Marie n'a fait que venir).
\item \textbf{Vérifier la négation double} : avec \esp{no… más que}, le \esp{no} précède le verbe et \esp{más que} précède l'élément restreint : \esp{No leí más que dos capítulos.}
\item \textbf{Traduire « à peine… que »} (temporel) par \esp{apenas… cuando} ou \esp{en cuanto} : \esp{Apenas llegó, empezó a llover.} « À peine fut-il arrivé qu'il se mit à pleuvoir. »
\end{enumerate}
\end{methode}

\begin{propriete}[Tableau comparatif des restrictions]
\begin{center}
\begin{tabularx}{\linewidth}{|X|X|X|}
\hline
\rowcolor{popblueL} Français & Español & Remarque \\
\hline
seulement & \esp{sólo / solamente} & L'accent sur \esp{sólo} est facultatif depuis 2010, mais recommandé s'il y a ambiguïté avec l'adjectif \esp{solo} « seul ». \\
\hline
uniquement & \esp{únicamente} & Registre soutenu, écrit administratif. \\
\hline
ne… que & \esp{no… más que} & Ne jamais oublier le \esp{no} devant le verbe. \\
\hline
non pas… mais & \esp{no… sino} & Opposition entre deux éléments ; \esp{sino que} devant une proposition verbale : \esp{No vino, sino que llamó.} \\
\hline
à peine & \esp{apenas} & Inversion fréquente du sujet : \esp{Apenas llegó María…} \\
\hline
\end{tabularx}
\end{center}
\end{propriete}

\begin{exoresolu}[Version : la phrase restrictive]
Traduisez en français : \esp{En el mercado de Mokolo, no vendían más que frutas y verduras. Sólo los lunes había pescado fresco, y apenas llegaban los clientes cuando ya no quedaba nada. No buscaban lujo, sino productos de calidad.}
\tcblower
\textbf{Analyse des structures restrictives.}
\begin{itemize}
\item \esp{no vendían más que frutas y verduras} : structure \esp{no… más que} = « ne… que ».
\item \esp{Sólo los lunes} : restriction temporelle = « seulement le lundi / uniquement les lundis ».
\item \esp{apenas llegaban los clientes cuando…} : « à peine les clients arrivaient-ils que… » ; noter l'inversion du sujet après \esp{apenas}, fréquente en espagnol.
\item \esp{no buscaban lujo, sino productos de calidad} : \esp{no… sino} = « non pas… mais ».
\end{itemize}
\textbf{Traduction rédigée.} « Au marché de Mokolo, on ne vendait que des fruits et des légumes. Ce n'est que le lundi qu'il y avait du poisson frais, et à peine les clients arrivaient-ils qu'il ne restait déjà plus rien. Ils ne recherchaient pas le luxe, mais des produits de qualité. »
\textbf{Conclusion.} Chaque mot restrictif espagnol a un équivalent français précis ; repérer la structure négative complète (\esp{no… más que}, \esp{no… sino}) avant de traduire évite les contresens.
\end{exoresolu}

\courssub{Savoir-faire 3 : Passer du discours direct au discours indirect}

\begin{methode}[Transformer une phrase du style direct au style indirect]
\begin{enumerate}
\item \textbf{Identifier le verbe introducteur et son temps} : \esp{decir, afirmar, explicar, preguntar, pedir, mandar, aconsejar…} Au présent ou futur, les temps ne changent pas ; au passé (\esp{dijo, preguntó, mandó}), il y a concordance.
\item \textbf{Supprimer les guillemets et introduire la subordonnée} : \esp{que} pour une déclaration, \esp{si} pour une question totale (oui/non), le mot interrogatif (\esp{qué, quién, dónde, cuándo, por qué}) pour une question partielle, \esp{que} + subjonctif ou infinitif pour un ordre.
\item \textbf{Transposer les personnes} : 1\iere{} et 2\ieme{} personnes s'adaptent au nouveau contexte (je $\to$ il/elle, tu $\to$ je ou il), ainsi que les possessifs (\esp{mi} $\to$ \esp{su}) et les pronoms.
\item \textbf{Appliquer la concordance des temps} (verbe introducteur au passé) : présent $\to$ imparfait (\esp{trabajo} $\to$ \esp{trabajaba}) ; passé composé / imparfait $\to$ plus-que-parfait (\esp{he trabajado} $\to$ \esp{había trabajado}) ; futur $\to$ conditionnel (\esp{trabajaré} $\to$ \esp{trabajaría}) ; impératif $\to$ subjonctif imparfait ou infinitif (\esp{¡trabaja!} $\to$ \esp{que trabajara / trabajar}).
\item \textbf{Adapter les indicateurs spatio-temporels} : \esp{aquí} $\to$ \esp{allí} ; \esp{ahora} $\to$ \esp{entonces} ; \esp{hoy} $\to$ \esp{aquel día} ; \esp{ayer} $\to$ \esp{el día anterior} ; \esp{mañana} $\to$ \esp{al día siguiente} ; \esp{este} $\to$ \esp{aquel/ese}.
\end{enumerate}
\end{methode}

\begin{aretenir}[Tableau de concordance des temps (verbe introducteur au passé)]
\begin{center}
\begin{tabular}{|l|l|}
\hline
\rowcolor{popblueL} Discours direct & Discours indirect \\
\hline
présent : \esp{trabajo} & imparfait : \esp{trabajaba} \\
\hline
passé composé : \esp{he trabajado} & plus-que-parfait : \esp{había trabajado} \\
\hline
imparfait : \esp{trabajaba} & plus-que-parfait : \esp{había trabajado} \\
\hline
futur : \esp{trabajaré} & conditionnel : \esp{trabajaría} \\
\hline
impératif : \esp{¡trabaja!} & subjonctif imparfait : \esp{que trabajara} \\
\hline
\end{tabular}
\end{center}
Si le verbe introducteur est au présent (\esp{dice que…}), aucun changement de temps : \esp{Dice que trabaja hoy.}
\end{aretenir}

\begin{exoresolu}[Transformation : discours indirect]
Mettez au discours indirect, en commençant par \esp{El profesor dijo que…} / \esp{preguntó si…} / \esp{mandó que…} : 1. \esp{El profesor dijo: «Mañana corregiré vuestros ejercicios aquí».} 2. \esp{Preguntó: «¿Habéis entendido la lección?»} 3. \esp{Mandó: «¡Abrid los libros ahora!»}
\tcblower
\textbf{Étape 1.} Verbe introducteur au passé simple (\esp{dijo, preguntó, mandó}) : concordance obligatoire.
\textbf{Étape 2.} Phrase 1 : déclaration $\to$ \esp{que}. Futur \esp{corregiré} $\to$ conditionnel \esp{corregiría} ; \esp{vuestros} $\to$ \esp{nuestros} (les élèves rapportent) ou \esp{sus} ; \esp{mañana} $\to$ \esp{al día siguiente} ; \esp{aquí} $\to$ \esp{allí}. Solution : \esp{El profesor dijo que corregiría nuestros ejercicios allí al día siguiente.}
Phrase 2 : question totale $\to$ \esp{si}. Passé composé \esp{habéis entendido} $\to$ plus-que-parfait \esp{habíamos entendido} (ou \esp{habían entendido} selon le rapporteur). Solution : \esp{Preguntó si habíamos entendido la lección.}
Phrase 3 : ordre $\to$ \esp{que} + subjonctif imparfait : \esp{abrir} $\to$ \esp{abriéramos} ; \esp{ahora} $\to$ \esp{entonces}. Solution : \esp{Mandó que abriéramos los libros entonces.} (Variante acceptée : \esp{Nos mandó abrir los libros enseguida.})
\textbf{Conclusion.} Trois transformations systématiques : personnes, temps, indicateurs spatio-temporels. Dresser mentalement le tableau de concordance avant d'écrire.
\end{exoresolu}

\courssub{Savoir-faire 4 : Choisir le verbe introducteur et la bonne subordonnée}

\begin{methode}[Déclarer, interroger, ordonner au discours indirect]
\begin{enumerate}
\item \textbf{Classer la phrase de départ.} Affirmation $\to$ verbe déclaratif + \esp{que} + indicatif : \esp{decir, afirmar, anunciar, explicar, contestar, añadir que…}
\item \textbf{Question totale} (réponse oui/non) $\to$ \esp{preguntar si} + indicatif : \esp{Me preguntó si estaba cansado.}
\item \textbf{Question partielle} $\to$ \esp{preguntar} + mot interrogatif (avec accent !) + indicatif : \esp{Preguntó dónde vivía yo.} Attention : \esp{dónde, qué, quién, cuándo, cuánto} gardent l'accent au discours indirect.
\item \textbf{Ordre, conseil, souhait} $\to$ verbe de volonté + \esp{que} + subjonctif : \esp{mandar, ordenar, pedir, aconsejar, recomendar, prohibir que…} : \esp{El médico le aconsejó que descansara.}
\item \textbf{Alternative avec infinitif} : quand le sujet de l'ordre est un COI, l'espagnol préfère souvent \esp{verbe + COI + infinitif} : \esp{Le pidió callarse} « il lui demanda de se taire » ; \esp{Nos mandó salir.}
\end{enumerate}
Réflexe Bac : dans une version, repérer le verbe introducteur espagnol pour choisir le bon verbe français (« il affirma que », « il demanda si », « il ordonna de ») ; dans un thème, faire l'inverse.
\end{methode}

\begin{exoresolu}[Thème : verbes introducteurs et subordonnées]
Traduisez en espagnol : 1. Le maire a annoncé que le marché ouvrirait le lendemain. 2. Elle m'a demandé si j'avais visité la Guinée équatoriale. 3. Le professeur nous a conseillé de réviser les verbes défectifs. 4. Il m'a demandé où j'habitais.
\tcblower
\textbf{Raisonnement.}
\begin{enumerate}
\item « Annoncer que » $\to$ \esp{anunciar que} ; futur français « ouvrirait » (futur dans le passé) $\to$ conditionnel espagnol \esp{abriría} ; « le lendemain » $\to$ \esp{al día siguiente}. Solution : \esp{El alcalde anunció que el mercado abriría al día siguiente.}
\item Question totale rapportée $\to$ \esp{preguntar si} ; passé composé $\to$ plus-que-parfait \esp{había visitado}. Solution : \esp{Me preguntó si había visitado Guinea Ecuatorial.}
\item Conseil $\to$ \esp{aconsejar} + COI + infinitif (structure la plus naturelle) : \esp{El profesor nos aconsejó repasar los verbos defectivos.} (Variante : \esp{…aconsejó que repasáramos los verbos defectivos.})
\item Question partielle $\to$ mot interrogatif accentué \esp{dónde} ; « j'habitais » (imparfait rapporté) $\to$ \esp{vivía}. Solution : \esp{Me preguntó dónde vivía.}
\end{enumerate}
\textbf{Conclusion.} Le choix du verbe introducteur détermine toute la structure : \esp{que} + indicatif, \esp{si} + indicatif, mot interrogatif accentué, ou \esp{que} + subjonctif / infinitif pour la volonté.
\end{exoresolu}

\courssub{Savoir-faire 5 : Entraînement intégré — version de type Baccalauréat}

\begin{methode}[Aborder une version mêlant les trois notions]
\begin{enumerate}
\item \textbf{Lire tout le texte} et souligner : verbes défectifs, mots de restriction, verbes introducteurs et discours rapportés.
\item \textbf{Traduire d'abord le sens global}, phrase par phrase, sans calquer mot à mot.
\item \textbf{Appliquer les trois réflexes} : sujet grammatical pour les défectifs ; structure négative complète pour la restriction ; concordance des temps pour le discours indirect.
\item \textbf{Rédiger en français soigné}, naturel, sans espagnolismes.
\item \textbf{Relire} : accords, temps français (imparfait, plus-que-parfait, conditionnel du discours rapporté), ponctuation.
\end{enumerate}
\end{methode}

\begin{exoresolu}[Version intégrée]
Traduisez en français : \esp{El alcalde anunció que sólo quedaban dos días para inscribirse en el centro cultural. Explicó que a los jóvenes les gustaban mucho los talleres de teatro y que no pedía más que un poco de entusiasmo. Añadió que solía reunirse con los vecinos los viernes y les aconsejó que no esperaran al último momento.}
\tcblower
\textbf{Repérage.} Défectifs : \esp{les gustaban} (sujet = \esp{los talleres}), \esp{solía reunirse}. Restriction : \esp{sólo}, \esp{no pedía más que}. Discours indirect : \esp{anunció que}, \esp{explicó que}, \esp{añadió que}, \esp{aconsejó que} + subjonctif imparfait \esp{esperaran}.
\textbf{Analyse grammaticale.}
\begin{itemize}
\item \esp{quedaban dos días} : « il restait deux jours » (imparfait de concordance).
\item \esp{a los jóvenes les gustaban los talleres} : sujet pluriel \esp{los talleres} $\to$ « les jeunes aimaient beaucoup les ateliers ».
\item \esp{no pedía más que un poco de entusiasmo} : « il ne demandait qu'un peu d'enthousiasme ».
\item \esp{les aconsejó que no esperaran} : subjonctif imparfait après verbe de conseil $\to$ « il leur conseilla de ne pas attendre ».
\end{itemize}
\textbf{Traduction rédigée.} « Le maire annonça qu'il ne restait que deux jours pour s'inscrire au centre culturel. Il expliqua que les jeunes aimaient beaucoup les ateliers de théâtre et qu'il ne demandait qu'un peu d'enthousiasme. Il ajouta qu'il avait l'habitude de rencontrer les habitants le vendredi et leur conseilla de ne pas attendre le dernier moment. »
\textbf{Conclusion.} Une version réussie combine les trois réflexes de la leçon : identifier le sujet des verbes défectifs, restituer la restriction par « ne… que / seulement », et respecter la concordance des temps du discours rapporté en français (imparfait, conditionnel).
\end{exoresolu}

\begin{exercice}[Thème intégré : à faire en classe]
Traduisez en espagnol : 1. Mes parents aiment beaucoup la musique makossa. 2. Il ne me reste que deux jours pour réviser. 3. Le directeur a dit qu'il ne voulait pas de retard, mais de la ponctualité. 4. Elle m'a demandé si j'avais mal à la tête. 5. D'habitude, il pleut beaucoup à Limbé en juillet.
\end{exercice}

\begin{corrige}[Thème intégré]
\begin{enumerate}
\item \esp{A mis padres les gusta mucho la música makossa.} (sujet singulier : \esp{la música} ; COI \esp{les} renforcé par \esp{a mis padres}.)
\item \esp{Sólo me quedan dos días para repasar.} ou \esp{No me quedan más que dos días para repasar.} (sujet pluriel : \esp{dos días} $\to$ \esp{quedan}.)
\item \esp{El director dijo que no quería retrasos, sino puntualidad.} (concordance : présent $\to$ imparfait ; opposition $\to$ \esp{sino}.)
\item \esp{Me preguntó si me dolía la cabeza.} (question totale $\to$ \esp{si} ; \esp{doler} à l'imparfait, sujet \esp{la cabeza}.)
\item \esp{Suele llover mucho en Limbé en julio.} (\esp{soler} + infinitif ; verbe impersonnel au singulier.)
\end{enumerate}
\end{corrige}

\begin{attention}[Les 5 erreurs qui coûtent des points au Bac]
\begin{enumerate}
\item \textbf{Calquer « j'aime » par \esp{yo gusto}.} Faute grave : \esp{gustar} exige le COI et le sujet-chose. On écrit \esp{Me gusta el fútbol}, jamais \esp{Yo gusto el fútbol}. Même vigilance pour \esp{doler} : \esp{Me duele la cabeza}, pas \esp{Yo duelo}.
\item \textbf{Oublier la concordance des temps au discours indirect.} Après un verbe introducteur au passé, le présent devient imparfait : \esp{Dijo que venía} (et non \esp{dijo que viene}). Le futur devient conditionnel : \esp{Anunció que vendría}.
\item \textbf{Perdre l'accent des mots interrogatifs au style indirect.} On écrit \esp{Me preguntó dónde vivía}, avec accent sur \esp{dónde} ; sans accent, \esp{donde} est une simple relative. De même \esp{qué, quién, cuándo, cuánto}.
\item \textbf{Traduire « ne… que » par \esp{no… sólo} ou oublier le \esp{no}.} La structure correcte est \esp{no… más que} : \esp{No tengo más que diez euros.} Et attention au faux ami : \esp{actualmente} signifie « actuellement », pas « en fait » ; \esp{únicamente} = « uniquement ».
\item \textbf{Négliger les indicateurs spatio-temporels et les accords.} \esp{mañana} rapporté devient \esp{al día siguiente}, \esp{aquí} devient \esp{allí}, \esp{mi} devient \esp{su}. Enfin, accorder le verbe défectif avec son vrai sujet : \esp{Me gustan las novelas} (pluriel), jamais \esp{me gusta las novelas}.
\end{enumerate}
\end{attention}

\newpage

\coursec{Exercices}

\courssub{Je vérifie mes connaissances}

\begin{exercice}[QCM : dépiste les pièges]
Choisis la bonne réponse pour chaque phrase.
\begin{enumerate}
\item \esp{A mí \_\_\_\_ los partidos de fútbol.}
\begin{itemize}
\item[a)] \esp{me gusto}
\item[b)] \esp{me gusta}
\item[c)] \esp{me gustan}
\end{itemize}
\item \esp{En la reunión no habló \_\_\_\_ el director.}
\begin{itemize}
\item[a)] \esp{más que}
\item[b)] \esp{sino}
\item[c)] \esp{apenas}
\end{itemize}
\item \esp{No compré pan \_\_\_\_ arroz.}
\begin{itemize}
\item[a)] \esp{más que}
\item[b)] \esp{sino}
\item[c)] \esp{solamente}
\end{itemize}
\item La directrice a dit : \esp{« Recibiré a los delegados mañana »}. Au discours indirect :
\begin{itemize}
\item[a)] \esp{La directora dijo que recibirá a los delegados mañana.}
\item[b)] \esp{La directora dijo que recibiría a los delegados al día siguiente.}
\item[c)] \esp{La directora dijo que recibió a los delegados ayer.}
\end{itemize}
\item \esp{Ayer \_\_\_\_ mucho en Douala.}
\begin{itemize}
\item[a)] \esp{lloví}
\item[b)] \esp{llovió}
\item[c)] \esp{llovieron}
\end{itemize}
\end{enumerate}
\end{exercice}

\begin{exercice}[Vrai ou faux ? Justifie]
Réponds par vrai ou faux, puis justifie ta réponse en une phrase en français.
\begin{enumerate}
\item Le verbe \esp{gustar} se construit comme le verbe français « aimer » : \esp{Yo gusto el café}.
\item \esp{Me duelen los pies} signifie « J'ai mal aux pieds ».
\item \esp{Soler} exprime une habitude et se traduit souvent par « avoir l'habitude de ».
\item Au discours indirect, \esp{aquí} devient \esp{allí} et \esp{ahora} devient \esp{entonces}.
\item \esp{No tengo más que dos hermanos} signifie « Je n'ai pas deux frères, mais trois ».
\end{enumerate}
\end{exercice}

\begin{exercice}[Vocabulaire : les verbes introducteurs]
Associe chaque verbe espagnol à son équivalent français, puis complète les phrases avec le verbe qui convient au passé simple.
\begin{center}
\begin{tabularx}{\linewidth}{|X|X|}
\hline
\rowcolor{popblueL} \textbf{Español} & \textbf{Français} \\
\hline
\esp{decir} & ordonner \\
\hline
\esp{preguntar} & répondre \\
\hline
\esp{contestar} & dire \\
\hline
\esp{mandar} & demander \\
\hline
\esp{exclamar} & s'exclamer \\
\hline
\end{tabularx}
\end{center}
\begin{enumerate}
\item El profesor \_\_\_\_ que el examen sería al día siguiente.
\item Los alumnos \_\_\_\_ si podían salir antes.
\item El jefe \_\_\_\_ que cerraran la puerta.
\item Ella \_\_\_\_ : « \esp{¡Qué alegría verte!} »
\end{enumerate}
\end{exercice}

\begin{exercice}[Conjugaison à compléter]
Complète avec la forme correcte du verbe entre parenthèses.
\begin{enumerate}
\item A nosotros \_\_\_\_ (gustar) las clases de español.
\item A mí \_\_\_\_ (doler) la cabeza desde ayer.
\item ¿A ti te \_\_\_\_ (parecer) bien la idea?
\item Mi abuela \_\_\_\_ (soler) ir al mercado de Mvog-Mbi los sábados.
\item Este asunto no \_\_\_\_ (concernir) a los alumnos de primero.
\item En diciembre \_\_\_\_ (llover) poco en Yaoundé.
\end{enumerate}
\end{exercice}

\courssub{Je m'entraîne}

\begin{exercice}[Traduction : phrases isolées]
Traduis en espagnol les phrases suivantes.
\begin{enumerate}
\item J'aime beaucoup la musique camerounaise.
\item Mes frères ont mal aux jambes après le match.
\item Nous n'avons que dix minutes pour finir l'exercice.
\item Ce n'est pas un stylo, mais un crayon.
\item Il a l'habitude de se lever à cinq heures du matin.
\end{enumerate}
\end{exercice}

\begin{exercice}[Du discours direct au discours indirect]
Transpose les phrases suivantes au discours indirect en commençant par l'indication donnée. Attention aux personnes, aux temps et aux indicateurs spatio-temporels.
\begin{enumerate}
\item \esp{« Vendré mañana a tu casa »} $\rightarrow$ \esp{Juan dijo que...}
\item \esp{« ¿Dónde está mi libro? »} $\rightarrow$ \esp{María preguntó...}
\item \esp{« No toquéis nada aquí »} $\rightarrow$ \esp{El policía ordenó que...}
\item \esp{« Ahora estoy muy ocupado »} $\rightarrow$ \esp{El director contestó que...}
\item \esp{« Ayer fuimos al mercado »} $\rightarrow$ \esp{Los niños dijeron que...}
\item \esp{« ¿Has terminado el trabajo hoy? »} $\rightarrow$ \esp{La profesora preguntó a Pedro...}
\end{enumerate}
\end{exercice}

\begin{exercice}[Texte à trous]
Complète le texte suivant avec la forme correcte des verbes \esp{gustar}, \esp{doler}, \esp{parecer} et \esp{soler}, ou avec le déictique qui convient au discours indirect (\esp{ese}, \esp{allí}, \esp{aquel}, \esp{entonces}).

\begin{textebox}[Una entrevista en el instituto]
Ayer, una periodista visitó nuestro instituto de Yaoundé. Nos dijo que \_\_\_\_ (1) día quería hablar con los alumnos de Terminale. A nosotros nos \_\_\_\_ (2) mucho la idea, porque nos \_\_\_\_ (3) interesante hablar con una profesional. Ella explicó que \_\_\_\_ (4) entrevistar a jóvenes en las escuelas de la capital. Un alumno contestó que no podía quedarse porque le \_\_\_\_ (5) mucho la cabeza. La periodista dijo que volvería \_\_\_\_ (6) semana siguiente y que \_\_\_\_ (7) podríamos hacerle todas las preguntas.
\end{textebox}
\end{exercice}

\begin{exercice}[Reformulation : la restriction]
Réécris chaque phrase en remplaçant \esp{sólo} par \esp{no... más que}, ou l'inverse, sans changer le sens.
\begin{enumerate}
\item \esp{Sólo tengo cinco mil francos.}
\item \esp{No vinieron más que tres delegados.}
\item \esp{Sólo hablamos español en esta clase.}
\item \esp{No queda más que un poco de arroz.}
\item \esp{Únicamente los delegados pueden entrar.}
\item \esp{No leí más que el primer capítulo.}
\end{enumerate}
\end{exercice}

\courssub{Je me prépare au Bac}

\begin{exercice}[Version : une réunion de parents d'élèves]
Traduis en français le texte suivant, puis justifie trois de tes choix de traduction (un verbe défectif, une tournure restrictive, une transposition au discours indirect).

\begin{textebox}[La reunión de padres]
El sábado pasado hubo una reunión de padres de alumnos en el liceo de Biyem-Assi. El director abrió la sesión y dijo que aquel año los resultados eran mejores, pero que todavía quedaba mucho trabajo. Añadió que sólo los padres inscritos podían votar al final de la reunión. Una madre preguntó si el comedor funcionaría al día siguiente, porque a su hijo no le gustaba la comida de la calle y le dolía a menudo el estómago. El director contestó que no se trataba de un problema de cocina, sino de organización, y que apenas tenían dinero para contratar a más personal. Otro padre declaró que a él le parecía urgente reparar las aulas antes de la temporada de lluvias. Al final, el director prometió que estudiaría todas las propuestas y que daría una respuesta la semana siguiente.
\end{textebox}

\emph{Barème indicatif : traduction 14 points ; justification des trois choix 6 points.}
\end{exercice}

\begin{exercice}[Thème type Bac]
Traduis en espagnol les cinq phrases suivantes. Soigne les verbes défectifs, les tournures restrictives et le discours indirect.
\begin{enumerate}
\item La directrice a dit qu'elle ne recevrait que les délégués le lendemain.
\item Mes camarades aiment beaucoup ce professeur parce que ses cours sont intéressants.
\item Ce matin-là, il pleuvait beaucoup et les élèves sont arrivés en retard.
\item Le professeur a demandé aux élèves s'ils avaient fini leurs devoirs ce jour-là.
\item Ce n'est pas la fatigue qui m'inquiète, mais ta santé : tu as l'habitude de dormir très peu.
\end{enumerate}
\emph{Barème indicatif : 4 points par phrase.}
\end{exercice}

\begin{exercice}[Compréhension et rédaction]
Lis le texte suivant, réponds aux questions en espagnol, puis traite le sujet de rédaction.

\begin{textebox}[El consejo de disciplina]
Ayer por la mañana, el consejo de disciplina del colegio se reunió para estudiar el caso de dos alumnos. El presidente del consejo anunció que no castigaría más que a los responsables de la pelea, y no a los simples testigos. Uno de los alumnos explicó que no había empezado la pelea, sino que sólo había intentado separar a los dos grupos. Añadió que le dolía el brazo desde aquel día y que no podía escribir. El presidente le preguntó si tenía un certificado médico. El alumno contestó que su madre lo llevaría al hospital al día siguiente. Después de dos horas de discusión, el consejo decidió que nadie sería excluido, pero que todos deberían pedir disculpas delante de la clase.
\end{textebox}

\textbf{Questions de compréhension} (réponds par des phrases complètes en espagnol) :
\begin{enumerate}
\item \esp{¿Por qué se reunió el consejo de disciplina?}
\item \esp{¿Qué anunció el presidente sobre los castigos?}
\item \esp{¿Qué explicó uno de los alumnos para defenderse?}
\end{enumerate}

\textbf{Questions de langue} :
\begin{enumerate}
\item Relève dans le texte deux verbes défectifs et explique leur construction.
\item Relève deux tournures restrictives et traduis-les en français.
\item Mets au discours direct la phrase : \esp{El alumno contestó que su madre lo llevaría al hospital al día siguiente.}
\end{enumerate}

\textbf{Sujet de rédaction} : \esp{Tu clase ha organizado una reunión sobre la disciplina en el instituto. Escribe un artículo para el periódico escolar en el que relates lo que dijeron el director y los delegados, utilizando el discurso indirecto.} (120 à 150 mots)
\end{exercice}

\newpage

\coursec{Corrigés des exercices}

\coursec{Corrigés des exercices — Leçon E14 : Traduction (verbes défectifs, restriction, discours indirect)}

\begin{corrige}[Exercice 1 — QCM : dépiste les pièges]
\begin{enumerate}
\item \textbf{c)} \esp{A mí me gustan los partidos de fútbol.} Le verbe \esp{gustar} s'accorde avec ce qui plaît (\esp{los partidos}, pluriel), pas avec la personne qui éprouve le goût. « J'aime » se dit \esp{me gusta(n)}, jamais \esp{me gusto} (qui signifierait « je plais »).
\item \textbf{a)} \esp{En la reunión no habló más que el director.} « À la réunion, seul le directeur a parlé. » \esp{No... más que} = « seulement ». \esp{Sino} exigerait une seconde partie contrastée (\esp{no habló el director sino el profesor}) ; \esp{apenas} ne se combine pas avec \esp{no} (valeur négative intrinsèque).
\item \textbf{b)} \esp{No compré pan sino arroz.} « Je n'ai pas acheté du pain, mais du riz. » Après une négation, la rectification « mais » se traduit par \esp{sino}. \esp{No... más que} signifierait « je n'ai acheté que » et exigerait \esp{más que arroz} seul, sans premier élément nié.
\item \textbf{b)} \esp{La directora dijo que recibiría a los delegados al día siguiente.} Futur \esp{recibiré} $\rightarrow$ conditionnel \esp{recibiría} (concordance des temps après \esp{dijo}) ; \esp{mañana} $\rightarrow$ \esp{al día siguiente}.
\item \textbf{b)} \esp{Ayer llovió mucho en Douala.} \esp{Llover} est un verbe impersonnel : il ne s'emploie qu'à la troisième personne du singulier (\esp{llueve}, \esp{llovió}, \esp{llovía}).
\end{enumerate}
\end{corrige}

\begin{corrige}[Exercice 2 — Vrai ou faux ? Justifie]
\begin{enumerate}
\item \textbf{Faux.} \esp{Gustar} fonctionne à l'envers du français « aimer » : la chose aimée est le sujet du verbe et la personne est un complément introduit par \esp{a}. On dit \esp{Me gusta el café} (« le café me plaît »), jamais \esp{yo gusto el café}.
\item \textbf{Vrai.} \esp{Doler} se construit comme \esp{gustar} : \esp{me duelen los pies} = « les pieds me font mal », c'est-à-dire « j'ai mal aux pieds ». Le verbe est au pluriel parce que \esp{los pies} est le sujet.
\item \textbf{Vrai.} \esp{Soler} + infinitif exprime une habitude : \esp{Suelo levantarme a las cinco} = « J'ai l'habitude de me lever à cinq heures ». On peut aussi traduire par « d'habitude » + présent.
\item \textbf{Vrai.} Au discours indirect, les déictiques de proximité reculent : \esp{aquí} $\rightarrow$ \esp{allí}, \esp{ahora} $\rightarrow$ \esp{entonces}, \esp{hoy} $\rightarrow$ \esp{aquel día}, \esp{mañana} $\rightarrow$ \esp{al día siguiente}.
\item \textbf{Faux.} \esp{No tengo más que dos hermanos} signifie « Je n'ai que deux frères » (restriction : deux, pas plus). La tournure \esp{no... más que} équivaut à \esp{solo} et ne nie pas le nombre indiqué.
\end{enumerate}
\end{corrige}

\begin{corrige}[Exercice 3 — Vocabulaire : les verbes introducteurs]
\textbf{Associations :} \esp{decir} = dire ; \esp{preguntar} = demander ; \esp{contestar} = répondre ; \esp{mandar} = ordonner ; \esp{exclamar} = s'exclamer.

\textbf{Phrases complétées (passé simple) :}
\begin{enumerate}
\item \esp{El profesor dijo que el examen sería al día siguiente.} « Le professeur a dit que l'examen aurait lieu le lendemain. »
\item \esp{Los alumnos preguntaron si podían salir antes.} « Les élèves ont demandé s'ils pouvaient sortir plus tôt. »
\item \esp{El jefe mandó que cerraran la puerta.} « Le chef a ordonné qu'on ferme la porte. » (\esp{mandar que} + subjonctif imparfait : \esp{cerraran}.)
\item \esp{Ella exclamó: «¡Qué alegría verte!»} « Elle s'est exclamée : "Quelle joie de te voir !" »
\end{enumerate}
\end{corrige}

\begin{attention}[Passé simple irrégulier]
\esp{decir} $\rightarrow$ \esp{dije, dijiste, dijo, dijimos, dijisteis, dijeron} ; \esp{preguntar}, \esp{contestar}, \esp{mandar}, \esp{exclamar} sont réguliers en \esp{-ar}. Ne confonds pas \esp{preguntar} (demander) et \esp{preguntar por} (demander des nouvelles de).
\end{attention}


\begin{corrige}[Exercice 4 — Conjugaison à compléter]
\begin{enumerate}
\item \esp{A nosotros nos gustan las clases de español.} « Nous aimons les cours d'espagnol. » Sujet pluriel : \esp{gustan} ; pronom pluriel : \esp{nos}.
\item \esp{A mí me duele la cabeza desde ayer.} « J'ai mal à la tête depuis hier. » Sujet singulier (\esp{la cabeza}) : \esp{duele} (verbe à diphtongaison \esp{o $\rightarrow$ ue}).
\item \esp{¿A ti te parece bien la idea?} « L'idée te semble-t-elle bonne ? » Sujet : \esp{la idea} $\rightarrow$ \esp{parece}.
\item \esp{Mi abuela suele ir al mercado de Mvog-Mbi los sábados.} « Ma grand-mère a l'habitude d'aller au marché de Mvog-Mbi le samedi. » \esp{Soler} : diphtongaison \esp{o $\rightarrow$ ue} (\esp{suelo, sueles, suele, solemos, soléis, suelen}).
\item \esp{Este asunto no concierne a los alumnos de primero.} « Cette affaire ne concerne pas les élèves de première. » \esp{Concernir} est défectif : il s'emploie surtout à la 3\textsuperscript{e} personne ; sujet singulier : \esp{concierne} (\esp{e $\rightarrow$ ie}).
\item \esp{En diciembre llueve poco en Yaoundé.} « En décembre, il pleut peu à Yaoundé. » Verbe impersonnel, 3\textsuperscript{e} personne du singulier : \esp{llueve}.
\end{enumerate}
\end{corrige}

\begin{corrige}[Exercice 5 — Traduction : phrases isolées]
\begin{enumerate}
\item \esp{Me gusta mucho la música camerunesa.} $\rightarrow$ \esp{Me gusta mucho la música camerunesa.} Attention : \esp{gustar} au singulier (\esp{la música} est le sujet) ; adjectif \esp{camerunés, camerunesa}.
\item \esp{A mis hermanos les duelen las piernas después del partido.} Pluriel \esp{les duelen} car le sujet est \esp{las piernas} ; ne pas oublier la reprise \esp{a mis hermanos... les}.
\item \esp{Solo tenemos diez minutos para terminar el ejercicio.} ou \esp{No tenemos más que diez minutos para terminar el ejercicio.} Les deux tournures restrictives sont correctes.
\item \esp{No es un bolígrafo, sino un lápiz.} Rectification après négation : \esp{sino}, jamais \esp{pero} ici.
\item \esp{Suele levantarse a las cinco de la mañana.} ou \esp{Tiene la costumbre de levantarse a las cinco de la mañana.} \esp{Soler} + infinitif, diphtongaison \esp{suele}.
\end{enumerate}
\end{corrige}

\begin{corrige}[Exercice 6 — Du discours direct au discours indirect]
\begin{enumerate}
\item \esp{Juan dijo que vendría a mi casa al día siguiente.} Futur \esp{vendré} $\rightarrow$ conditionnel \esp{vendría} ; \esp{tu} $\rightarrow$ \esp{mi} ; \esp{mañana} $\rightarrow$ \esp{al día siguiente}.
\item \esp{María preguntó dónde estaba su libro.} Question partielle : pas de \esp{si}, on garde le mot interrogatif avec accent (\esp{dónde}) ; présent \esp{está} $\rightarrow$ imparfait \esp{estaba} ; \esp{mi} $\rightarrow$ \esp{su}.
\item \esp{El policía ordenó que no tocaran nada allí.} Impératif $\rightarrow$ \esp{ordenó que} + subjonctif imparfait (\esp{tocaran}) ; \esp{aquí} $\rightarrow$ \esp{allí}.
\item \esp{El director contestó que entonces estaba muy ocupado.} Présent $\rightarrow$ imparfait ; \esp{ahora} $\rightarrow$ \esp{entonces}.
\item \esp{Los niños dijeron que el día anterior habían ido al mercado.} Passé simple \esp{fuimos} $\rightarrow$ plus-que-parfait \esp{habían ido} ; \esp{ayer} $\rightarrow$ \esp{el día anterior} ; 1\up{re} personne $\rightarrow$ 3\up{e}.
\item \esp{La profesora preguntó a Pedro si había terminado el trabajo aquel día.} Question totale introduite par \esp{si} ; passé composé \esp{has terminado} $\rightarrow$ plus-que-parfait \esp{había terminado} ; \esp{hoy} $\rightarrow$ \esp{aquel día}.
\end{enumerate}
\end{corrige}

\begin{attention}[Points de vigilance]
Ne jamais garder les guillemets ni les signes \esp{¿ ?} au discours indirect. Vérifier systématiquement les trois transformations : personnes (pronoms et possessifs), temps du verbe, indicateurs spatio-temporels.
\end{attention}


\begin{corrige}[Exercice 7 — Texte à trous]
\begin{enumerate}
\item \esp{ese} : \esp{ese día} (« ce jour-là ») — au discours indirect, \esp{este} $\rightarrow$ \esp{ese}.
\item \esp{gustó} : \esp{a nosotros nos gustó mucho la idea} — passé simple de \esp{gustar}, singulier car le sujet est \esp{la idea}.
\item \esp{parecía} : \esp{nos parecía interesante} — imparfait (appréciation dans le passé), sujet : l'infinitif \esp{hablar}.
\item \esp{solía} : \esp{solía entrevistar a jóvenes} — habitude passée : imparfait de \esp{soler} + infinitif.
\item \esp{dolía} : \esp{le dolía mucho la cabeza} — imparfait de \esp{doler}, singulier (\esp{la cabeza}).
\item \esp{la} : \esp{volvería la semana siguiente} — \esp{la próxima semana} devient \esp{la semana siguiente} au discours indirect.
\item \esp{entonces} : \esp{entonces podríamos hacerle todas las preguntas} — \esp{ahora} $\rightarrow$ \esp{entonces}.
\end{enumerate}
\textbf{Texte complet :} \esp{Ayer, una periodista visitó nuestro instituto de Yaoundé. Nos dijo que ese día quería hablar con los alumnos de Terminale. A nosotros nos gustó mucho la idea, porque nos parecía interesante hablar con una profesional. Ella explicó que solía entrevistar a jóvenes en las escuelas de la capital. Un alumno contestó que no podía quedarse porque le dolía mucho la cabeza. La periodista dijo que volvería la semana siguiente y que entonces podríamos hacerle todas las preguntas.}

\emph{Traduction : « Hier, une journaliste a visité notre lycée de Yaoundé. Elle nous a dit que ce jour-là elle voulait parler avec les élèves de Terminale. L'idée nous a beaucoup plu, car il nous semblait intéressant de parler avec une professionnelle. Elle a expliqué qu'elle avait l'habitude d'interviewer des jeunes dans les écoles de la capitale. Un élève a répondu qu'il ne pouvait pas rester parce qu'il avait très mal à la tête. La journaliste a dit qu'elle reviendrait la semaine suivante et qu'alors nous pourrions lui poser toutes nos questions. »}
\end{corrige}

\begin{corrige}[Exercice 8 — Reformulation : la restriction]
\begin{enumerate}
\item \esp{No tengo más que cinco mil francos.} « Je n'ai que cinq mille francs. »
\item \esp{Solo vinieron tres delegados.} ou \esp{Únicamente vinieron tres delegados.} « Seuls trois délégués sont venus. »
\item \esp{No hablamos más que español en esta clase.} « Nous ne parlons qu'espagnol dans cette classe. »
\item \esp{Solo queda un poco de arroz.} « Il ne reste qu'un peu de riz. »
\item \esp{No pueden entrar más que los delegados.} « Seuls les délégués peuvent entrer. »
\item \esp{Solo leí el primer capítulo.} « Je n'ai lu que le premier chapitre. »
\end{enumerate}
\end{corrige}

\begin{attention}[Attention]
La négation \esp{no} doit précéder le verbe dans la tournure \esp{no... más que} : \esp{No tengo más que...}. Ne dis jamais \esp{tengo no más que} à l'écrit scolaire. Rappel : depuis 2010, l'adverbe \esp{solo} s'écrit sans accent, mais l'accent traditionnel (\esp{sólo}) reste toléré quand il y a ambiguïté.
\end{attention}


\begin{corrige}[Exercice 9 — Version : une réunion de parents d'élèves]
\textbf{Traduction proposée :}

« Samedi dernier, une réunion de parents d'élèves a eu lieu au lycée de Biyem-Assi. Le directeur a ouvert la séance et a dit que cette année-là les résultats étaient meilleurs, mais qu'il restait encore beaucoup de travail. Il a ajouté que seuls les parents inscrits pouvaient voter à la fin de la réunion. Une mère a demandé si la cantine fonctionnerait le lendemain, parce que son fils n'aimait pas la nourriture de la rue et qu'il avait souvent mal à l'estomac. Le directeur a répondu qu'il ne s'agissait pas d'un problème de cuisine, mais d'organisation, et qu'ils n'avaient guère d'argent pour embaucher du personnel supplémentaire. Un autre père a déclaré qu'il lui semblait urgent de réparer les salles de classe avant la saison des pluies. À la fin, le directeur a promis qu'il étudierait toutes les propositions et qu'il donnerait une réponse la semaine suivante. »

\textbf{Justification de trois choix :}
\begin{enumerate}
\item \emph{Verbe défectif} : \esp{a su hijo no le gustaba la comida de la calle} se traduit « son fils n'aimait pas la nourriture de la rue ». \esp{Gustar} a pour sujet la chose (\esp{la comida}) et pour complément la personne (\esp{le}) ; le français inverse la construction avec « aimer ».
\item \emph{Tournure restrictive} : \esp{solo los padres inscritos podían votar} = « seuls les parents inscrits pouvaient voter » ; on aurait pu dire \esp{no podían votar más que los padres inscritos}. De même, \esp{apenas tenían dinero} = « ils n'avaient guère / presque pas d'argent ».
\item \emph{Discours indirect} : \esp{el director prometió que estudiaría todas las propuestas y que daría una respuesta la semana siguiente} : le futur du discours direct (\esp{estudiaré, daré}) passe au conditionnel (\esp{estudiaría, daría}) après un verbe introducteur au passé, et \esp{la próxima semana} devient \esp{la semana siguiente} (« la semaine suivante »).
\end{enumerate}

\textbf{Barème indicatif :} traduction 14 points (exactitude du sens 8 pts, qualité du français 3 pts, rendu des trois phénomènes grammaticaux 3 pts) ; justifications 6 points (2 pts par choix justifié).
\end{corrige}

\begin{corrige}[Exercice 10 — Thème type Bac]
\begin{enumerate}
\item \esp{La directora dijo que solo recibiría a los delegados al día siguiente.} ou \esp{...que no recibiría más que a los delegados al día siguiente.} Futur dans le passé : conditionnel \esp{recibiría} ; « le lendemain » : \esp{al día siguiente}.
\item \esp{A mis compañeros les gusta mucho este profesor porque sus clases son interesantes.} \esp{Gustar} au singulier (sujet : \esp{este profesor}) ; pronom pluriel \esp{les} avec reprise \esp{a mis compañeros}.
\item \esp{Aquel día por la mañana llovía mucho y los alumnos llegaron tarde.} « Ce matin-là » : \esp{aquel día por la mañana} ; pluie en arrière-plan : imparfait \esp{llovía} ; arrivée ponctuelle : passé simple \esp{llegaron}.
\item \esp{El profesor preguntó a los alumnos si habían terminado sus deberes aquel día.} Question totale : \esp{si} ; concordance : plus-que-parfait \esp{habían terminado} ; « ce jour-là » : \esp{aquel día}.
\item \esp{No es el cansancio lo que me preocupa, sino tu salud: sueles dormir muy poco.} Rectification : \esp{no... sino} ; « tu as l'habitude de » : \esp{sueles} + infinitif.
\end{enumerate}
\textbf{Barème indicatif :} 4 points par phrase (construction correcte du phénomène visé 2 pts, morphologie et orthographe 1 pt, accents et ponctuation 1 pt).

\textbf{Points de vigilance :} ne jamais traduire « aimer » par \esp{amar} pour des choses ou des personnes sans nuance affective forte — \esp{gustar} est la tournure naturelle ; ne pas oublier le \esp{a} personnel devant \esp{los delegados} et \esp{los alumnos} ; vérifier l'absence d'accent sur l'adjectif \esp{aquel} (accent possible seulement sur le pronom \esp{aquél} en cas d'ambiguïté) et l'absence de \esp{¿} au discours indirect.
\end{corrige}

\begin{corrige}[Exercice 11 — Compréhension et rédaction]
\textbf{Questions de compréhension (réponses modèles) :}
\begin{enumerate}
\item \esp{El consejo de disciplina se reunió para estudiar el caso de dos alumnos.} « Le conseil de discipline s'est réuni pour étudier le cas de deux élèves. »
\item \esp{El presidente anunció que no castigaría más que a los responsables de la pelea, y no a los simples testigos.} « Le président a annoncé qu'il ne punirait que les responsables de la bagarre, et non les simples témoins. »
\item \esp{Uno de los alumnos explicó que no había empezado la pelea, sino que solo había intentado separar a los dos grupos.} « L'un des élèves a expliqué qu'il n'avait pas commencé la bagarre, mais qu'il avait seulement essayé de séparer les deux groupes. »
\end{enumerate}

\textbf{Questions de langue :}
\begin{enumerate}
\item Deux verbes défectifs : \esp{dolía} dans \esp{le dolía el brazo} (« il avait mal au bras ») — le sujet est \esp{el brazo} et la personne est exprimée par le pronom \esp{le} ; et \esp{parecía} dans \esp{les parecía una buena idea} (« cela leur semblait une bonne idée ») — sujet : l'infinitif / la proposition. \esp{Se reunió} n'est pas défectif (verbe pronominal voix passive/réfléchie).
\item Deux tournures restrictives : \esp{no castigaría más que a los responsables} = « il ne punirait que les responsables » ; \esp{solo había intentado separar a los dos grupos} = « il avait seulement essayé de séparer les deux groupes ». On peut ajouter la rectification \esp{no... sino que} = « non pas... mais ».
\item Discours direct : \esp{El alumno contestó: «Mi madre me llevará al hospital mañana».} Conditionnel \esp{llevaría} $\rightarrow$ futur \esp{llevará} ; \esp{su/lo} $\rightarrow$ \esp{mi/me} ; \esp{al día siguiente} $\rightarrow$ \esp{mañana}.
\end{enumerate}

\textbf{Proposition de rédaction modèle :}

\begin{textebox}[Una reunión sobre la disciplina]
El lunes pasado, nuestra clase organizó una reunión sobre la disciplina en el instituto. El director abrió la sesión y dijo que aquel año había menos problemas que antes, pero que todavía quedaba mucho por hacer. Añadió que no castigaría más que a los alumnos que no respetaran el reglamento. Un delegado preguntó si podían organizarse actividades deportivas los sábados, porque a muchos alumnos les gustaba el fútbol y les parecía una buena manera de canalizar su energía. El director contestó que le parecía una idea excelente, pero que apenas tenían dinero para comprar material. Otra delegada explicó que a ella le dolía ver las aulas sucias y propuso que cada clase limpiara su sala al final de la semana. Al final, el director prometió que estudiaría todas las propuestas y que daría una respuesta la semana siguiente. Todos los alumnos salieron contentos de aquella reunión.
\end{textebox}

\emph{(Environ 150 mots.) Traduction : « Lundi dernier, notre classe a organisé une réunion sur la discipline au lycée. Le directeur a ouvert la séance et a dit que cette année-là il y avait moins de problèmes qu'avant, mais qu'il restait encore beaucoup à faire. Il a ajouté qu'il ne punirait que les élèves qui ne respecteraient pas le règlement. Un délégué a demandé si des activités sportives pouvaient être organisées le samedi, parce que beaucoup d'élèves aimaient le football et que cela leur semblait une bonne manière de canaliser leur énergie. Le directeur a répondu que l'idée lui semblait excellente, mais qu'ils n'avaient guère d'argent pour acheter du matériel. Une autre déléguée a expliqué que cela lui faisait mal de voir les salles sales et a proposé que chaque classe nettoie sa salle en fin de semaine. À la fin, le directeur a promis qu'il étudierait toutes les propositions et qu'il donnerait une réponse la semaine suivante. Tous les élèves sont sortis contents de cette réunion. »}

\textbf{Points de vigilance pour la rédaction :} employer au moins trois verbes introducteurs variés (\esp{decir, añadir, preguntar, contestar, explicar, prometer}) ; vérifier la concordance des temps (présent $\rightarrow$ imparfait, futur $\rightarrow$ conditionnel, passé $\rightarrow$ plus-que-parfait) ; transposer les déictiques (\esp{este $\rightarrow$ aquel}, \esp{mañana $\rightarrow$ al día siguiente}) ; soigner les accents et la ponctuation \esp{¿ ¡}.
\end{corrige}

\newpage

\coursec{Fiche bilan}

\begin{popfigure}
\begin{tikzpicture}[mindmap, grow cyclic, every node/.style={concept, text=white, font=\small\bfseries, align=center}, root concept/.append style={concept color=popblue, font=\large\bfseries}, level 1 concept/.append style={sibling angle=51, level distance=3.6cm, minimum size=2.1cm}, level 2 concept/.append style={level distance=2.2cm, minimum size=1.5cm, font=\scriptsize\bfseries}]
\node[root concept, align=center]{Traduction\\ E14}
  child[concept color=poporange]{ node[align=center]{Verbes\\ défectifs}
    child[concept color=poporange!70]{ node[align=center]{gustar\\ doler\\ parecer} }
    child[concept color=poporange!70]{ node[align=center]{soler\\ llover\\ atañer} }
  }
  child[concept color=popgreen]{ node[align=center]{Structure\\ gustar}
    child[concept color=popgreen!70]{ node[align=center]{A mí\\ me gusta} }
  }
  child[concept color=popred]{ node{Restriction}
    child[concept color=popred!70]{ node[align=center]{sólo\\ únicamente} }
    child[concept color=popred!70]{ node[align=center]{no... más que\\ no... sino} }
  }
  child[concept color=poppurple]{ node[align=center]{Discours\\ indirect}
    child[concept color=poppurple!70]{ node[align=center]{decir que\\ preguntar si} }
    child[concept color=poppurple!70]{ node[align=center]{mandar que\\ + subj.} }
  }
  child[concept color=popgold]{ node[align=center]{Concordance\\ des temps}
    child[concept color=popgold!70]{ node{prés. $\to$ imparf.} }
    child[concept color=popgold!70]{ node{pret. perf. $\to$ plusc.} }
  }
  child[concept color=poppink]{ node[align=center]{Repères\\ spatio-temp.}
    child[concept color=poppink!70]{ node[align=center]{hoy $\to$ ese día\\ aquí $\to$ allí} }
  }
  child[concept color=popblue!60]{ node[align=center]{Méthode\\ thème / version}
    child[concept color=popblue!40]{ node[align=center]{repérer,\\ transposer,\\ vérifier} }
  };
\end{tikzpicture}
\legende{Carte mentale de la leçon : verbes défectifs, restriction et discours indirect}
\end{popfigure}

\begin{bilan}
\textbf{1. Lexique clé (verbes défectifs et restriction)}

\begin{center}
\begin{tabularx}{\linewidth}{|X|X|}\hline
\rowcolor{popblueL} \textbf{Español} & \textbf{Français} \\ \hline
\esp{gustar} & plaire (aimer) \\ \hline
\esp{doler (ue)} & faire mal, souffrir \\ \hline
\esp{parecer} & sembler, paraître \\ \hline
\esp{interesar / importar} & intéresser / importer \\ \hline
\esp{encantar / fascinar} & adorer / passionner \\ \hline
\esp{soler (ue) + inf.} & avoir l'habitude de, d'habitude \\ \hline
\esp{llover (ue), nevar, granizar} & pleuvoir, neiger, grêler (3\textsuperscript{e} pers. sing.) \\ \hline
\esp{atañer / concernir} & concerner (3\textsuperscript{e} pers. seulement) \\ \hline
\esp{sólo / únicamente / solamente} & seulement \\ \hline
\esp{no... más que / no... sino} & ne... que / ne... mais \\ \hline
\esp{apenas} & à peine \\ \hline
\end{tabularx}
\end{center}

\textbf{2. Grammaire à connaître par cœur}

\begin{center}
\begin{tabularx}{\linewidth}{|X|X|}\hline
\rowcolor{popgreenL} \textbf{Règle} & \textbf{Exemple} \\ \hline
\esp{gustar} se conjugue selon la chose aimée : \esp{me gusta} + sing./inf., \esp{me gustan} + plur. & \esp{Me gusta el cacao. Me gustan las novelas.} \\ \hline
La personne est au COI : \esp{a mí me, a ti te, a él le, a nosotros nos, a vosotros os, a ellos les} & \esp{A Juan le duelen los pies.} \\ \hline
\esp{soler} + infinitif = habitude présente ; imparfait = habitude passée & \esp{Suelo leer por la noche. Solía leer...} \\ \hline
Verbes météo : 3\textsuperscript{e} pers. sing. uniquement & \esp{Llueve mucho en Douala. Ayer llovió.} \\ \hline
\esp{no... más que} = ne... que ; \esp{no... sino} = non pas... mais (correction) & \esp{No tengo más que cinco mil francos. No es médico, sino enfermero.} \\ \hline
Style indirect déclaratif : \esp{decir que} + indicatif, sans guillemets & \esp{Dice que está cansado.} \\ \hline
Interrogation totale : \esp{preguntar si} ; partielle : \esp{preguntar qué/dónde...} & \esp{Preguntó si venía. Me preguntó dónde vivía.} \\ \hline
Ordre/conseil : \esp{mandar/decir/pedir que} + subjonctif & \esp{Le dijo que viniera.} \\ \hline
Introducteur au passé $\Rightarrow$ recul des temps : prés. $\to$ imparf., pret. perf. $\to$ plusc., fut. $\to$ conditionnel & \esp{Dijo: «voy» $\to$ Dijo que iba.} \\ \hline
Repères : \esp{hoy $\to$ ese día, ayer $\to$ el día anterior, mañana $\to$ el día siguiente, aquí $\to$ allí, este $\to$ ese} & \esp{«Vendré mañana aquí» $\to$ que vendría al día siguiente allí.} \\ \hline
\end{tabularx}
\end{center}

\textbf{3. Méthodes express}
\begin{itemize}
  \item \textbf{Version :} repérer le verbe défectif $\to$ identifier le vrai sujet grammatical $\to$ accorder le verbe avec la chose, pas avec la personne.
  \item \textbf{Thème :} « aimer » suivi d'une chose $\Rightarrow$ penser \esp{gustar} (jamais \esp{amar} pour les choses) ; « ne... que » $\Rightarrow$ \esp{no... más que} ou \esp{sólo}.
  \item \textbf{Discours indirect :} 1) choisir le verbe introducteur ; 2) supprimer guillemets et deux-points ; 3) transposer les personnes ; 4) reculer les temps si l'introducteur est au passé ; 5) adapter les repères spatio-temporels.
  \item \textbf{Vérification finale :} relire en contrôlant accord du verbe, subjonctif après \esp{mandar que}, et ponctuation espagnole (¿ ¡).
\end{itemize}
\end{bilan}

\begin{aretenir}[Checklist avant le Bac]
\begin{itemize}
  \item $\square$ Je sais conjuguer \esp{gustar, doler, parecer} avec les pronoms COI (\esp{me, te, le, nos, os, les}).
  \item $\square$ Je sais que le verbe s'accorde avec la chose aimée : \esp{me gustan los libros}.
  \item $\square$ Je n'écris jamais \esp{yo gusto} : je dis \esp{a mí me gusta}.
  \item $\square$ Je sais employer \esp{soler} + infinitif pour traduire « d'habitude ».
  \item $\square$ Je sais que les verbes météo n'existent qu'à la 3\textsuperscript{e} personne du singulier : \esp{llueve, llovía, llovió}.
  \item $\square$ Je traduis « ne... que » par \esp{no... más que} ou \esp{sólo}, et « non pas... mais » par \esp{no... sino}.
  \item $\square$ Je sais choisir le verbe introducteur : \esp{decir que, preguntar si, mandar que, pedir que}.
  \item $\square$ J'applique la concordance des temps quand l'introducteur est au passé : présent $\to$ imparfait, futur $\to$ conditionnel.
  \item $\square$ Je mets le subjonctif après \esp{mandar que / decir que / pedir que} quand il y a ordre ou demande.
  \item $\square$ Je transpose les repères : \esp{hoy $\to$ ese día}, \esp{mañana $\to$ el día siguiente}, \esp{aquí $\to$ allí}.
  \item $\square$ Je change les personnes et les possessifs selon le nouveau locuteur (\esp{mi $\to$ su}, \esp{yo $\to$ él}).
  \item $\square$ Je relis ma traduction en vérifiant les accents, le \esp{ñ} et la ponctuation espagnole ¿ ¡.
\end{itemize}
\end{aretenir}

\findecours

\end{document}
